% Generated from the public Markdown study guide. Edit this file directly for print-specific changes.
% Recommended build command: tectonic -X compile ecd-exam-study-guide.tex
% Keep the local PNG icon assets in the adjacent icons/ directory when compiling.
\documentclass[11pt,letterpaper,oneside,openany]{book}
\usepackage[margin=0.82in,top=0.88in,bottom=0.84in,headheight=22pt]{geometry}
\usepackage{libertinus}
\usepackage{microtype}
\usepackage{setspace}
\usepackage{xcolor}
\usepackage{graphicx}
\usepackage{tikz}
\usetikzlibrary{calc,positioning}
\usepackage{titlesec}
\usepackage{fancyhdr}
\usepackage{enumitem}
\usepackage{booktabs}
\usepackage{array}
\usepackage{longtable}
\usepackage{ragged2e}
\usepackage{colortbl}
\usepackage[skins,breakable]{tcolorbox}
\usepackage{fvextra}
\usepackage{needspace}
\usepackage{ulem}
\usepackage{hyperref}
\usepackage{bookmark}
\usepackage{tocloft}

\definecolor{DeepNavy}{HTML}{203444}
\definecolor{Moss}{HTML}{667D6D}
\definecolor{MutedGold}{HTML}{C58A3A}
\definecolor{WarmIvory}{HTML}{F7F4ED}
\definecolor{SoftGray}{HTML}{E9ECEA}
\definecolor{Ink}{HTML}{26323A}

\hypersetup{
  colorlinks=true,
  linkcolor=DeepNavy,
  urlcolor=Moss,
  citecolor=MutedGold,
  pdftitle={EDAC Ec.D. Exam Study Guide},
  pdfauthor={Xiao},
  pdfsubject={Economic development certification exam preparation},
  pdfkeywords={EDAC, Ec.D., economic development, Canada, exam study guide}
}

\setstretch{1.07}
\setlength{\parindent}{0pt}
\setlength{\parskip}{0.58em}
\setcounter{secnumdepth}{2}
\setcounter{tocdepth}{2}
\setlist{itemsep=0.22em,topsep=0.38em,leftmargin=1.6em}
\renewcommand{\arraystretch}{1.25}
\setlength{\tabcolsep}{4pt}
\setlength{\emergencystretch}{3em}
\hbadness=2000
\renewcommand{\chaptername}{Chapter}
\newcommand{\chaptericon}[1]{\raisebox{-0.14em}{\includegraphics[height=1.12em]{icons/#1.png}}\hspace{0.42em}}
\newif\ifsectionhasicon
\newcommand{\currentsectionicon}{}
\newcommand{\marginsectionicon}{%
  \ifsectionhasicon
    \llap{\raisebox{-0.12em}{\includegraphics[height=0.92em]{icons/\currentsectionicon.png}}\hspace{0.14in}}%
  \fi
}
\newcommand{\marginsection}[2]{%
  \begingroup
  \sectionhasicontrue
  \def\currentsectionicon{#1}%
  \section[#2]{#2}%
  \endgroup
}

\titleformat{\chapter}[display]
  {\normalfont\sffamily\bfseries\color{DeepNavy}\raggedright}
  {\filleft\fontsize{13}{15}\selectfont\color{MutedGold}\MakeUppercase{\chaptertitlename}\ \thechapter}
  {0.5ex}
  {\titlerule\vspace{1.1ex}\Huge}
  [\vspace{0.7ex}\color{Moss}\titlerule]
\titlespacing*{\chapter}{0pt}{-18pt}{22pt}
\titleformat{\section}
  {\Large\sffamily\bfseries\color{DeepNavy}}
  {\marginsectionicon\thesection}{0.7em}{}
\titleformat{\subsection}
  {\large\sffamily\bfseries\color{Moss}}
  {\thesubsection}{0.7em}{}
\titleformat{\subsubsection}
  {\normalsize\sffamily\bfseries\color{Ink}}
  {\thesubsubsection}{0.7em}{}

\pagestyle{fancy}
\fancyhf{}
\fancyhead[L]{\small\sffamily\color{Moss}\nouppercase{\leftmark}}
\fancyhead[R]{\small\sffamily\color{Moss}EDAC Ec.D. Exam Study Guide}
\fancyfoot[C]{\small\sffamily\color{Ink}\thepage}
\renewcommand{\headrulewidth}{0.4pt}
\renewcommand{\headrule}{\hbox to\headwidth{\color{SoftGray}\leaders\hrule height \headrulewidth\hfill}}
\fancypagestyle{plain}{
  \fancyhf{}
  \fancyfoot[C]{\small\sffamily\color{Ink}\thepage}
  \renewcommand{\headrulewidth}{0pt}
}

\newtcolorbox{guidequote}{
  enhanced,
  breakable,
  colback=WarmIvory,
  colframe=Moss,
  boxrule=0pt,
  borderline west={2pt}{0pt}{Moss},
  left=10pt,right=8pt,top=6pt,bottom=6pt,
  before skip=8pt,after skip=8pt
}

\renewcommand{\cfttoctitlefont}{\Huge\sffamily\bfseries\color{DeepNavy}}
\renewcommand{\cftchapfont}{\sffamily\bfseries\color{DeepNavy}}
\renewcommand{\cftchappagefont}{\sffamily\bfseries\color{DeepNavy}}
\renewcommand{\cftsecfont}{\small\sffamily}
\renewcommand{\cftsecpagefont}{\small\sffamily}
\setlength{\cftbeforechapskip}{0.42em}
\setlength{\cftbeforesecskip}{0.08em}

\begin{document}
\frontmatter
\hypersetup{pageanchor=false}

\begin{titlepage}
\begin{tikzpicture}[remember picture,overlay]
  \fill[DeepNavy] (current page.south west) rectangle (current page.north east);
  \fill[Moss] ($(current page.south west)+(0,0)$) rectangle ($(current page.south east)+(0,0.34in)$);
  \fill[MutedGold] ($(current page.north east)+(-1.55in,0)$) rectangle ($(current page.south east)+(-1.23in,0)$);
  \draw[WarmIvory,opacity=0.18,line width=0.55pt]
    ($(current page.north west)+(0.75in,-1.0in)$) -- ($(current page.north east)+(-1.37in,-2.35in)$)
    -- ($(current page.south east)+(-1.37in,2.05in)$)
    -- ($(current page.south west)+(0.72in,1.22in)$);
  \foreach \x/\y in {1.1/7.9,2.05/8.55,3.05/7.35,4.2/8.15,5.35/7.1,6.55/7.85} {
    \fill[WarmIvory,opacity=0.7] ($(current page.south west)+(\x in,\y in)$) circle (2.2pt);
  }
  \draw[WarmIvory,opacity=0.38,line width=0.65pt]
    ($(current page.south west)+(1.1in,7.9in)$) --
    ($(current page.south west)+(2.05in,8.55in)$) --
    ($(current page.south west)+(3.05in,7.35in)$) --
    ($(current page.south west)+(4.2in,8.15in)$) --
    ($(current page.south west)+(5.35in,7.1in)$) --
    ($(current page.south west)+(6.55in,7.85in)$);
\end{tikzpicture}

\color{WarmIvory}
\vspace*{0.63in}
{\sffamily\small\bfseries\color{MutedGold}\MakeUppercase{Public study edition}\par}
\vspace{1.05in}
{\sffamily\fontsize{38}{43}\selectfont\bfseries EDAC Ec.D.\par}
\vspace{0.08in}
{\sffamily\fontsize{34}{39}\selectfont\bfseries Exam Study Guide\par}
\vspace{0.38in}
{\color{WarmIvory}\rule{0.78\linewidth}{1.2pt}\par}
\vspace{0.32in}
{\sffamily\Large A practical, independently created guide for\par}
{\sffamily\Large Canadian economic development professionals\par}
\vfill
{\sffamily\large\bfseries Xiao\par}
\vspace{0.10in}
{\sffamily\normalsize July 2026\par}
\vspace{0.34in}
{\sffamily\small This guide is independently created and is not endorsed by EDAC.\par}
\vspace{0.20in}
{\hypersetup{urlcolor=WarmIvory}\sffamily\small\href{https://www.linkedin.com/in/xxiao27/}{linkedin.com/in/xxiao27}\par}
\vspace{0.08in}
{\hypersetup{urlcolor=WarmIvory}\sffamily\small\href{https://github.com/B1ackCat7/ecd-exam-study-guide}{github.com/B1ackCat7/ecd-exam-study-guide}\par}
\vspace{0.18in}
\end{titlepage}

\thispagestyle{empty}
\vspace*{\fill}
{\sffamily\bfseries\color{DeepNavy}EDAC Ec.D. Exam Study Guide\par}
\vspace{0.7em}
Copyright \textcopyright\ 2026 \href{https://www.linkedin.com/in/xxiao27/}{Xiao}.

\vspace{1em}
\begin{tcolorbox}[
  enhanced,
  colback=WarmIvory,
  colframe=MutedGold,
  boxrule=1pt,
  arc=1mm,
  left=10pt,
  right=10pt,
  top=9pt,
  bottom=9pt
]
{\sffamily\bfseries\large\color{DeepNavy}Open licence: CC BY 4.0\par}
\vspace{0.45em}
This guide is licensed under the Creative Commons Attribution 4.0 International License. You are welcome to copy, share, and adapt it for any purpose, including commercially, provided you give appropriate credit to Xiao, link to the licence, and indicate whether changes were made.

\href{https://creativecommons.org/licenses/by/4.0/}{creativecommons.org/licenses/by/4.0}
\end{tcolorbox}

\vspace{1em}
This is an independently created study resource. It is not an official publication of, sponsored by, or endorsed by the Economic Developers Association of Canada (EDAC). Exam requirements may change; candidates should confirm current details with EDAC.

\vspace{1em}
Public source links are included throughout the guide for verification and further learning.

\vspace{1em}
Version: July 2026
\vspace*{\fill}
\clearpage

\hypersetup{pageanchor=true}
\tableofcontents
\clearpage
\chapter*{\chaptericon{book-open}About This Guide}
\addcontentsline{toc}{chapter}{About This Guide}
\label{chap:about}
I'm Xiao. After several people asked how I prepared for the Economic Developers Association of Canada (EDAC) Certified Economic Developer (Ec.D.) exam, I put together this open, independently created study guide. I hope it helps fellow candidates approach the designation with more confidence.


The guide is written for Canadian economic development professionals. Its short Markdown chapters focus on understanding, judgment, and application to Canadian communities rather than memorized model answers.


\begin{guidequote}
This resource is independently created. It is not an official EDAC publication and is not endorsed by EDAC.

\end{guidequote}



\needspace{5\baselineskip}
\section[What this Ec.D. exam guide covers]{What this Ec.D. exam guide covers}
The study guide includes:


\begin{itemize}
\item current public EDAC Ec.D. exam requirements and preparation guidance;
\item economic development foundations, professional ethics, governance, and performance measurement;
\item strategic planning, business retention and expansion (BR\&E), investment attraction, partnerships, and networking;
\item communications, public engagement, land-use planning, marketing, sales, and site selection;
\item business finance, economic impact analysis, calculations, and formulas;
\item 24 original practice multiple-choice and true/false questions, plus a purpose-built written mock exam based on the study-guide topics;
\item Canadian economic development issues for contemporary-issue preparation; and
\item an 84-term glossary with public references for every definition.
\end{itemize}


\needspace{5\baselineskip}
\marginsection{route}{How to use this study guide}
\begin{enumerate}
\item Begin with the \hyperref[chap:exam-overview]{EDAC Ec.D. exam overview}.
\item Choose a schedule in the \hyperref[chap:study-plan]{study plan and exam strategy}.
\item Work through the foundations chapter and seven competency chapters:
\begin{itemize}
\item \hyperref[chap:foundations]{Economic development foundations and ethics}
\item \hyperref[chap:processes]{Economic development processes}
\item \hyperref[chap:governance]{Governance and performance measurement}
\item \hyperref[chap:operations]{Operations, partnerships, and networking}
\item \hyperref[chap:communications]{Communications and public engagement}
\item \hyperref[chap:planning]{Planning and land use}
\item \hyperref[chap:marketing]{Marketing, sales, and site selection}
\item \hyperref[chap:finance]{Finance, business planning, and economic impact analysis}
\end{itemize}

\item Practise the quantitative material in \hyperref[chap:calculations]{calculations and formulas}.
\item Use the \hyperref[chap:practice]{mock exam for practice} and \hyperref[chap:contemporary]{contemporary economic development issues} for timed practice.
\item Use the \hyperref[chap:glossary]{economic development glossary} for quick review and term-by-term authoritative references.
\end{enumerate}


\needspace{5\baselineskip}
\marginsection{file-text}{Print edition}
The \href{https://github.com/B1ackCat7/ecd-exam-study-guide/blob/main/output/pdf/ecd-exam-study-guide.pdf}{print-ready PDF} collects the full guide in one document. The \href{https://github.com/B1ackCat7/ecd-exam-study-guide/blob/main/print-edition/ecd-exam-study-guide.tex}{single-file LaTeX source} is also available if you want to adapt the typography, cover, or page layout.


\needspace{5\baselineskip}
\marginsection{lightbulb}{What a strong exam answer usually does}
A well-developed exam answer:


\begin{itemize}
\item identifies the objective and the local context;
\item uses a logical process rather than an isolated tactic;
\item names the stakeholders and explains their roles;
\item distinguishes facts, assumptions, risks, and trade-offs;
\item addresses confidentiality, public interest, inclusion, and ethics;
\item provides an implementation sequence, responsibilities, and resources;
\item defines outputs, outcomes, and success measures;
\item includes communication, monitoring, and a contingency path; and
\item connects professional concepts to the candidate's own work experience.
\end{itemize}


\needspace{5\baselineskip}
\marginsection{triangle-alert}{Important limits}
\begin{itemize}
\item This is not an official EDAC publication and is not endorsed by EDAC.
\item It uses original explanations and purpose-built mock questions and links readers to EDAC for authoritative requirements and preparation resources.
\item Mock questions use constructed study scenarios and figures developed from the guide's subject matter.
\item Exam procedures, fees, contacts, delivery platforms, and eligibility rules can change. Verify them with EDAC before relying on them.
\item Provincial and territorial planning law varies. Treat the planning chapter as a conceptual framework and verify the law in the relevant jurisdiction.
\item Examples are fictional and are provided for learning only.
\end{itemize}


\needspace{5\baselineskip}
\marginsection{external-link}{Learn more and references}
Each substantive chapter ends with a \textbf{Learn more} section that connects its concepts to relevant public sources. The \hyperref[chap:glossary]{glossary} pairs every term with its definition and public sources, while \hyperref[chap:sources]{14-public-sources-and-disclaimer.md} collects official examination, data, and contemporary-issue sources.


Definitions in this guide are plain-language study syntheses rather than quotations. When professional conventions differ, especially in financial ratios, SMART wording, pipeline stages, planning terminology, and evaluation language, the guide identifies the convention or asks the reader to state it explicitly.


\needspace{5\baselineskip}
\marginsection{info}{Current exam information and public sources}
The online requirements were checked on August 12, 2026. Candidates should confirm current eligibility and examination rules through EDAC's \href{https://edac.ca/professional-development/ecd-designation/}{Ec.D. Designation} and \href{https://exam.edac.ca/exam-prep/}{Exam Preparation} pages. See the \hyperref[chap:sources]{public sources and disclaimer} for detailed references, interpretation notes, and limitations.


\needspace{5\baselineskip}
\marginsection{sparkles}{AI assistance acknowledgement}
Xiao prepared the underlying study notes and made the final decisions about the guide's scope, content, and wording. OpenAI's Codex, using the GPT-5.6 model, helped turn those notes into a public study guide, research additional public sources, edit and format the Markdown and LaTeX files, and prepare the PDF edition.


AI-assisted content can contain errors or omissions. Xiao reviewed the guide and takes responsibility for its published content. Readers should continue to verify examination requirements and substantive information with EDAC and the cited public sources.


\needspace{5\baselineskip}
\marginsection{user-round}{About the author and feedback}
This guide was created by \href{https://www.linkedin.com/in/xxiao27/}{Xiao}, who publishes projects on GitHub as \href{https://github.com/B1ackCat7}{B1ackCat7}. Constructive feedback, corrections, comments, and professional connections are welcome. The easiest way to reach Xiao is through \href{https://www.linkedin.com/in/xxiao27/}{LinkedIn}.


\needspace{5\baselineskip}
\section[Contributing]{Contributing}
Corrections that improve accuracy, Canadian relevance, accessibility, or source quality are welcome through GitHub or \href{https://www.linkedin.com/in/xxiao27/}{LinkedIn}. Contributions must be original, respect examination integrity and copyright, and contain no confidential or personal information. Unless stated otherwise, contributions accepted into this repository are published under the same licence as the guide.


\needspace{5\baselineskip}
\marginsection{badge-check}{Licence and attribution}
Copyright \textcopyright{} 2026 \href{https://www.linkedin.com/in/xxiao27/}{Xiao}.


This guide is licensed under the \href{https://creativecommons.org/licenses/by/4.0/}{Creative Commons Attribution 4.0 International License} (CC BY 4.0). You may share and adapt the guide for any purpose, including commercially, provided you give appropriate credit, link to the licence, and indicate whether changes were made.


Suggested attribution:


\begin{guidequote}
\emph{EDAC Ec.D. Exam Study Guide} by \href{https://www.linkedin.com/in/xxiao27/}{Xiao}, published on GitHub by \href{https://github.com/B1ackCat7}{B1ackCat7} and licensed under \href{https://creativecommons.org/licenses/by/4.0/}{CC BY 4.0}.

\end{guidequote}


See \hyperref[chap:license]{LICENSE.md} for the full repository licensing notice. The licence applies only to original content in this repository; third-party names, marks, linked resources, and quoted or otherwise separately identified material remain subject to their respective rights.


\mainmatter
\part{Exam Preparation}
\chapter[Exam Overview]{\chaptericon{compass}Exam Overview}
\label{chap:exam-overview}

\needspace{5\baselineskip}
\section[Current public eligibility requirements]{Current public eligibility requirements}
As of August 12, 2026, EDAC's public \href{https://edac.ca/professional-development/ecd-designation/}{Ec.D. Designation page} and \href{https://exam.edac.ca/exam-prep/}{Exam Preparation page} present the following baseline route:


\begin{enumerate}
\item Be an EDAC member in good standing.
\item Have at least three years of economic development experience.
\item Accumulate 45 professional-development points and have eligibility verified by EDAC.
\item Complete the certification process and pass the examination with at least 75\%.
\end{enumerate}


The designation page also describes alternative academic and experienced-professional routes. Requirements depend on the program completed, experience, and other conditions. After certification, annual membership and re-certification every three years are required to continue using the designation.


Use the certification and exam-preparation pages as starting points, not as an individual eligibility decision. Ask EDAC to confirm that a specific degree, program, experience record, and point total qualify. The public pages contain some inconsistent administrative details, so also confirm the current contact, fee, meeting platform, permitted materials, accommodation process, and scheduling instructions directly with EDAC.


\needspace{5\baselineskip}
\section[Published exam structure]{Published exam structure}
\begin{center}
\small
\setlength{\LTpre}{0.4em}
\setlength{\LTpost}{0.8em}
\begin{longtable}{>{\RaggedRight\arraybackslash}p{0.120\linewidth}>{\RaggedRight\arraybackslash}p{0.430\linewidth}>{\RaggedRight\arraybackslash}p{0.180\linewidth}>{\RaggedRight\arraybackslash}p{0.150\linewidth}}
\rowcolor{DeepNavy}
\color{white}\bfseries Section & \color{white}\bfseries Published format & \color{white}\bfseries Recommended maximum time & \color{white}\bfseries Marks \\
\endfirsthead
\rowcolor{DeepNavy}
\color{white}\bfseries Section & \color{white}\bfseries Published format & \color{white}\bfseries Recommended maximum time & \color{white}\bfseries Marks \\
\endhead
1 & Multiple choice; answer every question & 40 minutes & 12 \\
2(a) & Contemporary essay; maximum 3,000 words & 70 minutes & 15 \\
2(b) & Eight short-answer sections, A through H; choose one of three questions in each section & 160 minutes & 64 \\
\textbf{Published written subtotal} &  & \textbf{270 minutes} & \textbf{91} \\
\end{longtable}
\end{center}
The multiple-choice section awards 0.5 mark per question, which implies 24 questions. The short-answer section gives 8 marks to each of eight answers. Point form is permitted for the short answers, but the response must still answer every part of the prompt.


For the online examination, EDAC currently lists a webcam-enabled computer, a modern browser, stable internet, screen sharing, and Zoom. The three published written sections above total 270 minutes and 91 raw marks; the public pages do not clearly explain the current normalization or full-session timing. Confirm the scoring, schedule, breaks, and technical setup directly with EDAC because its public pages are not fully synchronized.


\needspace{5\baselineskip}
\section[Published competency areas]{Published competency areas}
EDAC's current designation page identifies three pillars:


\begin{itemize}
\item business retention and expansion;
\item investment attraction; and
\item strategic planning.
\end{itemize}


For study purposes, this guide organizes the broader professional knowledge base into seven technical and practical skill areas:


\begin{enumerate}
\item Economic Development Processes
\begin{itemize}
\item trade and investment
\item business retention and expansion
\end{itemize}

\item Governance
\begin{itemize}
\item performance measures
\item priority setting
\item operational structures
\end{itemize}

\item Operations
\begin{itemize}
\item mentoring
\item partnerships
\item networking
\end{itemize}

\item Communications
\begin{itemize}
\item presentations
\item media relations
\item communications plans
\end{itemize}

\item Planning
\begin{itemize}
\item strategic planning
\item municipal plans
\end{itemize}

\item Marketing/Sales
\begin{itemize}
\item marketing plans
\item sales plans
\end{itemize}

\item Finance
\begin{itemize}
\item financial analysis
\item budgeting
\item impact analysis
\item business planning
\end{itemize}

\end{enumerate}


The public exam page does not label the eight short-answer sections or map them to named topics. Prepare across the full professional knowledge base, with particular fluency in the three current pillars. Treat business attraction and business retention and expansion as distinct applications.


\needspace{5\baselineskip}
\section[Candidate agreement and result process]{Candidate agreement and result process}
EDAC's \href{https://exam.edac.ca/policies/}{Policies page} says:


\begin{itemize}
\item exam marks are final and are provided only to the candidate;
\item an unsuccessful candidate can re-write later and receives an overall critique of strengths and weaknesses;
\item re-write arrangements and fees apply;
\item passing the exam is followed by EDAC Board approval; and
\item the designation should be used only after formal notification from EDAC.
\end{itemize}


Read the current policy before writing. Confirm current fees and administrative procedures directly with EDAC.


\needspace{5\baselineskip}
\section[The exam is testing applied judgment]{The exam is testing applied judgment}
Each part of the exam tests a different ability:


\begin{itemize}
\item \textbf{Recall and discrimination:} recognize the best concept or process quickly.
\item \textbf{Analysis:} connect a contemporary issue to local sectors, firms, workers, public finances, and community well-being.
\item \textbf{Professional application:} propose a realistic process with stakeholders, resources, risks, and measures.
\item \textbf{Communication:} explain the local economy clearly, accurately, and ethically under time pressure.
\end{itemize}


Knowledge matters, but a list of definitions is rarely enough. A dependable way to structure a response is:


\begin{guidequote}
Context -\textgreater{} objective -\textgreater{} evidence -\textgreater{} stakeholders -\textgreater{} options -\textgreater{} recommendation -\textgreater{} implementation -\textgreater{} measures -\textgreater{} risks and contingency

\end{guidequote}


\needspace{5\baselineskip}
\section[Items to confirm with EDAC]{Items to confirm with EDAC}
At least two weeks before the exam, confirm:


\begin{itemize}
\item the exact date, time zone, and delivery method;
\item current contact and proctor details;
\item exact fee and payment status;
\item permitted notes, reference materials, calculator, and scratch paper;
\item break rules;
\item accessibility accommodations;
\item identification requirements;
\item browser, webcam, screen-sharing, and meeting software requirements;
\item what happens after a technical interruption; and
\item how the published written marks are normalized and how the full session is scheduled.
\end{itemize}


\needspace{5\baselineskip}
\marginsection{external-link}{Learn more}
\begin{itemize}
\item \href{https://edac.ca/professional-development/ecd-designation/}{EDAC --- Ec.D. Designation}: public eligibility routes and designation expectations.
\item \href{https://exam.edac.ca/exam-prep/}{EDAC --- Exam Preparation}: current public exam format, preparation guidance, timing, and technical information.
\item \href{https://exam.edac.ca/policies/}{EDAC --- Exam Policies}: candidate agreement, result process, re-write rules, and administrative requirements.
\item \href{https://edac.ca/about/code-of-ethics/}{EDAC --- Code of Ethics}: professional standards that should inform answers throughout the examination.
\end{itemize}


EDAC is the authority for current requirements. If its public pages differ, ask EDAC to confirm the rule that applies to you.


\chapter[Study Plan and Exam Strategy]{\chaptericon{calendar-check}Study Plan and Exam Strategy}
\label{chap:study-plan}

\needspace{5\baselineskip}
\section[Build an evidence bank first]{Build an evidence bank first}
Exam answers are stronger when they draw on professional experience. Prepare a private evidence bank, and do not commit it to a public repository.


For each example, record:


\begin{center}
\small
\setlength{\LTpre}{0.4em}
\setlength{\LTpost}{0.8em}
\begin{longtable}{>{\RaggedRight\arraybackslash}p{0.290\linewidth}>{\RaggedRight\arraybackslash}p{0.610\linewidth}}
\rowcolor{DeepNavy}
\color{white}\bfseries Field & \color{white}\bfseries Notes to prepare \\
\endfirsthead
\rowcolor{DeepNavy}
\color{white}\bfseries Field & \color{white}\bfseries Notes to prepare \\
\endhead
Situation & Community, sector, problem, and why it mattered \\
Evidence & Data, research, engagement, and assumptions \\
Role & What you personally did and what others did \\
Stakeholders & Internal, external, rights-holders, decision-makers, and affected groups \\
Decision & Options considered and recommendation made \\
Delivery & Resources, timeline, governance, and communication \\
Result & Outputs, outcomes, and what was measured \\
Reflection & What worked, what did not, ethics, equity, and what you would change \\
\end{longtable}
\end{center}
Remove client names and confidential details from any study notes that could be exposed during screen sharing.


\needspace{5\baselineskip}
\section[Four-week plan]{Four-week plan}
\needspace{5\baselineskip}
\subsection[Week 1: Foundations and economic development processes]{Week 1: Foundations and economic development processes}
\begin{itemize}
\item Read the exam overview and current EDAC policies.
\item Learn the role of the economic developer and the ethics principles.
\item Compare attraction, retention and expansion, entrepreneurship, community development, and sector development.
\item Build one business attraction example, one BR\&E example, and one economic shock response.
\item Complete two timed 20-minute short answers.
\end{itemize}


\needspace{5\baselineskip}
\subsection[Week 2: Governance, operations, and communications]{Week 2: Governance, operations, and communications}
\begin{itemize}
\item Compare municipal, arm's-length, regional, and advisory structures.
\item Practise a logic model and SMART measures.
\item Review mentoring, partnerships, networking, and training.
\item Draft a communications plan and a crisis-media response.
\item Complete two short-answer drills and compare them with the quality checklist.
\end{itemize}


\needspace{5\baselineskip}
\subsection[Week 3: Planning, marketing, and sales]{Week 3: Planning, marketing, and sales}
\begin{itemize}
\item Learn the seven-step strategic planning cycle.
\item Review the relationship between economic development, an official or municipal plan, and land-use regulation.
\item Build a community profile and site-selection response checklist.
\item Practise location quotient and shift-share interpretation.
\item Write one 70-minute contemporary essay.
\end{itemize}


\needspace{5\baselineskip}
\subsection[Week 4: Finance, integration, and simulation]{Week 4: Finance, integration, and simulation}
\begin{itemize}
\item Review the three financial statements and the main ratio families.
\item Practise budgeting, business planning, and development impact analysis.
\item Complete formula drills without notes.
\item Complete one timed written simulation using the purpose-built mock exam questions and current EDAC instructions.
\item Review weak areas, then repeat a shorter mixed simulation.
\item Refresh current events during the final week.
\end{itemize}


\needspace{5\baselineskip}
\section[Two-week compressed plan]{Two-week compressed plan}
\begin{center}
\small
\setlength{\LTpre}{0.4em}
\setlength{\LTpost}{0.8em}
\begin{longtable}{>{\RaggedRight\arraybackslash}p{0.290\linewidth}>{\RaggedRight\arraybackslash}p{0.610\linewidth}}
\rowcolor{DeepNavy}
\color{white}\bfseries Day & \color{white}\bfseries Focus \\
\endfirsthead
\rowcolor{DeepNavy}
\color{white}\bfseries Day & \color{white}\bfseries Focus \\
\endhead
1 & Exam structure, ethics, role of the EDO \\
2 & Attraction, trade, investment \\
3 & BR\&E, entrepreneurship, crisis response \\
4 & Governance, structures, priority setting \\
5 & Performance measurement and logic models \\
6 & Operations, partnerships, networking \\
7 & Communications, media, public engagement \\
8 & Strategic and land-use planning \\
9 & Marketing, community profiles, site selection \\
10 & Finance and statements \\
11 & Ratios, budgeting, impact analysis \\
12 & Contemporary essay and mixed question practice \\
13 & Full simulation \\
14 & Error review, light recall, technical check \\
\end{longtable}
\end{center}
\needspace{5\baselineskip}
\section[Daily study loop]{Daily study loop}
Use a 60- to 90-minute loop:


\begin{enumerate}
\item \textbf{Recall for 10 minutes:} write what you remember without notes.
\item \textbf{Learn for 20 to 30 minutes:} read one focused section.
\item \textbf{Apply for 20 minutes:} answer one scenario in point form.
\item \textbf{Check for 10 minutes:} compare the response to the competency checklist.
\item \textbf{Explain for 5 minutes:} teach the concept aloud in plain language.
\end{enumerate}


Spaced recall and application are more useful than repeatedly rereading the same chapter.


\needspace{5\baselineskip}
\section[Time strategy]{Time strategy}
\needspace{5\baselineskip}
\subsection[Multiple choice: 40 minutes]{Multiple choice: 40 minutes}
\begin{itemize}
\item Average available time is about 100 seconds per question.
\item Answer every question.
\item First pass: answer clear questions and flag uncertain ones.
\item Second pass: eliminate options that conflict with ethics, evidence, role clarity, or a complete process.
\item Do not spend several minutes defending one half-mark answer.
\end{itemize}


\needspace{5\baselineskip}
\subsection[Contemporary essay: 70 minutes]{Contemporary essay: 70 minutes}
\begin{itemize}
\item 5 minutes: choose the issue and define the local context.
\item 8 minutes: outline the thesis, impact channels, sectors, response, partners, and measures.
\item 50 minutes: write.
\item 7 minutes: check completeness, logic, and clarity.
\end{itemize}


Build the essay in this order:


\begin{enumerate}
\item event and thesis;
\item local exposure and baseline;
\item impact pathways;
\item affected sectors, firms, workers, and residents;
\item actions already taken or recommended;
\item partners, resources, and communication;
\item measures, risks, equity, and next steps.
\end{enumerate}


\needspace{5\baselineskip}
\subsection[Eight short answers: 160 minutes]{Eight short answers: 160 minutes}
Budget 20 minutes per answer:


\begin{itemize}
\item 2 minutes to choose one of the three questions and identify every verb;
\item 15 minutes to write a structured point-form answer; and
\item 3 minutes to check that all parts were answered.
\end{itemize}


Do not borrow time from later answers unless you have a clear recovery plan. Eight complete answers usually outperform several polished answers and one unfinished response.


\needspace{5\baselineskip}
\section[Decoding question verbs]{Decoding question verbs}
\begin{center}
\small
\setlength{\LTpre}{0.4em}
\setlength{\LTpost}{0.8em}
\begin{longtable}{>{\RaggedRight\arraybackslash}p{0.290\linewidth}>{\RaggedRight\arraybackslash}p{0.610\linewidth}}
\rowcolor{DeepNavy}
\color{white}\bfseries Verb & \color{white}\bfseries What the marker needs \\
\endfirsthead
\rowcolor{DeepNavy}
\color{white}\bfseries Verb & \color{white}\bfseries What the marker needs \\
\endhead
Identify & A relevant list, usually concise \\
Describe & Characteristics or sequence with enough detail to understand \\
Explain & Cause, reason, or mechanism \\
Compare & Similarities, differences, and implications \\
Analyze & Components, relationships, evidence, and consequences \\
Evaluate & Criteria, benefits, risks, trade-offs, and judgment \\
Recommend & A decision, rationale, implementation, and measures \\
Develop & A coherent plan with roles, sequence, resources, and controls \\
\end{longtable}
\end{center}
Circle or note every instruction. A prompt that says "identify stakeholders and explain their roles" needs both names and roles.


\needspace{5\baselineskip}
\section[Short-answer quality checklist]{Short-answer quality checklist}
Before submitting, ask:


\begin{itemize}
\item Did I answer the exact question?
\item Did I organize the response around a process?
\item Did I name specific stakeholders and assign roles?
\item Did I use current and defensible evidence?
\item Did I address confidentiality, conflicts, public interest, and inclusion?
\item Did I distinguish municipal authority from partner authority?
\item Did I include resources, timing, and responsibility?
\item Did I define outputs and outcomes?
\item Did I acknowledge risks and a contingency?
\item Did I connect the concept to practice?
\end{itemize}


\needspace{5\baselineskip}
\marginsection{external-link}{Learn more}
\begin{itemize}
\item \href{https://exam.edac.ca/exam-prep/}{EDAC --- Exam Preparation}: use the current published structure to shape timed practice.
\item \href{https://exam.edac.ca/policies/}{EDAC --- Exam Policies}: confirm current candidate, technical, result, and re-write requirements.
\item \href{https://www.canada.ca/en/treasury-board-secretariat/services/audit-evaluation/evaluation-government-canada/evaluation-101-backgrounder.html}{Treasury Board of Canada Secretariat --- Evaluation 101}: a useful primer for organizing answers around inputs, activities, outputs, and outcomes.
\item \hyperref[chap:glossary]{Glossary with definition sources}: review unfamiliar terminology and follow the external references before practising.
\end{itemize}


\chapter[Foundations and Ethics]{\chaptericon{scale}Foundations and Ethics}
\label{chap:foundations}

\needspace{5\baselineskip}
\section[A working definition]{A working definition}
Economic development is a deliberate, evidence-informed process that improves a community's economic well-being and resilience. It can include job creation and retention, income and productivity growth, entrepreneurship, investment, tax-base development, infrastructure, workforce development, sector development, and quality of life.


No single program works everywhere. The appropriate strategy depends on local assets, constraints, institutions, geography, rights, risks, and community priorities.


\needspace{5\baselineskip}
\section[Role of the economic developer]{Role of the economic developer}
An economic developer may act as:


\begin{itemize}
\item analyst and researcher;
\item adviser to public decision-makers;
\item business navigator and problem solver;
\item facilitator and consensus builder;
\item connector among firms, governments, institutions, and community organizations;
\item project manager;
\item investment marketer and salesperson;
\item advocate for a competitive and fair business environment;
\item communicator and spokesperson, subject to organizational policy;
\item steward of public resources and confidential information; and
\item ethical ambassador for the profession.
\end{itemize}


The role is usually to enable and coordinate decisions, not to make every decision. Land-use approvals, budgets, grants, procurement, incentives, and corporate investments each have their own authority.


\needspace{5\baselineskip}
\section[Core functions]{Core functions}
Economic development work commonly includes:


\begin{itemize}
\item business attraction;
\item business retention and expansion;
\item business support and entrepreneurship;
\item business advocacy;
\item community development;
\item sector development;
\item communication and collaboration;
\item marketing;
\item research and analysis; and
\item ethical professional conduct.
\end{itemize}


These functions should reinforce one another. For example, BR\&E intelligence can reveal workforce gaps, guide a sector strategy, improve a community profile, and generate qualified attraction leads.


\needspace{5\baselineskip}
\section[Public value and the triple lens]{Public value and the triple lens}
A recommendation should consider:


\begin{itemize}
\item \textbf{Economic:} jobs, incomes, productivity, investment, diversification, tax base, and business viability.
\item \textbf{Social:} inclusion, accessibility, housing, services, health, distributional effects, and community cohesion.
\item \textbf{Environmental:} land, water, emissions, climate risk, infrastructure resilience, and cumulative effects.
\end{itemize}


The goal is development that is lawful, ethical, financially sustainable, aligned with community objectives, and resilient over time, rather than growth at any cost.


\needspace{5\baselineskip}
\section[EDAC Code of Ethics: study themes]{EDAC Code of Ethics: study themes}
Read the current \href{https://edac.ca/about/code-of-ethics/}{EDAC Code of Ethics} in full. The following is a study summary, not a replacement:


\begin{itemize}
\item protect the reputation of the profession and serve the public interest;
\item maintain competent and current professional knowledge;
\item compete fairly and reject deception;
\item disclose or avoid interests and relationships that impair, or appear to impair, objectivity;
\item protect confidential information and never exploit it for personal advantage;
\item act honestly and with integrity;
\item improve the profession through research and knowledge sharing;
\item treat colleagues and participants with respect;
\item demonstrate good citizenship;
\item reject bullying, harassment, and sexual harassment;
\item follow EDAC bylaws, policies, and professional rules; and
\item accept that breaches may lead to discipline.
\end{itemize}


\needspace{5\baselineskip}
\section[Common ethical scenarios]{Common ethical scenarios}
\needspace{5\baselineskip}
\subsection[Confidential expansion]{Confidential expansion}
A firm reveals an unannounced expansion and possible land purchase.


Response priorities:


\begin{enumerate}
\item clarify what is confidential and who is authorized to receive it;
\item store and share information on a need-to-know basis;
\item explain public-sector approval and disclosure limits honestly;
\item avoid promising an outcome outside your authority;
\item separate the firm's private information from aggregate BR\&E reporting; and
\item document commitments and follow-up.
\end{enumerate}


\needspace{5\baselineskip}
\subsection[Conflict of interest]{Conflict of interest}
An EDO or board member has a personal relationship with a developer.


Response priorities:


\begin{enumerate}
\item disclose the actual, potential, or perceived conflict;
\item follow the organization's conflict policy;
\item recuse the affected person where required;
\item preserve an independent, documented decision process; and
\item avoid preferential access or information.
\end{enumerate}


\needspace{5\baselineskip}
\subsection[Incentive competition]{Incentive competition}
A prospect requests an incentive while comparing multiple communities.


Response priorities:


\begin{itemize}
\item apply a published, lawful, and consistent policy;
\item assess incremental public benefit and opportunity cost;
\item consider performance conditions, clawbacks, and monitoring;
\item avoid misleading claims or disparaging competing communities;
\item protect negotiating information while meeting public accountability requirements; and
\item recommend only what produces acceptable risk-adjusted public value.
\end{itemize}


\needspace{5\baselineskip}
\subsection[Pressure to overstate data]{Pressure to overstate data}
A decision-maker asks for a stronger number than the evidence supports.


Response priorities:


\begin{itemize}
\item report the best available data and its source date;
\item state assumptions, range, and limitations;
\item separate observed results from estimates and projections;
\item offer scenarios rather than false precision; and
\item document professional advice.
\end{itemize}


\needspace{5\baselineskip}
\section[Ethical decision framework]{Ethical decision framework}
Use this sequence:


\begin{enumerate}
\item \textbf{Facts:} What is known, unknown, confidential, or disputed?
\item \textbf{Authority:} Who has the legal and organizational mandate?
\item \textbf{Stakeholders and rights-holders:} Who benefits, bears risk, or has rights that must be respected?
\item \textbf{Duties:} What do law, policy, contract, and the code require?
\item \textbf{Options:} What reasonable courses of action exist?
\item \textbf{Impacts:} What are the economic, social, environmental, fiscal, and reputational effects?
\item \textbf{Fairness:} Is the process consistent, transparent where possible, and free from improper influence?
\item \textbf{Decision:} What option best serves the public interest while respecting legitimate private interests?
\item \textbf{Record:} What rationale, approvals, conditions, and follow-up must be documented?
\end{enumerate}


\needspace{5\baselineskip}
\section[Exam application]{Exam application}
Apply ethics throughout the exam, not only in a standalone ethics answer. Connect it naturally to other topics:


\begin{itemize}
\item BR\&E: confidentiality and aggregated reporting;
\item site selection: accurate data and fair dealing;
\item governance: conflicts, accountability, and role clarity;
\item partnerships: procurement, risk allocation, and public interest;
\item communications: accuracy, privacy, and authorized spokespersons;
\item engagement: inclusion, transparency, and clear influence over the decision;
\item finance: defensible assumptions and disclosure of limitations; and
\item performance: honest attribution and reporting of both successes and failures.
\end{itemize}


\needspace{5\baselineskip}
\marginsection{external-link}{Learn more}
\begin{itemize}
\item \href{https://edac.ca/about/code-of-ethics/}{EDAC --- Code of Ethics}: the primary professional reference for this chapter.
\item \href{https://www.tbs-sct.canada.ca/pol/doc-eng.aspx?id=25049}{Treasury Board of Canada Secretariat --- Values and Ethics Code for the Public Sector}: public-interest, integrity, stewardship, and conflict principles.
\item \href{https://ised-isde.canada.ca/site/corporate-social-responsibility/en/sme-sustainability-roadmap}{Innovation, Science and Economic Development Canada --- SME Sustainability Roadmap}: economic, environmental, and social dimensions of responsible business.
\item \href{https://justice.canada.ca/eng/cons/def.html}{Justice Canada --- Consultation terminology}: legal context for Indigenous consultation and rights-holder terminology.
\end{itemize}


For definitions of conflict of interest, equity, rights-holder, stakeholder, and triple bottom line, see the \hyperref[chap:glossary]{glossary}.


\part{Professional Practice}
\chapter[Economic Development Processes]{\chaptericon{workflow}Economic Development Processes}
\label{chap:processes}

\needspace{5\baselineskip}
\section[Portfolio approach]{Portfolio approach}
Most communities need a mix of:


\begin{itemize}
\item retention and expansion of existing firms;
\item entrepreneurship and small business support;
\item attraction of compatible investment;
\item workforce and talent development;
\item sector and cluster development;
\item export development and market diversification;
\item infrastructure and site readiness;
\item downtown, rural, tourism, or community development as locally relevant; and
\item research, policy, and advocacy.
\end{itemize}


Attraction is visible, but existing firms and entrepreneurs often provide more immediate intelligence and attainable opportunities. Avoid relying on an unverified percentage claim about where "most" jobs come from. Use current local and national evidence.


\needspace{5\baselineskip}
\section[Business attraction and investment]{Business attraction and investment}
Business attraction is the process of identifying and converting investment opportunities that fit the community.


\needspace{5\baselineskip}
\subsection[Readiness before promotion]{Readiness before promotion}
Confirm:


\begin{itemize}
\item available, correctly zoned land and buildings;
\item ownership, price, constraints, and development timeline;
\item infrastructure capacity and upgrade responsibility;
\item transportation and digital connectivity;
\item labour force, occupations, wages, commuting, and training capacity;
\item supply chain and market access;
\item permitting and approval processes;
\item taxes, fees, and lawful incentives;
\item environmental, archaeological, and servicing considerations;
\item housing, health, education, and other quality-of-life factors;
\item community support and compatibility; and
\item a fast, confidential response process.
\end{itemize}


\needspace{5\baselineskip}
\subsection[Targeting process]{Targeting process}
\begin{enumerate}
\item Establish community goals and screening criteria.
\item Analyze local strengths, gaps, risks, and sector trends.
\item Identify sectors and business functions that fit the assets.
\item Select geographic markets and decision-makers.
\item Develop an evidence-based value proposition.
\item Choose lead-generation and relationship tactics.
\item Qualify prospects before spending heavily.
\item Prepare sites, data, partners, and response protocols.
\item Manage the prospect through inquiry, site visit, due diligence, decision, and implementation.
\item Track conversion, investment, outcomes, and aftercare.
\end{enumerate}


\needspace{5\baselineskip}
\subsection[Qualifying a lead]{Qualifying a lead}
Ask:


\begin{itemize}
\item Is there a verified investment project?
\item What business function, scale, capital, employment, and timing are involved?
\item What are the must-have and preferred criteria?
\item Who is the decision-maker and who represents the prospect?
\item Is confidentiality required?
\item What locations are being considered?
\item What information or action is needed next?
\item Does the project fit community objectives and risk tolerance?
\end{itemize}


Distinguish:


\begin{itemize}
\item \textbf{Contact:} a person or organization in the network.
\item \textbf{Lead:} a possible opportunity with limited validation.
\item \textbf{Prospect:} a target or lead that merits active follow-up.
\item \textbf{Qualified project:} a verified investment with criteria, authority, and timing.
\item \textbf{Completed investment:} a project that has reached the locally defined completion milestone.
\end{itemize}


\needspace{5\baselineskip}
\section[Trade and global investment]{Trade and global investment}
Global events reach a community through both inputs and outputs.


\needspace{5\baselineskip}
\subsection[Impact channels]{Impact channels}
\begin{itemize}
\item export demand and market access;
\item tariffs and non-tariff barriers;
\item imported input price and availability;
\item exchange rates;
\item commodity and energy prices;
\item interest rates and access to capital;
\item foreign direct investment;
\item tourism and travel flows;
\item labour and immigration;
\item technology and productivity;
\item transportation and logistics; and
\item confidence and expectations.
\end{itemize}


\needspace{5\baselineskip}
\subsection[Analysis sequence]{Analysis sequence}
\begin{enumerate}
\item Map local firms, sectors, employment, ownership, inputs, customers, and markets.
\item Identify exposure to the event.
\item Separate immediate, medium-term, and structural effects.
\item Identify winners, losers, and distributional effects.
\item Validate with firms, workers, institutions, and current data.
\item Match problems to the correct municipal, provincial, territorial, federal, or private response.
\item Communicate clearly without overstating certainty.
\item Monitor leading indicators and adapt.
\end{enumerate}


Currency effects are not one-directional. A lower Canadian dollar may support some exporters and Canadian tourism while raising the cost of imported inputs. Analyze the actual business model.


\needspace{5\baselineskip}
\section[Business retention and expansion]{Business retention and expansion}
BR\&E is an ongoing, confidential relationship and intelligence process designed to help existing firms remain competitive, resolve barriers, expand, and invest.


\needspace{5\baselineskip}
\subsection[Objectives]{Objectives}
Short-term objectives:


\begin{itemize}
\item build trusted relationships;
\item create or improve a business database;
\item identify urgent issues and expansion opportunities;
\item improve two-way communication; and
\item coordinate timely support.
\end{itemize}


Long-term objectives:


\begin{itemize}
\item improve firm competitiveness and resilience;
\item retain and create quality employment;
\item strengthen supply chains and sector ecosystems;
\item improve the local business climate;
\item diversify the economy; and
\item inform strategy, policy, infrastructure, workforce, and attraction.
\end{itemize}


\needspace{5\baselineskip}
\subsection[Ten-step BR\&E process]{Ten-step BR\&E process}
\begin{enumerate}
\item Define objectives, scope, privacy rules, and success measures.
\item Establish leadership, delivery roles, and a resource network.
\item Build and segment the business universe.
\item Select firms using a transparent method and a useful sector/size mix.
\item Design and test a concise survey and interview guide.
\item Train interviewers in listening, confidentiality, escalation, and record keeping.
\item Conduct visits and distinguish urgent issues from strategic intelligence.
\item Triage individual issues and analyze aggregate patterns.
\item Develop, resource, and implement an action plan.
\item Report aggregate findings, close the feedback loop, measure outcomes, and continue regular visits.
\end{enumerate}


\needspace{5\baselineskip}
\subsection[Issue triage]{Issue triage}
Use categories such as:


\begin{itemize}
\item immediate risk of closure or job loss;
\item expansion or investment opportunity;
\item workforce;
\item financing;
\item market or export development;
\item property, infrastructure, or utilities;
\item permitting and regulation;
\item succession and ownership transition;
\item innovation, technology, or productivity;
\item supply chain; and
\item referral or information request.
\end{itemize}


Assign an owner, next action, due date, confidentiality level, and status. Do not promise that every issue will be resolved.


\needspace{5\baselineskip}
\subsection[Confidentiality controls]{Confidentiality controls}
\begin{itemize}
\item explain the purpose and use of information before the interview;
\item obtain appropriate consent;
\item restrict access to firm-level records;
\item report only aggregated or sufficiently anonymized findings;
\item avoid elected-official participation in visits unless the firm is comfortable;
\item separate service records from public reporting; and
\item follow applicable privacy, records, and freedom-of-information requirements.
\end{itemize}


\needspace{5\baselineskip}
\section[Entrepreneurship and business support]{Entrepreneurship and business support}
Common interventions include:


\begin{itemize}
\item navigation and referrals;
\item business planning and financial literacy;
\item mentoring and peer networks;
\item incubators, accelerators, and shared space;
\item market access and procurement readiness;
\item financing and investor connections;
\item succession planning;
\item innovation, commercialization, and productivity support;
\item digital adoption; and
\item regulatory and permit navigation.
\end{itemize}


Match the intervention to the business stage. A start-up, a scaling firm, a distressed business, and a succession candidate require different support.


\needspace{5\baselineskip}
\section[Sector and cluster development]{Sector and cluster development}
A sector strategy should:


\begin{enumerate}
\item define the sector and geography;
\item map firms, institutions, workers, suppliers, customers, infrastructure, and research assets;
\item measure concentration, growth, productivity, exports, and risk;
\item engage sector participants;
\item identify shared constraints and opportunities;
\item prioritize interventions;
\item assign partners, resources, and timing; and
\item monitor outcomes and revisit assumptions.
\end{enumerate}


An industry can have a high location quotient but be declining. It can have a low location quotient but be growing quickly. Combine quantitative measures with firm intelligence and community priorities.


\needspace{5\baselineskip}
\section[Economic shock and recovery framework]{Economic shock and recovery framework}
For a closure, disaster, trade shock, infrastructure failure, or public-health event:


\begin{enumerate}
\item activate the appropriate emergency or corporate structure;
\item define the event, geography, sectors, and time horizon;
\item establish a baseline and rapid intelligence process;
\item identify affected businesses, workers, residents, and critical services;
\item communicate verified information and available supports;
\item stabilize urgent needs;
\item coordinate workforce, finance, supply-chain, and market responses;
\item develop recovery and diversification actions;
\item monitor leading and lagging indicators; and
\item capture lessons and improve resilience plans.
\end{enumerate}


The EDO coordinates economic recovery but does not replace emergency management, public health, social services, or other authorities.


\needspace{5\baselineskip}
\section[Measures]{Measures}
Track a mix of:


\begin{itemize}
\item firms contacted and response rate;
\item issues identified, resolved, and time to resolution;
\item expansions and investments supported;
\item jobs retained and created, with method and confidence level;
\item referrals and partner contributions;
\item firms entering new markets;
\item productivity, innovation, and training outcomes;
\item closures, layoffs, and risk signals;
\item client satisfaction and trust; and
\item policy, infrastructure, or program changes resulting from aggregate intelligence.
\end{itemize}


\needspace{5\baselineskip}
\marginsection{external-link}{Learn more}
\begin{itemize}
\item \href{https://www.ontario.ca/page/business-retention-and-expansion-program}{Ontario --- Business Retention and Expansion Program}: a Canadian public-sector BR\&E framework.
\item \href{https://www.worldbank.org/en/topic/investment-climate/brief/investment-policy-and-promotion}{World Bank --- Investment Policy and Promotion}: investment attraction, facilitation, retention, and expansion.
\item \href{https://www.oecd.org/en/publications/oecd-benchmark-definition-of-foreign-direct-investment-fifth-edition\_7f05c0a3-en/full-report/main-concepts-and-definitions-of-foreign-direct-investment\_c3240e7d.html}{OECD --- Benchmark Definition of Foreign Direct Investment}: concepts and definitions for FDI.
\item \href{https://www.eda.gov/about/economic-development-glossary}{U.S. Economic Development Administration --- Economic Development Glossary}: clusters, resilience, and related economic-development concepts.
\item \href{https://www.eda.gov/resources/comprehensive-economic-development-strategy/content/economic-resilience}{U.S. Economic Development Administration --- Economic Resilience}: preparedness, diversification, response, and recovery.
\end{itemize}


Pipeline terms such as \emph{lead}, \emph{prospect}, and \emph{project} vary among organizations. Define each stage and qualification rule before reporting conversion.


\chapter[Governance and Performance]{\chaptericon{chart-no-axes-combined}Governance and Performance}
\label{chap:governance}

\needspace{5\baselineskip}
\section[Governance is an operating system]{Governance is an operating system}
Governance defines:


\begin{itemize}
\item mandate and strategic direction;
\item authority and accountability;
\item board, council, committee, and staff roles;
\item decision and reporting processes;
\item funding and risk oversight;
\item conflicts and ethical controls; and
\item relationships with partners and the public.
\end{itemize}


A structure without a clear mandate, strategy, resources, authority, and reporting relationship will underperform even if it has talented people.


\needspace{5\baselineskip}
\section[Common structures]{Common structures}
\begin{center}
\small
\setlength{\LTpre}{0.4em}
\setlength{\LTpost}{0.8em}
\begin{longtable}{>{\RaggedRight\arraybackslash}p{0.180\linewidth}>{\RaggedRight\arraybackslash}p{0.400\linewidth}>{\RaggedRight\arraybackslash}p{0.300\linewidth}}
\rowcolor{DeepNavy}
\color{white}\bfseries Structure & \color{white}\bfseries Potential strengths & \color{white}\bfseries Common risks \\
\endfirsthead
\rowcolor{DeepNavy}
\color{white}\bfseries Structure & \color{white}\bfseries Potential strengths & \color{white}\bfseries Common risks \\
\endhead
Municipal department & Stable access to corporate services, policy alignment, direct accountability & Political cycles, slower decisions, unclear boundaries with other departments \\
Arm's-length corporation or agency & Business-led governance, specialized delivery, potential operating flexibility & Overhead, funding insecurity, accountability or procurement ambiguity \\
Regional partnership & Larger asset base, shared cost, broader labour and supply-chain view & Multiple funders, competing priorities, lead-sharing disputes, complex reporting \\
Advisory committee & Low cost, community input, political buffer, support to staff & Limited authority, volunteer capacity, unclear role, risk of operational interference \\
\end{longtable}
\end{center}
There is no universally correct model. Select a structure based on goals, scale, geography, legislation, funding, partner capacity, and the need for confidentiality and speed.


\needspace{5\baselineskip}
\section[Principles for a sound structure]{Principles for a sound structure}
\begin{enumerate}
\item Include the groups needed for the economic development process.
\item Coordinate the major economic development efforts in the community.
\item Include formal decision-makers and credible community leaders.
\item Give the governing body enough attention and authority to act within its mandate.
\item Balance continuity with renewal, diversity, and succession.
\item Meet regularly and maintain active participation.
\item Provide adequate funding, a clear mandate, strategy, and terms of reference.
\item Do not consume staff capacity that should serve clients and deliver the strategy.
\end{enumerate}


Sound governance also covers accessibility, inclusion, Indigenous rights and relationships, risk oversight, data governance, and transparent conflict-of-interest controls.


\needspace{5\baselineskip}
\section[Board, council, and staff roles]{Board, council, and staff roles}
\needspace{5\baselineskip}
\subsection[Governing body]{Governing body}
\begin{itemize}
\item approves mandate, strategy, policy, budget, and risk tolerance;
\item appoints or supervises senior leadership as applicable;
\item monitors outcomes and financial stewardship;
\item manages conflicts and accountability; and
\item represents the organization's legitimacy.
\end{itemize}


\needspace{5\baselineskip}
\subsection[Economic development staff]{Economic development staff}
\begin{itemize}
\item provide evidence and professional advice;
\item design and deliver programs;
\item manage clients, projects, data, and partners;
\item implement approved strategy and budget;
\item report performance and risk; and
\item maintain professional competence and ethics.
\end{itemize}


The board governs; staff manage and deliver. Individual board members should not direct staff or promise client outcomes unless formally authorized.


\needspace{5\baselineskip}
\section[Priority setting]{Priority setting}
Use a transparent scoring method. Possible criteria:


\begin{itemize}
\item strategic fit;
\item evidence of need or opportunity;
\item economic, social, and environmental impact;
\item equity and inclusion;
\item feasibility and organizational capacity;
\item partner readiness;
\item cost and funding;
\item risk and reversibility;
\item time to benefit;
\item additionality;
\item community support; and
\item ability to measure results.
\end{itemize}


Document the weighting and assumptions. A scoring tool supports judgment; it does not replace it.


\needspace{5\baselineskip}
\section[Logic model]{Logic model}
\begin{center}
\small
\setlength{\LTpre}{0.4em}
\setlength{\LTpost}{0.8em}
\begin{longtable}{>{\RaggedRight\arraybackslash}p{0.180\linewidth}>{\RaggedRight\arraybackslash}p{0.400\linewidth}>{\RaggedRight\arraybackslash}p{0.300\linewidth}}
\rowcolor{DeepNavy}
\color{white}\bfseries Element & \color{white}\bfseries Question & \color{white}\bfseries Example \\
\endfirsthead
\rowcolor{DeepNavy}
\color{white}\bfseries Element & \color{white}\bfseries Question & \color{white}\bfseries Example \\
\endhead
Inputs & What resources are used? & Staff time, budget, data, partner expertise \\
Activities & What do we do? & Firm visits, workshops, site preparation \\
Outputs & What is delivered? & Visits completed, firms served, sites profiled \\
Short-term outcomes & What changes first? & Better intelligence, resolved barriers, new skills \\
Intermediate outcomes & What changes next? & Firm expansion, investment decisions, market entry \\
Long-term outcomes & What public value results? & Resilience, quality jobs, incomes, diversification \\
\end{longtable}
\end{center}
Do not confuse activity with impact. Ten meetings are an output; a resolved workforce barrier is an outcome.


\needspace{5\baselineskip}
\section[SMART measures]{SMART measures}
Measures should be:


\begin{itemize}
\item \textbf{Specific:} clearly defined;
\item \textbf{Measurable:} supported by a reliable collection method;
\item \textbf{Achievable:} realistic but meaningful;
\item \textbf{Relevant:} connected to the mandate and strategy; and
\item \textbf{Time-bound:} attached to a reporting period or deadline.
\end{itemize}


Also ask whether the measure is attributable, timely, comparable, affordable to collect, privacy-safe, and resistant to gaming.


\needspace{5\baselineskip}
\section[Performance framework]{Performance framework}
Track measures across four levels:


\needspace{5\baselineskip}
\subsection[Activity and service]{Activity and service}
\begin{itemize}
\item inquiries and response time;
\item business visits;
\item prospects and qualified projects;
\item workshops and participants;
\item referrals and issue-resolution cases.
\end{itemize}


\needspace{5\baselineskip}
\subsection[Conversion and effectiveness]{Conversion and effectiveness}
\begin{itemize}
\item contact-to-project conversion;
\item project-to-investment conversion;
\item expansions completed;
\item issues resolved;
\item client satisfaction;
\item cost per qualified project or completed outcome.
\end{itemize}


\needspace{5\baselineskip}
\subsection[Economic outcomes]{Economic outcomes}
\begin{itemize}
\item capital investment;
\item jobs retained or created;
\item wage or job-quality measures;
\item assessment or tax-base change;
\item business starts, survival, and closures;
\item exports, productivity, or market diversification.
\end{itemize}


\needspace{5\baselineskip}
\subsection[Community outcomes]{Community outcomes}
\begin{itemize}
\item economic diversification;
\item inclusion and distribution of benefits;
\item housing and infrastructure capacity;
\item downtown occupancy;
\item resilience and climate readiness;
\item resident or business confidence.
\end{itemize}


\needspace{5\baselineskip}
\section[Attribution and additionality]{Attribution and additionality}
Economic results have many causes. Interest rates, exchange rates, business cycles, firm decisions, and other governments can dominate local efforts.


Report:


\begin{itemize}
\item \textbf{Contribution:} what the organization did that plausibly helped.
\item \textbf{Attribution:} what can reasonably be assigned to the program.
\item \textbf{Additionality:} what likely would not have occurred, or would have occurred later or at smaller scale, without the intervention.
\item \textbf{Deadweight:} benefits that likely would have happened anyway.
\item \textbf{Displacement:} activity moved from elsewhere rather than newly created.
\item \textbf{Leakage:} benefits that flow outside the intended community.
\end{itemize}


Use confidence levels or ranges rather than claiming every announced job as fully caused by the EDO.


\needspace{5\baselineskip}
\section[Evaluation cycle]{Evaluation cycle}
\begin{enumerate}
\item Define the decision the evaluation must inform.
\item Confirm the program theory and baseline.
\item Choose measures and collection responsibilities.
\item Establish data quality, privacy, and reporting rules.
\item Monitor implementation.
\item Compare results with target, baseline, and relevant peers.
\item Explain variance and external factors.
\item Decide whether to continue, change, scale, or stop.
\item Communicate results honestly.
\end{enumerate}


\needspace{5\baselineskip}
\section[Governance answer pattern]{Governance answer pattern}
Structure a governance response in this order:


\begin{enumerate}
\item diagnose the mandate and accountability problem;
\item map authorities and stakeholders;
\item present viable structure options;
\item compare benefits, risks, cost, and speed;
\item recommend a structure and transition;
\item define board and staff roles;
\item add funding, conflict, confidentiality, and reporting controls; and
\item define outcomes and a review date.
\end{enumerate}


\needspace{5\baselineskip}
\marginsection{external-link}{Learn more}
\begin{itemize}
\item \href{https://www.canada.ca/en/treasury-board-secretariat/services/audit-evaluation/evaluation-government-canada/evaluation-101-backgrounder.html}{Treasury Board of Canada Secretariat --- Evaluation 101}: logic models, theories of change, and results chains.
\item \href{https://www.tbs-sct.canada.ca/pol/doc-eng.aspx?id=31300\&section=glossary}{Treasury Board of Canada Secretariat --- Policy on Results glossary}: Canadian federal performance and evaluation terminology.
\item \href{https://www.oecd.org/en/publications/glossary-of-key-terms-in-evaluation-and-results-based-management-for-sustainable-development-second-edition\_632da462-en-fr-es.html}{OECD --- Glossary of Key Terms in Evaluation and Results-Based Management}: internationally used definitions of activities, outputs, outcomes, impacts, and evaluation.
\item \href{https://www.gov.uk/government/publications/the-green-book-appraisal-and-evaluation-in-central-government/the-green-book-2026}{UK Government --- The Green Book}: additionality, deadweight, displacement, leakage, appraisal, and evaluation.
\item \href{https://www.canada.ca/en/treasury-board-secretariat/services/guidance-crown-corporations/guiding-principles-management-crown-corporations.html}{Treasury Board of Canada Secretariat --- Guiding Principles for the Management of Crown Corporations}: arm's-length governance and accountability.
\end{itemize}


SMART wording varies across sources. This guide uses \emph{specific, measurable, achievable, relevant, and time-bound}; identify the version you use rather than treating another established variant as automatically wrong.


\chapter[Operations, Partnerships, and Networking]{\chaptericon{users}Operations, Partnerships, and Networking}
\label{chap:operations}

\needspace{5\baselineskip}
\section[Mentoring]{Mentoring}
Mentoring pairs experience with a development need. It should support capability and judgment rather than create dependency.


\needspace{5\baselineskip}
\subsection[Common forms]{Common forms}
\begin{itemize}
\item \textbf{Casual:} learning from observed behaviour without a formal arrangement.
\item \textbf{Informal:} a relationship that develops naturally.
\item \textbf{Structured, non-facilitated:} mentor and mentee set goals and expectations directly.
\item \textbf{Facilitated:} a coordinator recruits, matches, supports, and evaluates pairs.
\item \textbf{Group:} mentors support a group with shared needs.
\item \textbf{Multiple mentor:} a mentee works with different mentors for different capabilities.
\end{itemize}


\needspace{5\baselineskip}
\subsection[Program design]{Program design}
\begin{enumerate}
\item Define target users, needs, and outcomes.
\item Recruit suitable mentors and screen conflicts.
\item Set boundaries, confidentiality, duration, and escalation rules.
\item Match on goals, sector, experience, and working style.
\item Create a learning agreement and meeting rhythm.
\item Provide orientation and resources.
\item Check progress without intruding on confidential discussion.
\item Recognize contributions and evaluate outcomes.
\end{enumerate}


Mentoring is not therapy, professional legal/accounting advice, a guarantee of financing, or a power relationship. Refer specialized issues to qualified professionals.


\needspace{5\baselineskip}
\section[Partnerships]{Partnerships}
Economic development normally requires municipalities, other governments, Indigenous communities and organizations, firms, sector groups, educational institutions, workforce agencies, utilities, financial institutions, non-profits, and residents.


\needspace{5\baselineskip}
\subsection[Partnership test]{Partnership test}
Before forming a partnership, ask:


\begin{itemize}
\item Is there a shared problem or opportunity?
\item Does each partner add authority, assets, expertise, trust, reach, or funding?
\item Are interests compatible enough to act?
\item Are roles, decisions, resources, risks, data, and public communication clear?
\item Is a partnership better than a contract, referral, consultation, or informal network?
\item How will success, disputes, change, and exit be handled?
\end{itemize}


\needspace{5\baselineskip}
\subsection[Partnership agreement checklist]{Partnership agreement checklist}
\begin{itemize}
\item purpose, scope, and intended outcomes;
\item partners and authorized representatives;
\item governance and decision rules;
\item roles and deliverables;
\item financial and in-kind contributions;
\item timeline and milestones;
\item intellectual property and data;
\item privacy and confidentiality;
\item communications and branding;
\item procurement and conflict rules;
\item risk allocation and insurance;
\item measurement and reporting;
\item dispute resolution, amendment, and termination.
\end{itemize}


A RACI table can clarify delivery roles:


\begin{itemize}
\item \textbf{Responsible:} does the work.
\item \textbf{Accountable:} owns the result and approves.
\item \textbf{Consulted:} provides input before the action.
\item \textbf{Informed:} receives updates.
\end{itemize}


Only one role should normally be accountable for a deliverable.


\needspace{5\baselineskip}
\section[Public-private partnerships]{Public-private partnerships}
A public-private partnership allocates the design, construction, financing, operation, maintenance, or ownership of a public asset or service among public and private parties.


Potential benefits:


\begin{itemize}
\item specialized expertise;
\item lifecycle integration;
\item schedule discipline;
\item access to capital;
\item risk transfer where the private partner can manage the risk;
\item innovation and service improvements.
\end{itemize}


Potential risks:


\begin{itemize}
\item loss of public control;
\item higher financing or transaction cost;
\item weak competition;
\item inflexible long-term terms;
\item unclear accountability;
\item optimistic demand assumptions;
\item risk nominally transferred but ultimately borne by taxpayers;
\item inequitable access or user fees.
\end{itemize}


A sound evaluation includes public-interest objectives, full lifecycle cost, value-for-money comparison, procurement integrity, risk allocation, performance standards, affordability, transparency, contract management capacity, and exit provisions.


Government remains accountable for protecting the public interest even when delivery is contracted.


\needspace{5\baselineskip}
\section[Stakeholder and rights-holder mapping]{Stakeholder and rights-holder mapping}
Map:


\begin{itemize}
\item decision authority;
\item legal or constitutional rights;
\item degree of impact;
\item interest and influence;
\item knowledge and resources;
\item trust and relationship history;
\item preferred form of engagement;
\item potential conflict or alignment.
\end{itemize}


Indigenous governments and communities are rights-holders, not merely another stakeholder category. Determine the applicable legal duties, jurisdiction, protocols, and relationship expectations early.


\needspace{5\baselineskip}
\section[Networking]{Networking}
Networking builds working relationships. Collecting contacts alone is not enough.


\needspace{5\baselineskip}
\subsection[Effective cycle]{Effective cycle}
\begin{enumerate}
\item Define the purpose and the people relevant to it.
\item Research their role and interests.
\item Offer a clear, useful introduction.
\item Listen for needs and possible mutual value.
\item Record commitments with appropriate privacy.
\item Follow up promptly.
\item Provide value before asking for repeated help.
\item Maintain the relationship over time.
\end{enumerate}


Build connections across:


\begin{itemize}
\item local business and sector networks;
\item provincial and territorial associations;
\item EDAC and other professional groups;
\item regional development agencies;
\item site selectors and real-estate professionals;
\item workforce and education partners;
\item trade and investment representatives;
\item utilities, ports, airports, and transportation providers;
\item Indigenous economic development organizations; and
\item peers in comparable communities.
\end{itemize}


Track relationship quality and outcomes, not only event attendance.


\needspace{5\baselineskip}
\section[Training and continuous development]{Training and continuous development}
An EDO needs:


\begin{itemize}
\item current economic and business knowledge;
\item research, data, and financial skills;
\item facilitation and engagement;
\item project and risk management;
\item writing, media, presentation, and negotiation;
\item land-use and municipal process knowledge;
\item privacy, ethics, accessibility, and inclusion;
\item digital, cybersecurity, and AI literacy; and
\item sector-specific knowledge.
\end{itemize}


A personal development plan should record:


\begin{center}
\small
\setlength{\LTpre}{0.4em}
\setlength{\LTpost}{0.8em}
\begin{longtable}{>{\RaggedRight\arraybackslash}p{0.140\linewidth}>{\RaggedRight\arraybackslash}p{0.140\linewidth}>{\RaggedRight\arraybackslash}p{0.140\linewidth}>{\RaggedRight\arraybackslash}p{0.140\linewidth}>{\RaggedRight\arraybackslash}p{0.140\linewidth}>{\RaggedRight\arraybackslash}p{0.140\linewidth}}
\rowcolor{DeepNavy}
\color{white}\bfseries Capability & \color{white}\bfseries Current level & \color{white}\bfseries Evidence & \color{white}\bfseries Gap & \color{white}\bfseries Action & \color{white}\bfseries Deadline \\
\endfirsthead
\rowcolor{DeepNavy}
\color{white}\bfseries Capability & \color{white}\bfseries Current level & \color{white}\bfseries Evidence & \color{white}\bfseries Gap & \color{white}\bfseries Action & \color{white}\bfseries Deadline \\
\endhead
Example: financial analysis & Working & Can interpret statements & Ratio speed & Weekly drills & Month end \\
\end{longtable}
\end{center}
EDAC currently requires annual membership and re-certification every three years to maintain use of the Ec.D. designation. Verify current requirements on the \href{https://edac.ca/professional-development/ec-d-recertification/}{Re-certification page}.


\needspace{5\baselineskip}
\section[Downtown redevelopment stakeholder example]{Downtown redevelopment stakeholder example}
Possible participants and roles:


\begin{itemize}
\item council: mandate, budget, and policy decisions;
\item economic development: business intelligence, coordination, investment, and implementation;
\item planning/building: policy, zoning, permits, design, and compliance;
\item finance: fiscal analysis and funding;
\item engineering/public works: infrastructure and phasing;
\item communications: public information and media;
\item Indigenous rights-holders: rights, interests, knowledge, and partnership;
\item property and business owners: operational realities, investment, and delivery;
\item residents and community organizations: needs, impacts, ideas, and legitimacy;
\item accessibility, housing, heritage, environmental, and social-service groups: specialized impacts;
\item developers, investors, and lenders: feasibility, capital, and construction;
\item utilities and transportation providers: capacity and network changes;
\item educational institutions and workforce partners: talent and programming.
\end{itemize}


The response should also explain engagement method, decision authority, conflict management, phasing, business-disruption mitigation, financing, measures, and communication.


\needspace{5\baselineskip}
\marginsection{external-link}{Learn more}
\begin{itemize}
\item \href{https://justice.canada.ca/eng/rp-pr/other-autre/rr12\_1/p1.html}{Justice Canada --- Three Years On: Mentoring}: mentoring as a learning relationship and organizational practice.
\item \href{https://www.infrastructureontario.ca/en/what-we-do/capital-delivery/faqs---public-private-partnerships-p3s/}{Infrastructure Ontario --- Public-private partnerships frequently asked questions}: Canadian P3 delivery and risk-allocation context.
\item \href{https://justice.canada.ca/eng/cons/def.html}{Justice Canada --- Consultation terminology}: consultation and accommodation terminology in the Indigenous rights context.
\item \href{https://our-languages.canada.ca/en/blogue-blog/rethinking-the-word-stakeholder-eng}{Language Portal of Canada --- Rethinking the word ``stakeholder''}: why Indigenous Peoples should not be reduced to a stakeholder category.
\item \href{https://edac.ca/professional-development/ec-d-recertification/}{EDAC --- Ec.D. Re-certification}: verify current continuing-membership and re-certification requirements.
\end{itemize}


Stakeholder engagement does not replace legal duties, government-to-government relationships, or distinctions among Indigenous nations and rights-holding collectives.


\chapter[Communications and Public Engagement]{\chaptericon{messages-square}Communications and Public Engagement}
\label{chap:communications}

\needspace{5\baselineskip}
\section[Communication plan]{Communication plan}
Plan the communication before choosing channels. Cover:


\begin{enumerate}
\item \textbf{Situation:} decision, issue, context, evidence, and sensitivities.
\item \textbf{Objectives:} what the audience should know, feel, or do.
\item \textbf{Audiences:} affected groups, rights-holders, decision-makers, partners, media, and internal staff.
\item \textbf{Research:} public environment, stakeholder intelligence, media scan, and misinformation risks.
\item \textbf{Messages:} a few accurate, audience-relevant points with supporting facts.
\item \textbf{Channels:} direct contact, web, email, meeting, media, social, print, or partner channels.
\item \textbf{Timing:} sequence, dependencies, approvals, and response windows.
\item \textbf{Roles:} author, approver, spokesperson, monitor, and responder.
\item \textbf{Accessibility and language:} format, disability access, literacy, official languages, translation, and cultural relevance.
\item \textbf{Risk and response:} confidentiality, legal constraints, controversy, errors, and escalation.
\item \textbf{Measurement:} reach, understanding, participation, sentiment, action, and response time.
\end{enumerate}


Ask:


\begin{itemize}
\item Are the decision and degree of public influence clear?
\item Are departments and partners coordinated?
\item Are the messages accurate, defensible, and approved?
\item Does each priority audience have a suitable channel?
\item Can people respond or ask questions?
\item Are privacy, records, and accessibility requirements met?
\end{itemize}


\needspace{5\baselineskip}
\section[Message map]{Message map}
Use one central message supported by three points:


\begin{guidequote}
\textbf{Central message:} What happened, why it matters, and what is being done.

\end{guidequote}


\begin{itemize}
\item Supporting point 1: evidence and context.
\item Supporting point 2: action and responsibility.
\item Supporting point 3: timing, support, or next step.
\end{itemize}


Prepare proof, examples, and likely questions for each supporting point.


\needspace{5\baselineskip}
\section[Media relations]{Media relations}
\needspace{5\baselineskip}
\subsection[Before an interview]{Before an interview}
\begin{itemize}
\item follow the organization's spokesperson and approval policy;
\item clarify outlet, reporter, topic, format, deadline, and other sources;
\item identify the audience and likely angle;
\item verify facts and source dates;
\item prepare three key messages and difficult questions;
\item determine what is confidential or outside your authority;
\item arrange an appropriate setting and technical setup; and
\item decide who will follow up with supporting material.
\end{itemize}


Assume anything said may be used. Do not rely on "off the record" to protect confidential information.


\needspace{5\baselineskip}
\subsection[During an interview]{During an interview}
\begin{itemize}
\item be accurate, direct, respectful, and concise;
\item lead with the most important point;
\item use plain language and define technical terms;
\item answer the question, then bridge to the key message where appropriate;
\item correct a false premise calmly;
\item say when you do not know and commit only to follow-up you can deliver;
\item explain when law, privacy, or commercial confidentiality prevents comment;
\item do not speculate; and
\item stop when the answer is complete.
\end{itemize}


\needspace{5\baselineskip}
\subsection[After an interview]{After an interview}
\begin{itemize}
\item provide promised information promptly;
\item record material commitments;
\item notify affected partners before coverage triggers inquiries;
\item monitor for factual errors;
\item request a factual correction calmly when necessary; and
\item evaluate whether the message was understood.
\end{itemize}


\needspace{5\baselineskip}
\section[Crisis communication]{Crisis communication}
For a closure, layoff, disaster, cyber incident, trade shock, or controversial project:


\begin{enumerate}
\item activate the authorized response structure;
\item verify the known facts and unknowns;
\item identify affected groups and urgent needs;
\item coordinate one source of truth;
\item state what happened, what is being done, and when the next update will occur;
\item protect personal, security, and commercial information;
\item monitor misinformation and operational feedback;
\item update regularly, even when little has changed; and
\item transition from incident information to recovery support.
\end{enumerate}


Empathy is not an admission of liability. Acknowledge impact on people and businesses.


\needspace{5\baselineskip}
\section[Written communication]{Written communication}
Effective reports and briefings:


\begin{itemize}
\item put the decision or conclusion first;
\item distinguish background from analysis;
\item use headings that match the reader's questions;
\item cite data and dates;
\item explain assumptions and uncertainty;
\item compare options using consistent criteria;
\item make the recommendation explicit;
\item assign actions, timing, resources, and measures; and
\item use accessible plain language.
\end{itemize}


A council or board report often follows:


\begin{guidequote}
Purpose -\textgreater{} recommendation -\textgreater{} background -\textgreater{} analysis and options -\textgreater{} financial/risk implications -\textgreater{} engagement -\textgreater{} implementation -\textgreater{} attachments

\end{guidequote}


\needspace{5\baselineskip}
\section[Presentations]{Presentations}
\needspace{5\baselineskip}
\subsection[Format]{Format}
\begin{itemize}
\item one-to-one: use dialogue and adapt to the listener;
\item small group: involve participants and manage differing interests;
\item large group: create a clear narrative and maintain attention;
\item online: use vocal signposting, shorter segments, and deliberate interaction.
\end{itemize}


\needspace{5\baselineskip}
\subsection[Audience]{Audience}
\begin{itemize}
\item receptive: do not assume agreement; still provide evidence;
\item challenging: answer questions and demonstrate reasoning;
\item hostile: remain calm, acknowledge concerns, enforce respectful rules, and do not overpromise.
\end{itemize}


\needspace{5\baselineskip}
\subsection[Structure]{Structure}
\begin{enumerate}
\item opening: purpose and why it matters;
\item roadmap: what will be covered;
\item evidence: a few relevant facts;
\item analysis: implications and trade-offs;
\item recommendation or request;
\item next steps;
\item conclusion that repeats the main point.
\end{enumerate}


The goal is audience understanding and retention, not display of every fact collected.


\needspace{5\baselineskip}
\section[Public engagement]{Public engagement}
Providing information alone does not create engagement.


\needspace{5\baselineskip}
\subsection[Levels]{Levels}
\begin{itemize}
\item \textbf{Inform:} a decision is made; the purpose is clear, accessible explanation.
\item \textbf{Consult:} gather feedback before a decision.
\item \textbf{Involve:} participants help define issues, criteria, or options.
\item \textbf{Collaborate:} participants share development of recommendations.
\item \textbf{Empower:} formal authority is delegated where law permits.
\end{itemize}


State what is open to influence and what is not. Do not invite participation if the decision is already fixed and present the session as engagement.


\needspace{5\baselineskip}
\subsection[Engagement design]{Engagement design}
\begin{enumerate}
\item define the decision and authority;
\item identify affected groups and rights-holders;
\item assess barriers, trust, and engagement history;
\item choose the appropriate level and methods;
\item provide balanced information early enough to use;
\item make participation accessible and safe;
\item record input accurately;
\item analyze themes without erasing minority views;
\item show how input affected the decision; and
\item evaluate reach, representation, quality, and trust.
\end{enumerate}


Possible methods include interviews, surveys, focus groups, workshops, roundtables, advisory groups, public meetings, online tools, and community-led processes.


\needspace{5\baselineskip}
\section[Digital and social channels]{Digital and social channels}
Platform popularity changes. Choose channels based on audience and purpose rather than memorizing platform statistics.


Controls should cover:


\begin{itemize}
\item authorized accounts and administrator access;
\item records retention;
\item privacy and consent;
\item accessibility and captions;
\item moderation and respectful conduct;
\item security and multi-factor authentication;
\item response time and escalation;
\item separation of personal and organizational accounts;
\item account continuity when staff leave;
\item misinformation and correction protocol;
\item measurement tied to objectives.
\end{itemize}


Social media complements a maintained website and direct relationships. It does not replace a reliable source of official information.


\needspace{5\baselineskip}
\section[Communications answer pattern]{Communications answer pattern}
Build a communications response around:


\begin{enumerate}
\item identify the issue and authority;
\item establish objective and audiences;
\item gather and verify facts;
\item develop messages;
\item choose channels and engagement level;
\item assign approvals and spokespersons;
\item address confidentiality, accessibility, and risk;
\item set timing and feedback loops; and
\item measure understanding and action.
\end{enumerate}


\needspace{5\baselineskip}
\marginsection{external-link}{Learn more}
\begin{itemize}
\item \href{https://www.canada.ca/en/health-canada/services/publications/health-system-services/health-canada-public-health-agency-canada-guidelines-public-engagement.html}{Health Canada and Public Health Agency of Canada --- Guidelines on Public Engagement}: planning, conducting, reporting, and evaluating engagement.
\item \href{https://www.canada.ca/en/health-canada/services/publications/health-system-services/health-canada-public-health-agency-canada-guidelines-public-engagement/glossary-terms.html}{Health Canada and Public Health Agency of Canada --- Public engagement glossary}: engagement and stakeholder terminology.
\item \href{https://justice.canada.ca/eng/cons/def.html}{Justice Canada --- Consultation terminology}: distinguish public engagement from the Crown's legal consultation obligations.
\item \href{https://accessible.canada.ca/standards-and-technical-guides/standards-and-technical-guides-database/can-asc-312025-plain-language?mode=full-html}{Accessibility Standards Canada --- CAN-ASC-3.1:2025 Plain Language}: principles and requirements for communication that audiences can find, understand, and use.
\item \href{https://edac.ca/about/code-of-ethics/}{EDAC --- Code of Ethics}: accuracy, confidentiality, integrity, and professional conduct.
\end{itemize}


Use \emph{stakeholder} for affected or interested parties where appropriate, but not as a substitute for \emph{Indigenous rights-holder}.


\chapter[Planning and Land Use]{\chaptericon{map}Planning and Land Use}
\label{chap:planning}

\needspace{5\baselineskip}
\section[Strategic economic planning]{Strategic economic planning}
A strategy aligns evidence, priorities, resources, actions, and measurement.


\needspace{5\baselineskip}
\subsection[Seven-step cycle]{Seven-step cycle}
\begin{enumerate}
\item Set broad goals and objectives.
\item Complete a SWOT or comparable evidence-based assessment.
\item Identify strategic alternatives.
\item Select specific objectives and priorities.
\item Develop and implement action plans.
\item Develop contingencies.
\item Monitor, evaluate, learn, and update.
\end{enumerate}


The process is iterative. A final document is the beginning of implementation, not the end of planning.


\needspace{5\baselineskip}
\section[Evidence for the strategy]{Evidence for the strategy}
Analyze:


\begin{itemize}
\item population, migration, households, and demographics;
\item labour force, occupations, skills, wages, and commuting;
\item businesses, sectors, ownership, and supply chains;
\item exports, imports, and investment;
\item entrepreneurship and innovation;
\item land, buildings, infrastructure, utilities, and broadband;
\item housing, transportation, health, education, and services;
\item municipal finance and development capacity;
\item climate, environmental, and hazard risk;
\item Indigenous rights, interests, and economic relationships;
\item inclusion and distribution of opportunity;
\item existing plans, legislation, and partner strategies; and
\item perceptions gathered through engagement.
\end{itemize}


Separate internal conditions from external trends. A weakness is not automatically a threat, and a trend is not automatically an opportunity.


\needspace{5\baselineskip}
\section[From goal to action]{From goal to action}
Each action should state:


\begin{itemize}
\item objective and rationale;
\item lead and partners;
\item scope and deliverable;
\item required authority;
\item budget and staff;
\item start, milestones, and completion;
\item dependencies;
\item risks and contingency;
\item output and outcome measures;
\item reporting and review.
\end{itemize}


Avoid vague actions such as "promote the community." Specify the target audience, message, method, owner, budget, timing, and intended result.


\needspace{5\baselineskip}
\section[Common strategy failures]{Common strategy failures}
\begin{itemize}
\item predetermined conclusions disguised as analysis;
\item planning only to satisfy a funder;
\item top-down process with weak community ownership;
\item perceptions treated as facts;
\item too many priorities for available resources;
\item actions without owners or budgets;
\item no feedback and performance system;
\item artificial deadlines;
\item failure to consider implementation capacity;
\item no contingency or update process.
\end{itemize}


\needspace{5\baselineskip}
\section[Municipal and land-use plans]{Municipal and land-use plans}
Terms and legal requirements vary by province and territory. The high-level relationship is generally:


\begin{itemize}
\item a strategic or official/municipal plan expresses long-term community direction and land-use policy;
\item secondary, area, transportation, servicing, or other plans add detail;
\item zoning or land-use bylaws/regulations establish permitted and discretionary uses and development standards;
\item subdivision, site-plan, development, building, environmental, and other approvals apply to specific proposals.
\end{itemize}


Always verify the applicable statute, policy hierarchy, approval authority, appeal process, consultation duty, and Indigenous rights in the relevant jurisdiction.


\needspace{5\baselineskip}
\section[What a municipal or official plan may address]{What a municipal or official plan may address}
\begin{itemize}
\item location of residential, commercial, industrial, institutional, agricultural, resource, and protected land;
\item growth sequence and settlement boundaries;
\item employment lands;
\item transportation and mobility;
\item water, wastewater, stormwater, energy, and telecommunications;
\item housing and community services;
\item natural hazards and climate resilience;
\item environmental and cultural resources;
\item urban design and public realm;
\item regional coordination; and
\item implementation and monitoring.
\end{itemize}


\needspace{5\baselineskip}
\section[Role of the economic developer]{Role of the economic developer}
The EDO should:


\begin{itemize}
\item communicate business, workforce, and investment needs to planners;
\item understand policy, zoning, servicing, and approval pathways;
\item maintain current land and building information;
\item identify employment-land and infrastructure needs;
\item connect prospects with the correct staff without promising approval;
\item support engagement and evidence gathering;
\item assess economic and business-climate implications;
\item coordinate implementation after approval; and
\item protect orderly planning, public interest, and procedural fairness.
\end{itemize}


The professional planner leads planning analysis and statutory process. The EDO contributes economic evidence and client navigation.


\needspace{5\baselineskip}
\section[Development readiness]{Development readiness}
For a site or employment area, assess:


\begin{itemize}
\item policy and zoning status;
\item ownership and willingness to transact;
\item title, easements, contamination, heritage, archaeology, and natural hazards;
\item parcel size, configuration, topography, and expansion;
\item transportation access;
\item water, wastewater, power, gas, and telecommunications capacity;
\item approvals and realistic timeline;
\item required studies;
\item development charges, taxes, and other costs;
\item housing and workforce access;
\item adjacent uses and community compatibility;
\item climate and infrastructure resilience.
\end{itemize}


"Shovel ready" should have a defined, evidence-based meaning. Do not use it as an unsupported marketing label.


\needspace{5\baselineskip}
\section[Public and private development roles]{Public and private development roles}
\needspace{5\baselineskip}
\subsection[Publicly owned employment land]{Publicly owned employment land}
The municipality may:


\begin{itemize}
\item acquire, plan, service, phase, market, lease, or sell land;
\item set public objectives and disposition conditions;
\item use transparent valuation, procurement, and approval processes;
\item manage infrastructure and long-term fiscal risk.
\end{itemize}


\needspace{5\baselineskip}
\subsection[Privately owned development]{Privately owned development}
The municipality generally:


\begin{itemize}
\item sets policy and regulation;
\item reviews and decides applications within its authority;
\item provides or coordinates public infrastructure;
\item applies fees, conditions, standards, and incentives lawfully;
\item consults and engages as required;
\item monitors compliance.
\end{itemize}


Avoid conflicts when the municipality is both landowner/promoter and regulator. Separate roles, document decisions, and maintain procedural fairness.


\needspace{5\baselineskip}
\section[Development impact analysis]{Development impact analysis}
Begin early, before major commitments.


\needspace{5\baselineskip}
\subsection[Impact categories]{Impact categories}
\begin{itemize}
\item \textbf{Fiscal:} service costs, capital costs, taxes, fees, timing, and net municipal position.
\item \textbf{Environmental:} land, water, air, habitat, emissions, waste, and cumulative effects.
\item \textbf{Socio-economic:} jobs, wages, housing, population, equity, services, health, culture, and displacement.
\item \textbf{Transportation:} traffic, transit, goods movement, safety, active transportation, and infrastructure.
\item \textbf{Economic:} supplier effects, competition, market demand, diversification, productivity, and resilience.
\end{itemize}


\needspace{5\baselineskip}
\subsection[Process]{Process}
\begin{enumerate}
\item screen consistency with plans and law;
\item define scope, geography, baseline, time horizon, and alternatives;
\item identify affected groups and rights-holders;
\item collect project and community data;
\item analyze direct, indirect, induced, cumulative, and distributional effects as appropriate;
\item test assumptions and scenarios;
\item identify mitigation, conditions, and monitoring;
\item communicate uncertainty and trade-offs;
\item integrate findings into the decision.
\end{enumerate}


Do not count gross benefits without costs, displacement, leakage, timing, and risk.


\needspace{5\baselineskip}
\section[Downtown redevelopment]{Downtown redevelopment}
Assess the project through:


\begin{itemize}
\item business and property-owner intelligence;
\item market and leakage analysis;
\item land-use and housing policy;
\item mobility, parking, loading, and public realm;
\item heritage, accessibility, safety, and climate resilience;
\item construction disruption mitigation;
\item programming, events, and place management;
\item investment and financing tools;
\item clear governance and phasing;
\item measures such as vacancy, foot traffic, investment, housing, business survival, and user experience.
\end{itemize}


\needspace{5\baselineskip}
\marginsection{external-link}{Learn more}
\begin{itemize}
\item \href{https://www.ontario.ca/document/citizens-guide-land-use-planning/official-plans}{Ontario --- Official plans}: one Canadian jurisdiction's explanation of long-term municipal land-use policy.
\item \href{https://www.ontario.ca/document/citizens-guide-land-use-planning/zoning-bylaws}{Ontario --- Zoning bylaws}: how zoning implements land-use direction through permitted uses and standards.
\item \href{https://www.canada.ca/en/impact-assessment-agency/services/policy-guidance/practitioners-guide-impact-assessment-act/analyzing-health-social-economic-effects-impact-assessment-act.html}{Impact Assessment Agency of Canada --- Analyzing health, social and economic effects}: federal impact-assessment guidance.
\item \href{https://www.oecd.org/en/topics/sub-issues/strategic-planning-and-regional-development.html}{OECD --- Strategic planning and regional development}: place-based strategic planning and regional development.
\item \href{https://www.calgary.ca/content/dam/www/cfod/finance/documents/corporate-economics/other-reports/municipal-fiscal-impact-model.pdf}{City of Calgary --- A Municipal Fiscal Impact Model}: municipal revenue-and-cost implications of growth.
\end{itemize}


Planning law and terminology differ across provinces and territories. \emph{Official plan}, \emph{municipal plan}, \emph{zoning bylaw}, and \emph{land-use bylaw} are not universally interchangeable legal terms; use the language of the applicable jurisdiction.


\chapter[Marketing, Sales, and Site Selection]{\chaptericon{target}Marketing, Sales, and Site Selection}
\label{chap:marketing}

\needspace{5\baselineskip}
\section[Marketing versus sales]{Marketing versus sales}
\begin{itemize}
\item \textbf{Marketing} identifies target audiences, builds awareness, positions the community, and generates leads.
\item \textbf{Sales} qualifies needs, builds trust, demonstrates fit, overcomes barriers, and moves a project toward a decision.
\end{itemize}


A brand cannot compensate for weak readiness or inaccurate data. The product is the community's actual investment environment.


\needspace{5\baselineskip}
\section[Community competitive analysis]{Community competitive analysis}
\needspace{5\baselineskip}
\subsection[Purpose]{Purpose}
\begin{itemize}
\item retain existing firms;
\item support expansion;
\item stimulate entrepreneurship;
\item identify sector opportunities;
\item target compatible investment;
\item guide infrastructure, workforce, and policy.
\end{itemize}


\needspace{5\baselineskip}
\subsection[Evidence]{Evidence}
Combine:


\begin{itemize}
\item demographics and migration;
\item labour force, occupations, wages, skills, and commuting;
\item business counts and sector employment;
\item exports, imports, customers, and suppliers;
\item location quotients and shift-share;
\item real estate and infrastructure;
\item costs, taxes, utilities, and approvals;
\item innovation and institutional assets;
\item climate and resilience;
\item competitor and benchmark communities;
\item business interviews and community engagement.
\end{itemize}


Quantitative analysis identifies patterns. Local intelligence explains why they exist and whether they represent a strategic opportunity.


\needspace{5\baselineskip}
\section[Community profile]{Community profile}
A community profile should allow an investor to conduct a first-stage comparison.


\needspace{5\baselineskip}
\subsection[Core categories]{Core categories}
\begin{itemize}
\item location, maps, and transportation access;
\item population, growth, age, households, income, and migration;
\item labour force, occupations, skills, wages, unemployment, participation, and commuting;
\item sector composition, major employers where public, business counts, and supply chains;
\item available land and buildings;
\item zoning, approvals, fees, and realistic timelines;
\item utilities, capacity, cost, reliability, and broadband;
\item taxes and lawful incentives;
\item education, training, research, health, housing, and community services;
\item quality of life relevant to workforce attraction;
\item sector profiles and case evidence;
\item current contact and response commitment.
\end{itemize}


\needspace{5\baselineskip}
\subsection[Data standards]{Data standards}
For each important figure, record:


\begin{itemize}
\item definition;
\item geography;
\item reference period;
\item source;
\item last-updated date;
\item limitations;
\item owner and next refresh date.
\end{itemize}


Use third-party evidence where possible. Do not present marketing estimates as official statistics.


\needspace{5\baselineskip}
\section[Targeting and positioning]{Targeting and positioning}
\needspace{5\baselineskip}
\subsection[Target selection]{Target selection}
Score possible sectors or business functions on:


\begin{itemize}
\item existing concentration and growth;
\item supply-chain fit;
\item labour and training;
\item land, utilities, and logistics;
\item market access;
\item innovation and institutional assets;
\item community support;
\item environmental and infrastructure capacity;
\item resilience and diversification;
\item realistic competition and project pipeline.
\end{itemize}


Choose a small number of defensible targets. "Technology" is usually too broad; a specific business function and subsector give the program a clearer focus.


\needspace{5\baselineskip}
\subsection[Value proposition]{Value proposition}
A value proposition should answer:


\begin{itemize}
\item For which target?
\item What business need?
\item Why this community?
\item What evidence proves it?
\item What action should the decision-maker take next?
\end{itemize}


Use credible differentiation. Do not claim to be the best without a defined comparison.


\needspace{5\baselineskip}
\section[Investment marketing plan]{Investment marketing plan}
\begin{enumerate}
\item Research the community and prepare the data foundation.
\item Select target sectors and business functions.
\item Identify geographic markets and decision-makers.
\item Develop positioning and evidence-based messages.
\item Select integrated tactics.
\item Assign partners, budget, timing, and lead management.
\item Execute, measure, learn, and adjust.
\end{enumerate}


Tactics may include:


\begin{itemize}
\item corporate aftercare and local firm calls;
\item referral networks;
\item direct outreach;
\item site-selector and broker relationships;
\item sector events and trade shows;
\item earned media and thought leadership;
\item digital content and search;
\item familiarization or inward missions;
\item partner and senior-leader outreach.
\end{itemize}


Choose tactics based on qualified lead generation and relationship value, not visibility alone.


\needspace{5\baselineskip}
\section[Lead and project management]{Lead and project management}
\needspace{5\baselineskip}
\subsection[Suggested fields]{Suggested fields}
\begin{itemize}
\item unique project number;
\item contact and organization;
\item source and referral;
\item date opened and last contact;
\item confidentiality level;
\item sector and business function;
\item investment, employment, site, utility, labour, and timing needs;
\item decision process and decision-makers;
\item current stage and probability;
\item owner, next action, and due date;
\item partners involved;
\item barriers and risk;
\item outcome and close reason.
\end{itemize}


Use access controls and retention rules appropriate to commercial confidentiality and privacy.


\needspace{5\baselineskip}
\subsection[Funnel]{Funnel}
\begin{guidequote}
Audience -\textgreater{} contact -\textgreater{} lead -\textgreater{} qualified project -\textgreater{} proposal/site visit -\textgreater{} due diligence -\textgreater{} decision -\textgreater{} implementation -\textgreater{} aftercare

\end{guidequote}


Track conversion and time at each stage. A high volume of unqualified leads can hide a weak program.


\needspace{5\baselineskip}
\section[Site-selection response]{Site-selection response}
\needspace{5\baselineskip}
\subsection[First response]{First response}
\begin{itemize}
\item acknowledge quickly;
\item confirm deadline, format, criteria, and confidentiality;
\item assign one project lead;
\item identify internal and external contributors;
\item clarify questions early;
\item create a source-controlled response schedule.
\end{itemize}


\needspace{5\baselineskip}
\subsection[Data-room checklist]{Data-room checklist}
\begin{itemize}
\item available properties with maps and constraints;
\item land use, zoning, and approvals;
\item utilities and capacity confirmation;
\item transportation and logistics;
\item labour, wages, skills, commuting, and training;
\item taxes, fees, and incentives;
\item construction and operating costs where defensible;
\item market, supplier, and customer access;
\item environmental and geotechnical information where available;
\item housing and community services;
\item project timeline and contacts;
\item source notes and dates.
\end{itemize}


\needspace{5\baselineskip}
\subsection[Site visit]{Site visit}
\begin{enumerate}
\item confirm the prospect's objectives and confidentiality;
\item design the route around decision criteria;
\item brief participants and control the attendee list;
\item validate sites and presentation details;
\item include useful business, workforce, utility, and government contacts;
\item answer accurately and log open questions;
\item send a concise follow-up with owners and deadlines;
\item continue disciplined follow-up without over-contacting.
\end{enumerate}


\needspace{5\baselineskip}
\section[Corporate aftercare]{Corporate aftercare}
Existing firms can generate:


\begin{itemize}
\item reinvestment and expansion;
\item referrals and testimonials;
\item supplier attraction;
\item export growth;
\item intelligence on competitiveness;
\item early warning of risk.
\end{itemize}


Aftercare requires regular, confidential contact and issue resolution. It is both a BR\&E function and an investment marketing function.


\needspace{5\baselineskip}
\section[Measures]{Measures}
\begin{itemize}
\item profile data freshness;
\item target decision-makers reached;
\item qualified leads by source;
\item cost per qualified lead;
\item proposals and site visits;
\item conversion by stage and target;
\item response time;
\item investment and quality jobs;
\item reinvestment from existing firms;
\item project loss reasons;
\item partner contribution;
\item client experience.
\end{itemize}


Interpret results over a realistic multi-year cycle. Major investment decisions may take years.


\needspace{5\baselineskip}
\marginsection{external-link}{Learn more}
\begin{itemize}
\item \href{https://www.trade.gov/selectusa-investor-guide}{SelectUSA --- Investor Guide}: public-sector guidance on the investor journey, location assessment, and establishment process.
\item \href{https://www.trade.gov/selectusa-reports-and-publications}{SelectUSA --- Reports and Publications}: competitive-analysis and value-proposition resources.
\item \href{https://www.trade.gov/selectusa-target-industries-dashboard}{SelectUSA --- Target Industries Dashboard}: an example of evidence used to examine target industries.
\item \href{https://www.worldbank.org/en/topic/investment-climate/brief/investment-policy-and-promotion}{World Bank --- Investment Policy and Promotion}: investment attraction, facilitation, retention, and aftercare.
\item \href{https://www.bdc.ca/en/articles-tools/entrepreneur-toolkit/templates-business-guides/glossary/marketing-plan}{Business Development Bank of Canada --- Marketing plan}: practical marketing-plan elements.
\item \href{https://trailhead.salesforce.com/content/learn/modules/lead-qualification-quick-look/get-to-know-lead-qualification}{Salesforce --- Lead qualification}: an introduction to explicit qualification criteria.
\end{itemize}


There is no universal lead-to-project vocabulary. Define the entry and exit criteria for each funnel stage, report conversion between named stages, and avoid comparing organizations that use different definitions.


\chapter[Finance, Business Planning, and Impact Analysis]{\chaptericon{chart-column}Finance, Business Planning, and Impact Analysis}
\label{chap:finance}

\needspace{5\baselineskip}
\section[Why finance matters]{Why finance matters}
An EDO may not be the firm's accountant, lender, or investment adviser, but should be able to:


\begin{itemize}
\item understand the business model;
\item read basic statements;
\item identify questions and warning signs;
\item assess public proposals and funding requests;
\item refer specialized work appropriately;
\item connect financial assumptions to economic development outcomes.
\end{itemize}


Always confirm the accounting basis, period, definitions, and source data before comparing ratios.


\needspace{5\baselineskip}
\section[The three financial statements]{The three financial statements}
\needspace{5\baselineskip}
\subsection[Income statement]{Income statement}
The income statement shows revenue, expenses, and profit over a period.


Basic relationship:


\begin{guidequote}
Revenue - Expenses = Net income

\end{guidequote}


It answers:


\begin{itemize}
\item Is the business profitable?
\item What drives revenue?
\item What costs are fixed or variable?
\item Are margins improving?
\item Is depreciation included?
\end{itemize}


Principal repayment is not an income-statement expense; interest generally is.


\needspace{5\baselineskip}
\subsection[Balance sheet]{Balance sheet}
The balance sheet shows financial position at one date.


\begin{guidequote}
Assets = Liabilities + Equity

\end{guidequote}


In IFRS terminology, an asset is a present economic resource controlled by the entity, a liability is a present obligation to transfer an economic resource, and equity is the residual interest after liabilities are deducted from assets. The familiar language of what a business ``owns and owes'' is useful shorthand, but it is not the complete accounting definition.


It includes:


\begin{itemize}
\item current assets such as cash, receivables, and inventory;
\item long-term assets such as equipment and buildings;
\item current liabilities due within the operating or reporting definition;
\item long-term debt;
\item owner or shareholder equity.
\end{itemize}


It answers:


\begin{itemize}
\item What does the business own and owe?
\item Can it meet near-term obligations?
\item How is it financed?
\item How much working capital and leverage does it have?
\end{itemize}


\needspace{5\baselineskip}
\subsection[Cash-flow statement]{Cash-flow statement}
The cash-flow statement shows cash received and paid through:


\begin{itemize}
\item operating activities;
\item investing activities;
\item financing activities.
\end{itemize}


A profitable firm can still fail if cash is tied up in receivables or inventory, capital spending is heavy, or debt payments arrive before cash receipts.


\needspace{5\baselineskip}
\section[Ratio interpretation]{Ratio interpretation}
Use ratios to ask better questions, not to produce a verdict from one number.


Compare with:


\begin{itemize}
\item the firm's own trend;
\item plan or budget;
\item firms with comparable sector, size, age, and accounting;
\item lender or industry norms where appropriate.
\end{itemize}


Investigate:


\begin{itemize}
\item liquidity;
\item operating efficiency;
\item leverage and debt service;
\item profitability;
\item cash generation;
\item customer or supplier concentration;
\item seasonality;
\item quality of underlying assets and receivables.
\end{itemize}


See \hyperref[chap:calculations]{10-calculations-and-formulas.md} for formulas and examples.


Ratio conventions are not universal. In particular, state:


\begin{itemize}
\item whether quick assets exclude both inventory and less-liquid prepaids;
\item whether days sales outstanding uses average or ending receivables, net credit or total sales, and 360 or 365 days;
\item whether debt-to-equity uses interest-bearing debt or total liabilities; and
\item which earnings, asset, and equity measures are used for coverage, return on assets, and return on equity.
\end{itemize}


\needspace{5\baselineskip}
\section[Warning signs]{Warning signs}
\begin{itemize}
\item profit rising while operating cash flow falls;
\item receivables growing faster than sales;
\item inventory buildup without demand evidence;
\item repeated short-term borrowing for persistent losses;
\item current obligations exceeding realistic liquid assets;
\item debt service with little margin for a downturn;
\item one customer, supplier, market, or owner dominates;
\item forecasts without documented assumptions;
\item no allowance for working capital, tax, contingency, or seasonality;
\item owner withdrawals obscuring operating performance.
\end{itemize}


Do not diagnose from a single ratio. Ask for context and qualified advice.


\needspace{5\baselineskip}
\section[Budgeting]{Budgeting}
A budget translates priorities into a controlled financial plan.


\needspace{5\baselineskip}
\subsection[Operating budget]{Operating budget}
Typical items:


\begin{itemize}
\item salaries and benefits;
\item professional services;
\item marketing and communications;
\item travel and events;
\item training and memberships;
\item data, software, and subscriptions;
\item office and facility costs;
\item program grants or contributions;
\item revenues, grants, and transfers.
\end{itemize}


\needspace{5\baselineskip}
\subsection[Capital budget]{Capital budget}
Used for long-lived assets such as:


\begin{itemize}
\item land;
\item buildings;
\item infrastructure;
\item major equipment;
\item site servicing;
\item information systems.
\end{itemize}


Assess full lifecycle cost, timing, financing, debt, operating impact, maintenance, risk, and disposal.


\needspace{5\baselineskip}
\subsection[Budget cycle]{Budget cycle}
\begin{enumerate}
\item connect requests to strategy and mandate;
\item estimate revenue and expenditure;
\item document assumptions;
\item prioritize and approve;
\item assign authority and controls;
\item monitor actual versus budget;
\item explain variance;
\item forecast year-end;
\item take corrective action;
\item report and learn.
\end{enumerate}


\needspace{5\baselineskip}
\section[Business plan]{Business plan}
A business plan should explain:


\begin{itemize}
\item executive summary;
\item business concept and value proposition;
\item ownership and legal structure;
\item products or services;
\item industry and market analysis;
\item target customer;
\item competition and advantage;
\item marketing and sales;
\item operations and supply chain;
\item location and facilities;
\item management and workforce;
\item approvals, licences, and intellectual property;
\item risk and contingency;
\item capital and funding;
\item financial assumptions;
\item projected income statement, balance sheet, and cash flow;
\item milestones and measures;
\item supporting documents.
\end{itemize}


An economic development organization's annual business or work plan should also include target sectors or clients, programs, staffing, partners, marketing, budget, deliverables, performance measures, and reporting.


\needspace{5\baselineskip}
\subsection[Forecast questions]{Forecast questions}
\begin{itemize}
\item What volume, price, and growth are assumed?
\item What evidence supports market share?
\item What are fixed and variable costs?
\item When is cash collected?
\item What working capital is required?
\item What capacity constraints arise?
\item Are wages, tax, inflation, financing, maintenance, and contingency included?
\item What happens under low, base, and high scenarios?
\end{itemize}


\needspace{5\baselineskip}
\section[Development impact analysis]{Development impact analysis}
Development impact analysis estimates the consequences of a project before a decision or major commitment.


\needspace{5\baselineskip}
\subsection[Fiscal]{Fiscal}
Consider:


\begin{itemize}
\item property and other public revenues;
\item fees and grants;
\item capital infrastructure;
\item operating and maintenance costs;
\item service demand;
\item timing of receipts and costs;
\item financing and debt;
\item lifecycle and replacement;
\item risk, contingencies, and agreement conditions.
\end{itemize}


The result is not simply "new taxes minus one year's costs." Timing and lifecycle matter.


\needspace{5\baselineskip}
\subsection[Environmental]{Environmental}
Consider:


\begin{itemize}
\item land and habitat;
\item water and wastewater;
\item air, emissions, and energy;
\item contamination and waste;
\item climate hazard and resilience;
\item cumulative effects;
\item mitigation, monitoring, and liability.
\end{itemize}


\needspace{5\baselineskip}
\subsection[Socio-economic]{Socio-economic}
Consider:


\begin{itemize}
\item direct, indirect, and induced employment where methodology supports them;
\item job quality, wages, and workforce access;
\item housing and population;
\item health, education, social, and cultural services;
\item local procurement and supply-chain effects;
\item equity, inclusion, displacement, and vulnerable groups;
\item community character and public acceptance.
\end{itemize}


\needspace{5\baselineskip}
\subsection[Transportation]{Transportation}
Consider:


\begin{itemize}
\item worker and goods movement;
\item road, rail, port, airport, transit, and active modes;
\item congestion and safety;
\item capital and operating needs;
\item emissions and neighbourhood impacts;
\item construction and long-term demand.
\end{itemize}


\needspace{5\baselineskip}
\section[Economic impact versus fiscal impact]{Economic impact versus fiscal impact}
\begin{itemize}
\item \textbf{Economic impact} estimates activity in the economy, often output, income, employment, or value added.
\item \textbf{Fiscal impact} estimates the effect on government revenues and costs.
\item \textbf{Financial feasibility} tests whether the project or business can fund itself.
\item \textbf{Cost-benefit analysis} compares social benefits and costs across alternatives.
\end{itemize}


Do not treat these as interchangeable.


\needspace{5\baselineskip}
\section[Impact-analysis discipline]{Impact-analysis discipline}
\begin{enumerate}
\item Define the decision, alternatives, geography, and time horizon.
\item Establish a baseline and counterfactual.
\item State data, model, and assumptions.
\item Separate direct, indirect, and induced effects.
\item Account for displacement, leakage, deadweight, and capacity constraints.
\item Show distribution and timing, not only totals.
\item Test scenarios and sensitivity.
\item Avoid adding overlapping benefits.
\item Identify non-quantified effects.
\item Recommend conditions, mitigation, monitoring, and review.
\end{enumerate}


\needspace{5\baselineskip}
\section[Finance answer pattern]{Finance answer pattern}
Structure a financial or business case in this order:


\begin{enumerate}
\item clarify the decision and information required;
\item review the business model and three statements;
\item calculate and interpret relevant ratios;
\item examine cash flow and assumptions;
\item compare scenarios and benchmarks;
\item identify risk and information gaps;
\item assess public impact and additionality if support is requested;
\item recommend action, conditions, referrals, and monitoring.
\end{enumerate}


\needspace{5\baselineskip}
\marginsection{external-link}{Learn more}
\begin{itemize}
\item \href{https://www.ifrs.org/issued-standards/list-of-standards/conceptual-framework/?headerSearchOption=phrase}{IFRS Foundation --- Conceptual Framework for Financial Reporting}: assets, liabilities, equity, recognition, and measurement concepts.
\item \href{https://www.ifrs.org/issued-standards/list-of-standards/ias-1-presentation-of-financial-statements.html/}{IFRS Foundation --- IAS 1 Presentation of Financial Statements}: financial-statement presentation.
\item \href{https://www.ifrs.org/issued-standards/list-of-standards/ias-7-statement-of-cash-flows.html/}{IFRS Foundation --- IAS 7 Statement of Cash Flows}: operating, investing, and financing cash flows.
\item \href{https://www.bdc.ca/en/articles-tools/money-finance/manage-finances/financial-ratios-what-are-how-use}{Business Development Bank of Canada --- Financial ratios}: interpretation and business uses of common ratios.
\item \href{https://www.tbs-sct.canada.ca/rtrap-parfa/analys/analys-eng.pdf}{Treasury Board of Canada Secretariat --- Canadian Cost-Benefit Analysis Guide}: public-sector cost-benefit analysis.
\item \href{https://www.statcan.gc.ca/en/statistical-programs/document/5115\_D6\_T9\_V1}{Statistics Canada --- Guide to Using the Input-Output Simulation Model}: direct, indirect, and induced economic effects.
\item \href{https://www.calgary.ca/content/dam/www/cfod/finance/documents/corporate-economics/other-reports/municipal-fiscal-impact-model.pdf}{City of Calgary --- A Municipal Fiscal Impact Model}: municipal fiscal-impact concepts.
\end{itemize}


\part{Practice and Review}
\chapter[Calculations and Formulas]{\chaptericon{calculator}Calculations and Formulas}
\label{chap:calculations}

\needspace{5\baselineskip}
\section[Calculation habits]{Calculation habits}
For every calculation:


\begin{enumerate}
\item write the formula;
\item label units and time period;
\item substitute values;
\item calculate;
\item interpret in plain language;
\item state limitations.
\end{enumerate}


Do not interpret a result without a benchmark, trend, or context.


\needspace{5\baselineskip}
\section[Location quotient]{Location quotient}
Location quotient compares a local industry's share of a regional measure with the same industry's share of that measure in a reference economy. The calculation below uses employment, the most common exam-style version. LQ can also be calculated with earnings, output, or GDP; do not mix measures in one calculation or comparison.


\begin{Verbatim}[fontsize=\small,breaklines=true,breakanywhere=true,frame=single,framesep=2mm]
LQ = (local industry employment / total local employment)
     ----------------------------------------------------
     (reference industry employment / total reference employment)
\end{Verbatim}

Example:


\begin{itemize}
\item local industry employment = 1,200
\item total local employment = 20,000
\item reference industry employment = 40,000
\item total reference employment = 1,000,000
\end{itemize}


\begin{Verbatim}[fontsize=\small,breaklines=true,breakanywhere=true,frame=single,framesep=2mm]
Local share     = 1,200 / 20,000 = 0.06
Reference share = 40,000 / 1,000,000 = 0.04
LQ              = 0.06 / 0.04 = 1.50
\end{Verbatim}

Interpretation:


\begin{itemize}
\item LQ = 1.00: same concentration as the reference area.
\item LQ \textgreater{} 1.00: more concentrated locally.
\item LQ \textless{} 1.00: less concentrated locally.
\end{itemize}


An LQ of 1.50 means the industry's local employment share is 50\% higher than its reference share. It does not prove productivity, competitiveness, exports, or future growth.


Check:


\begin{itemize}
\item consistent industry classification and geography;
\item same year and employment concept;
\item suppression and rounding;
\item self-employment and commuting limitations;
\item whether a few firms dominate.
\end{itemize}


\needspace{5\baselineskip}
\section[Employment growth rate]{Employment growth rate}
\begin{Verbatim}[fontsize=\small,breaklines=true,breakanywhere=true,frame=single,framesep=2mm]
Growth rate = (employment at end - employment at start)
              -----------------------------------------
                       employment at start
\end{Verbatim}

Multiply by 100 for percent.


\needspace{5\baselineskip}
\section[Shift-share analysis]{Shift-share analysis}
Shift-share decomposes local industry employment change into:


\begin{itemize}
\item reference-area growth effect;
\item industry-mix effect;
\item local or differential shift effect.
\end{itemize}


Let:


\begin{itemize}
\item \texttt{e0} = local industry employment at start;
\item \texttt{e1} = local industry employment at end;
\item \texttt{g} = growth rate of total reference employment;
\item \texttt{gi} = growth rate of reference industry employment;
\item \texttt{gl} = growth rate of local industry employment.
\end{itemize}


\begin{Verbatim}[fontsize=\small,breaklines=true,breakanywhere=true,frame=single,framesep=2mm]
Reference growth effect = e0 x g
Industry mix effect     = e0 x (gi - g)
Local shift effect      = e0 x (gl - gi)
Total local change      = reference effect + mix effect + local effect
\end{Verbatim}

Example:


\begin{itemize}
\item local industry: 1,000 to 1,150, so \texttt{gl = 15\%}
\item total reference economy grows 4\%, so \texttt{g = 4\%}
\item reference industry grows 8\%, so \texttt{gi = 8\%}
\end{itemize}


\begin{Verbatim}[fontsize=\small,breaklines=true,breakanywhere=true,frame=single,framesep=2mm]
Reference effect = 1,000 x 0.04          = 40
Industry mix     = 1,000 x (0.08 - 0.04) = 40
Local shift      = 1,000 x (0.15 - 0.08) = 70
Total change     = 40 + 40 + 70           = 150
\end{Verbatim}

Interpretation:


\begin{itemize}
\item 40 jobs are associated with overall reference growth.
\item 40 are associated with the industry's stronger reference trend.
\item 70 are associated with stronger local performance relative to that industry trend.
\end{itemize}


A positive local shift suggests a local advantage but does not identify its cause. Validate with firms and other evidence.


\needspace{5\baselineskip}
\section[Revenue, cost, and profit]{Revenue, cost, and profit}
\begin{Verbatim}[fontsize=\small,breaklines=true,breakanywhere=true,frame=single,framesep=2mm]
Net sales = Gross sales - Returns - Discounts

Cost of goods sold = Beginning inventory + Purchases - Ending inventory

Gross profit = Net sales - Cost of goods sold

Net income = Revenue - Expenses
\end{Verbatim}

\needspace{5\baselineskip}
\section[Liquidity ratios]{Liquidity ratios}
\needspace{5\baselineskip}
\subsection[Current ratio]{Current ratio}
\begin{Verbatim}[fontsize=\small,breaklines=true,breakanywhere=true,frame=single,framesep=2mm]
Current ratio = Current assets / Current liabilities
\end{Verbatim}

Example: \texttt{300 / 150 = 2.00}.


The firm reports \$2.00 of current assets for each \$1.00 of current liabilities. Asset quality and timing still matter.


\needspace{5\baselineskip}
\subsection[Quick ratio]{Quick ratio}
\begin{Verbatim}[fontsize=\small,breaklines=true,breakanywhere=true,frame=single,framesep=2mm]
Quick ratio = Quick assets / Current liabilities
\end{Verbatim}

Quick assets commonly include cash, short-term investments, and receivables. A simplified exam formula is:


\begin{Verbatim}[fontsize=\small,breaklines=true,breakanywhere=true,frame=single,framesep=2mm]
Quick ratio = (Current assets - Inventory - Less-liquid prepaids)
              ---------------------------------------------------
                         Current liabilities
\end{Verbatim}

If no prepaid amount is provided, the common approximation is \texttt{(Current assets - Inventory) / Current liabilities}.


Example with no prepaids given: \texttt{(300 - 100) / 150 = 1.33}.


This is stricter than the current ratio because inventory and other less-liquid current assets are excluded. State the definition used.


\needspace{5\baselineskip}
\section[Activity ratios]{Activity ratios}
\needspace{5\baselineskip}
\subsection[Inventory turnover]{Inventory turnover}
\begin{Verbatim}[fontsize=\small,breaklines=true,breakanywhere=true,frame=single,framesep=2mm]
Inventory turnover = Cost of goods sold / Average inventory

Average inventory = (Beginning inventory + Ending inventory) / 2
\end{Verbatim}

Example: \texttt{600 / 150 = 4.0 times}.


Compare with the sector and the firm's trend. High turnover can mean efficiency or insufficient stock.


\needspace{5\baselineskip}
\subsection[Days sales outstanding]{Days sales outstanding}
\begin{Verbatim}[fontsize=\small,breaklines=true,breakanywhere=true,frame=single,framesep=2mm]
DSO = Average accounts receivable / (Annual net credit sales / Days in period)
\end{Verbatim}

Use beginning and ending receivables to calculate an average when both are available. If only period-end receivables are given, use them as a proxy and say so. If net credit sales are unavailable, total sales may be used with a limitation.


Example using 365 days and period-end receivables as a proxy:


\begin{Verbatim}[fontsize=\small,breaklines=true,breakanywhere=true,frame=single,framesep=2mm]
DSO = 100 / (900 / 365) = 40.6 days
\end{Verbatim}

Use 360 or 365 consistently with the benchmark.


\needspace{5\baselineskip}
\section[Leverage and coverage]{Leverage and coverage}
\needspace{5\baselineskip}
\subsection[Debt-to-assets]{Debt-to-assets}
\begin{Verbatim}[fontsize=\small,breaklines=true,breakanywhere=true,frame=single,framesep=2mm]
Debt-to-assets = Total debt / Total assets
\end{Verbatim}

Example: \texttt{450 / 900 = 0.50}, or 50\%.


State what \emph{total debt} includes. Some sources use interest-bearing debt, while others use total liabilities.


\needspace{5\baselineskip}
\subsection[Debt-to-equity]{Debt-to-equity}
\begin{Verbatim}[fontsize=\small,breaklines=true,breakanywhere=true,frame=single,framesep=2mm]
Debt-to-equity = Defined debt measure / Defined equity measure
\end{Verbatim}

Two conventions are common:


\begin{Verbatim}[fontsize=\small,breaklines=true,breakanywhere=true,frame=single,framesep=2mm]
Interest-bearing debt-to-equity = Interest-bearing debt / Shareholders' equity

Liabilities-to-equity = Total liabilities / Shareholders' equity
\end{Verbatim}

Example using interest-bearing debt: \texttt{450 / 450 = 1.00}.


Do not mix the two conventions. Name the numerator and compare only with a benchmark using the same definition.


\needspace{5\baselineskip}
\subsection[Interest coverage]{Interest coverage}
\begin{Verbatim}[fontsize=\small,breaklines=true,breakanywhere=true,frame=single,framesep=2mm]
Interest coverage = Earnings before interest and tax / Interest expense
\end{Verbatim}

Example: \texttt{120 / 30 = 4.0 times}.


Definitions differ. State whether EBIT, EBITDA, or another measure is used.


\needspace{5\baselineskip}
\section[Profitability]{Profitability}
\needspace{5\baselineskip}
\subsection[Gross profit margin]{Gross profit margin}
\begin{Verbatim}[fontsize=\small,breaklines=true,breakanywhere=true,frame=single,framesep=2mm]
Gross margin = (Net sales - Cost of goods sold) / Net sales
\end{Verbatim}

Example: \texttt{(1,000 - 600) / 1,000 = 40\%}.


\needspace{5\baselineskip}
\subsection[Net profit margin]{Net profit margin}
\begin{Verbatim}[fontsize=\small,breaklines=true,breakanywhere=true,frame=single,framesep=2mm]
Net margin = Net income / Net sales
\end{Verbatim}

Example: \texttt{80 / 1,000 = 8\%}.


\needspace{5\baselineskip}
\subsection[Return on assets]{Return on assets}
\begin{Verbatim}[fontsize=\small,breaklines=true,breakanywhere=true,frame=single,framesep=2mm]
ROA = Net income / Average total assets
\end{Verbatim}

Some analysts use operating income instead of net income. State the numerator. If only one balance-sheet date is available, state that total assets are used as an approximation.


Example: \texttt{80 / 900 = 8.9\%}.


\needspace{5\baselineskip}
\subsection[Return on equity]{Return on equity}
\begin{Verbatim}[fontsize=\small,breaklines=true,breakanywhere=true,frame=single,framesep=2mm]
ROE = Net income available to common shareholders / Average common equity
\end{Verbatim}

Some sources use total net income and total shareholders' equity. State the convention. Example using year-end common equity: \texttt{80 / 450 = 17.8\%}.


\needspace{5\baselineskip}
\section[Shareholder ratios]{Shareholder ratios}
These may be less relevant to a small private business but can test financial literacy.


\needspace{5\baselineskip}
\subsection[Total shareholder return]{Total shareholder return}
\begin{Verbatim}[fontsize=\small,breaklines=true,breakanywhere=true,frame=single,framesep=2mm]
TSR = (Ending share price - Beginning share price + Dividends) / Beginning share price
\end{Verbatim}

\needspace{5\baselineskip}
\subsection[Price-earnings ratio]{Price-earnings ratio}
\begin{Verbatim}[fontsize=\small,breaklines=true,breakanywhere=true,frame=single,framesep=2mm]
P/E = Market price per share / Earnings per share
\end{Verbatim}

\needspace{5\baselineskip}
\subsection[Market-to-book ratio]{Market-to-book ratio}
\begin{Verbatim}[fontsize=\small,breaklines=true,breakanywhere=true,frame=single,framesep=2mm]
Market-to-book = Market price per share / Book value per share
\end{Verbatim}

Use \textbf{book value per share}, not earnings per share.


\needspace{5\baselineskip}
\subsection[Dividend yield]{Dividend yield}
\begin{Verbatim}[fontsize=\small,breaklines=true,breakanywhere=true,frame=single,framesep=2mm]
Dividend yield = Annual dividend per share / Market price per share
\end{Verbatim}

\needspace{5\baselineskip}
\section[Balance-sheet and cash relationships]{Balance-sheet and cash relationships}
\begin{Verbatim}[fontsize=\small,breaklines=true,breakanywhere=true,frame=single,framesep=2mm]
Assets = Liabilities + Equity

Ending cash = Beginning cash + Cash inflows - Cash outflows

Working capital = Current assets - Current liabilities
\end{Verbatim}

Positive working capital does not guarantee liquidity if receivables are uncollectible or inventory is obsolete.


\needspace{5\baselineskip}
\section[Program-performance calculations]{Program-performance calculations}
\needspace{5\baselineskip}
\subsection[Conversion rate]{Conversion rate}
\begin{Verbatim}[fontsize=\small,breaklines=true,breakanywhere=true,frame=single,framesep=2mm]
Conversion rate = Projects reaching next stage / Projects entering stage
\end{Verbatim}

\needspace{5\baselineskip}
\subsection[Cost per qualified lead]{Cost per qualified lead}
\begin{Verbatim}[fontsize=\small,breaklines=true,breakanywhere=true,frame=single,framesep=2mm]
Cost per qualified lead = Tactic cost / Qualified leads attributable to tactic
\end{Verbatim}

\needspace{5\baselineskip}
\subsection[Issue resolution rate]{Issue resolution rate}
\begin{Verbatim}[fontsize=\small,breaklines=true,breakanywhere=true,frame=single,framesep=2mm]
Issue resolution rate = Issues resolved / Issues eligible for resolution
\end{Verbatim}

Define "resolved" and exclude cases outside scope only by a documented rule.


\needspace{5\baselineskip}
\subsection[Funding leverage]{Funding leverage}
\begin{Verbatim}[fontsize=\small,breaklines=true,breakanywhere=true,frame=single,framesep=2mm]
Funding leverage = Partner and external funds / Organization funds
\end{Verbatim}

State whether in-kind contributions are included. Leverage is not proof of impact.


\needspace{5\baselineskip}
\section[Avoiding common errors]{Avoiding common errors}
\begin{itemize}
\item mixing percentages and percentage points;
\item using mismatched years or geographies;
\item dividing by zero or a very small base;
\item using total sales instead of credit sales without disclosure;
\item using ending inventory instead of average inventory without disclosure;
\item confusing market-to-book with P/E;
\item interpreting LQ as growth;
\item interpreting correlation as cause;
\item counting announced and completed outcomes together;
\item counting the same impact in multiple categories;
\item omitting inflation, timing, risk, or lifecycle cost.
\end{itemize}


\needspace{5\baselineskip}
\marginsection{external-link}{Learn more}
\begin{itemize}
\item \href{https://www150.statcan.gc.ca/n1/pub/21-006-x/2008007/def-eng.htm}{Statistics Canada --- Location quotient}: interpretation of employment concentration.
\item \href{https://www.dot.mn.gov/ofrw/freight/PDF/d1plan/wp2.pdf}{Minnesota Department of Transportation --- Regional Economic Analysis and Shift-Share Analysis}: reference growth, industry mix, and differential or local shift.
\item \href{https://www.bdc.ca/en/articles-tools/money-finance/manage-finances/financial-ratios-4-ways-assess-business}{Business Development Bank of Canada --- Four ways to assess a business using financial ratios}: liquidity and other ratio categories.
\item \href{https://www.bdc.ca/en/articles-tools/money-finance/manage-finances/financial-ratios-what-are-how-use}{Business Development Bank of Canada --- Financial ratios}: profitability, leverage, coverage, collection, and turnover measures.
\item \href{https://www.investor.gov/introduction-investing/investing-basics/glossary/price-earnings-pe-ratio}{Investor.gov --- Price-earnings ratio}: P/E definition.
\item \href{https://www.cfainstitute.org/insights/professional-learning/refresher-readings/2026/market-based-valuation-price-enterprise-value-multiples}{CFA Institute --- Market-Based Valuation}: price-to-book and other market multiples.
\item \href{https://www.ontario.ca/document/performance-measurement-agriculture-agri-food-and-economic-development-organizations/appendix-5-examples-identifying-performance-measures-and-data}{Ontario --- Examples of identifying performance measures and data}: economic-development performance measures.
\end{itemize}


When two recognized formulas differ, name the convention, use it consistently, and compare it only with a benchmark calculated the same way.


\chapter[Mock Exam for Practice]{\chaptericon{clipboard-check}Mock Exam for Practice}
\label{chap:practice}
This chapter contains a complete written practice exam. Its questions draw on the topics, frameworks, terminology, and public references used throughout the guide, organized around the publicly described examination format.



\needspace{5\baselineskip}
\marginsection{route}{How to use this file}
Complete the whole chapter as a written simulation, or choose one part for a shorter study session. Confirm the current timing, marks, and requirements with EDAC before setting up a full mock exam.


For a full practice session:


\begin{itemize}
\item complete the 24 practice multiple-choice and true/false questions;
\item choose one contemporary essay question;
\item answer one question from each of Mock Sets A through H;
\item complete the calculation practice if quantitative review is part of your study plan.
\end{itemize}


Change the place, industry, constraint, and stakeholder mix when repeating a question. This keeps the exercise focused on transferable judgment instead of memorized wording.


\needspace{5\baselineskip}
\marginsection{clipboard-check}{Mock written exam}
The following practice questions use constructed scenarios and figures developed from the study-guide chapters. The A-through-H labels simply organize the mock exam by knowledge area.


% practice-questions:start
\needspace{5\baselineskip}
\subsection[Practice multiple-choice and true/false questions]{Practice multiple-choice and true/false questions}
The 24 original questions below mix definitions, distinctions, short applications, professional judgment, and calculations. Allow about 40 minutes, answer every question, and record your choices before checking the answer key.

\subsubsection*{Questions}
\begin{enumerate}[leftmargin=*,itemsep=1.1em]
\needspace{7\baselineskip}
\item \textbf{Professional practice --- Multiple choice.} A councillor asks an economic development officer to announce an unverified major investment before the company has approved disclosure. What is the best response?
\begin{itemize}[leftmargin=1.7em]
\item \textbf{A.} Announce it immediately because elected officials set municipal priorities.
\item \textbf{B.} Confirm the facts, confidentiality obligations, and approval to disclose before preparing any announcement.
\item \textbf{C.} Give the information to one reporter on background so the story can be tested.
\item \textbf{D.} Refuse to communicate about the project at any time.
\end{itemize}
\needspace{7\baselineskip}
\item \textbf{Business retention and expansion --- Multiple choice.} Which activity is most central to a business retention and expansion program?
\begin{itemize}[leftmargin=1.7em]
\item \textbf{A.} Purchasing untargeted advertising in national media.
\item \textbf{B.} Conducting structured outreach to existing employers and coordinating responses to identified barriers.
\item \textbf{C.} Approving every expansion incentive requested by a local business.
\item \textbf{D.} Replacing the municipality's planning approval process.
\end{itemize}
\needspace{7\baselineskip}
\item \textbf{Performance measurement --- Multiple choice.} A program reports that staff completed 80 business visits. In a logic model, this figure is best classified as:
\begin{itemize}[leftmargin=1.7em]
\item \textbf{A.} an input
\item \textbf{B.} an activity
\item \textbf{C.} an output
\item \textbf{D.} a long-term outcome
\end{itemize}
\needspace{7\baselineskip}
\item \textbf{Performance measurement --- Multiple choice.} Which measure provides the strongest evidence that a workforce initiative improved participant employment outcomes?
\begin{itemize}[leftmargin=1.7em]
\item \textbf{A.} The number of social-media impressions generated.
\item \textbf{B.} The number of workshops delivered.
\item \textbf{C.} The share of participants employed after the program compared with a defined baseline and appropriate comparator.
\item \textbf{D.} A testimonial selected by the program manager.
\end{itemize}
\needspace{7\baselineskip}
\item \textbf{Governance --- Multiple choice.} An arm's-length economic development corporation is most likely to maintain public accountability when it has:
\begin{itemize}[leftmargin=1.7em]
\item \textbf{A.} no written mandate so it can respond flexibly
\item \textbf{B.} a clear mandate, delegated authorities, reporting requirements, and conflict-of-interest controls
\item \textbf{C.} complete exemption from municipal policies and public reporting
\item \textbf{D.} responsibility for approving its own municipal budget without oversight
\end{itemize}
\needspace{7\baselineskip}
\item \textbf{Partnerships --- Multiple choice.} Before launching a regional investment-attraction partnership, the participating organizations should first clarify:
\begin{itemize}[leftmargin=1.7em]
\item \textbf{A.} objectives, roles, decision rights, resources, data handling, and success measures
\item \textbf{B.} which partner will receive public credit for every lead
\item \textbf{C.} that all municipal differences will be removed
\item \textbf{D.} that confidential prospect information will be circulated to every employee
\end{itemize}
\needspace{7\baselineskip}
\item \textbf{Communications --- Multiple choice.} During a plant-closure response, which communication approach is most appropriate?
\begin{itemize}[leftmargin=1.7em]
\item \textbf{A.} Speculate about the cause so the municipality appears informed.
\item \textbf{B.} Use a designated spokesperson, separate confirmed facts from unknowns, show empathy, and state the next update time.
\item \textbf{C.} Avoid all communication until every long-term impact is known.
\item \textbf{D.} Publish personal details about affected workers to demonstrate impact.
\end{itemize}
\needspace{7\baselineskip}
\item \textbf{Public engagement --- Multiple choice.} An open house attracts mostly established property owners. What is the best next step before treating the feedback as representative?
\begin{itemize}[leftmargin=1.7em]
\item \textbf{A.} Close engagement because an open invitation was provided.
\item \textbf{B.} Weight every comment equally and claim it represents the whole community.
\item \textbf{C.} Identify missing groups and barriers, then add accessible and targeted ways for them to participate.
\item \textbf{D.} Discard the open-house feedback entirely.
\end{itemize}
\needspace{7\baselineskip}
\item \textbf{Strategic planning --- Multiple choice.} Which sequence best reflects a defensible strategic-planning process?
\begin{itemize}[leftmargin=1.7em]
\item \textbf{A.} Select projects, write a vision, collect data, then identify stakeholders.
\item \textbf{B.} Establish mandate, assess evidence and context, engage stakeholders, set priorities, build an action plan, implement, and evaluate.
\item \textbf{C.} Copy a neighbouring strategy, announce targets, and assign measures after implementation.
\item \textbf{D.} Start implementation, then seek approval if early results are positive.
\end{itemize}
\needspace{7\baselineskip}
\item \textbf{Planning and land use --- Multiple choice.} A manufacturer asks the economic development officer to guarantee rezoning of a site. The officer should:
\begin{itemize}[leftmargin=1.7em]
\item \textbf{A.} provide the guarantee to keep the project moving
\item \textbf{B.} explain the approval process and risks, coordinate early review, and avoid promising the outcome
\item \textbf{C.} tell the company that planning considerations are irrelevant to investment attraction
\item \textbf{D.} submit the application without the company's authorization
\end{itemize}
\needspace{7\baselineskip}
\item \textbf{Investment attraction --- Multiple choice.} Which information is most useful when first qualifying an investment lead?
\begin{itemize}[leftmargin=1.7em]
\item \textbf{A.} The prospect's operating requirements, decision criteria, timeline, authority, and confidentiality needs.
\item \textbf{B.} Only the number of jobs mentioned in the initial inquiry.
\item \textbf{C.} The community's preferred announcement date.
\item \textbf{D.} A promise that every requested incentive will be approved.
\end{itemize}
\needspace{7\baselineskip}
\item \textbf{Marketing and sales --- Multiple choice.} Which statement is the strongest investment value proposition?
\begin{itemize}[leftmargin=1.7em]
\item \textbf{A.} Our community is a great place to do business.
\item \textbf{B.} Our serviced employment area offers verified utility capacity, highway access, and a workforce program aligned with food-processing occupations.
\item \textbf{C.} We will beat every competing community on cost.
\item \textbf{D.} Our brochure contains information about all local sectors.
\end{itemize}
\needspace{7\baselineskip}
\item \textbf{Finance --- Multiple choice.} Which financial statement reports a business's assets, liabilities, and equity at a specific date?
\begin{itemize}[leftmargin=1.7em]
\item \textbf{A.} Income statement
\item \textbf{B.} Balance sheet
\item \textbf{C.} Cash-flow statement
\item \textbf{D.} Marketing plan
\end{itemize}
\needspace{7\baselineskip}
\item \textbf{Finance --- Multiple choice.} A business has current assets of \$240,000 and current liabilities of \$160,000. What is its current ratio?
\begin{itemize}[leftmargin=1.7em]
\item \textbf{A.} 0.67
\item \textbf{B.} 1.00
\item \textbf{C.} 1.50
\item \textbf{D.} 80,000
\end{itemize}
\needspace{7\baselineskip}
\item \textbf{Economic analysis --- Multiple choice.} A local industry has a location quotient of 1.40. Which interpretation is most defensible?
\begin{itemize}[leftmargin=1.7em]
\item \textbf{A.} The industry's local employment share is 40\% larger than the benchmark share.
\item \textbf{B.} The industry grew by 40\% during the period.
\item \textbf{C.} The industry is guaranteed to export 40\% of its output.
\item \textbf{D.} The industry is 40\% more productive than the national industry.
\end{itemize}
\needspace{7\baselineskip}
\item \textbf{Economic impact --- Multiple choice.} Which proposed benefit should be excluded from an estimate of net new local economic impact?
\begin{itemize}[leftmargin=1.7em]
\item \textbf{A.} Spending that would have occurred locally even without the project.
\item \textbf{B.} Verified purchases from local suppliers that occur because of the project.
\item \textbf{C.} New wages paid to local workers because of the project.
\item \textbf{D.} Incremental municipal revenue attributable to the project.
\end{itemize}
\needspace{7\baselineskip}
\item \textbf{Economic analysis --- Multiple choice.} In shift-share analysis, a positive local competitive effect most directly indicates that:
\begin{itemize}[leftmargin=1.7em]
\item \textbf{A.} the local industry outperformed the growth expected from national and industry-mix effects
\item \textbf{B.} all local job growth was caused by the economic development program
\item \textbf{C.} the industry has the largest employment base in the community
\item \textbf{D.} the industry will continue growing at the same rate
\end{itemize}
\needspace{7\baselineskip}
\item \textbf{Inclusive development --- Multiple choice.} A major project may affect the rights and interests of an Indigenous community. Which is the most appropriate economic development approach?
\begin{itemize}[leftmargin=1.7em]
\item \textbf{A.} Treat the community only as another stakeholder in a general open house.
\item \textbf{B.} Engage early and respectfully, understand rights and protocols, involve the proper authorities, and avoid treating consultation as a promotional exercise.
\item \textbf{C.} Wait until every project decision has been finalized before making contact.
\item \textbf{D.} Assume a benefits agreement eliminates all Crown consultation obligations.
\end{itemize}
\needspace{7\baselineskip}
\item \textbf{Performance measurement --- True or false.} True or false: Program outputs and program outcomes are interchangeable terms.
\begin{itemize}[leftmargin=1.7em]
\item \textbf{T.} True
\item \textbf{F.} False
\end{itemize}
\needspace{7\baselineskip}
\item \textbf{Finance --- True or false.} True or false: A positive net present value is sufficient by itself to approve a public economic development project.
\begin{itemize}[leftmargin=1.7em]
\item \textbf{T.} True
\item \textbf{F.} False
\end{itemize}
\needspace{7\baselineskip}
\item \textbf{Labour market --- True or false.} True or false: The unemployment rate counts every working-age person who does not have a job as unemployed.
\begin{itemize}[leftmargin=1.7em]
\item \textbf{T.} True
\item \textbf{F.} False
\end{itemize}
\needspace{7\baselineskip}
\item \textbf{Business retention and expansion --- True or false.} True or false: Aggregated BR\&E intelligence can inform workforce, infrastructure, and policy priorities as well as individual business follow-up.
\begin{itemize}[leftmargin=1.7em]
\item \textbf{T.} True
\item \textbf{F.} False
\end{itemize}
\needspace{7\baselineskip}
\item \textbf{Marketing and sales --- True or false.} True or false: Economic development marketing and investment sales are interchangeable because both promote the community.
\begin{itemize}[leftmargin=1.7em]
\item \textbf{T.} True
\item \textbf{F.} False
\end{itemize}
\needspace{7\baselineskip}
\item \textbf{Strategic planning --- True or false.} True or false: An implementable economic development strategy should assign responsibilities, timelines, resources, and performance measures.
\begin{itemize}[leftmargin=1.7em]
\item \textbf{T.} True
\item \textbf{F.} False
\end{itemize}
\end{enumerate}

\subsubsection*{Answer key and explanations}
\begin{enumerate}[leftmargin=*,itemsep=.8em]
\needspace{4\baselineskip}
\item \textbf{B --- Confirm the facts, confidentiality obligations, and approval to disclose before preparing any announcement.} Economic developers should protect confidential information, verify facts, and coordinate authorized communication before a public announcement.
\needspace{4\baselineskip}
\item \textbf{B --- Conducting structured outreach to existing employers and coordinating responses to identified barriers.} A sound BR\&E program builds relationships with existing firms, identifies opportunities and risks, and coordinates appropriate responses without promising outcomes outside the EDO's authority.
\needspace{4\baselineskip}
\item \textbf{C --- an output} Outputs count what a program directly delivers. Outcomes describe changes that follow, and should not be assumed from activity counts alone.
\needspace{4\baselineskip}
\item \textbf{C --- The share of participants employed after the program compared with a defined baseline and appropriate comparator.} Outcome evidence is stronger when it specifies the population, timing, baseline, and a reasonable comparison rather than relying only on outputs or anecdotes.
\needspace{4\baselineskip}
\item \textbf{B --- a clear mandate, delegated authorities, reporting requirements, and conflict-of-interest controls} An arm's-length structure still needs a documented mandate, appropriate delegated authority, transparent reporting, and controls for ethical and financial stewardship.
\needspace{4\baselineskip}
\item \textbf{A --- objectives, roles, decision rights, resources, data handling, and success measures} Effective partnerships begin with shared objectives and an operating framework that covers authority, resources, information, accountability, and evaluation.
\needspace{4\baselineskip}
\item \textbf{B --- Use a designated spokesperson, separate confirmed facts from unknowns, show empathy, and state the next update time.} Crisis communication should be accurate, empathetic, coordinated, and transparent about uncertainty and the timing of further information.
\needspace{4\baselineskip}
\item \textbf{C --- Identify missing groups and barriers, then add accessible and targeted ways for them to participate.} Representative engagement requires attention to who is missing, why they may be missing, and which additional methods can reduce participation barriers.
\needspace{4\baselineskip}
\item \textbf{B --- Establish mandate, assess evidence and context, engage stakeholders, set priorities, build an action plan, implement, and evaluate.} A credible strategy connects mandate, evidence, engagement, choices, resources, delivery, measurement, and adaptation.
\needspace{4\baselineskip}
\item \textbf{B --- explain the approval process and risks, coordinate early review, and avoid promising the outcome} Economic development staff can convene reviewers and clarify process, but they must distinguish facilitation from statutory decision-making.
\needspace{4\baselineskip}
\item \textbf{A --- The prospect's operating requirements, decision criteria, timeline, authority, and confidentiality needs.} Lead qualification establishes the decision need, project requirements, timing, authority, fit, and information rules before major resources are committed.
\needspace{4\baselineskip}
\item \textbf{B --- Our serviced employment area offers verified utility capacity, highway access, and a workforce program aligned with food-processing occupations.} A useful value proposition is specific to the target audience, evidence based, differentiated, and connected to the prospect's decision criteria.
\needspace{4\baselineskip}
\item \textbf{B --- Balance sheet} The balance-sheet equation is assets = liabilities + equity, measured at a stated point in time.
\needspace{4\baselineskip}
\item \textbf{C --- 1.50} Current ratio = current assets / current liabilities = \$240,000 / \$160,000 = 1.50. Interpretation still depends on industry, trends, and asset quality.
\needspace{4\baselineskip}
\item \textbf{A --- The industry's local employment share is 40\% larger than the benchmark share.} An LQ of 1.40 means the industry's share of local employment is 1.4 times its share in the benchmark economy. It does not by itself establish growth, productivity, or causation.
\needspace{4\baselineskip}
\item \textbf{A --- Spending that would have occurred locally even without the project.} Net impact should exclude deadweight and also test leakage, displacement, substitution, timing, and the geographic boundary of the analysis.
\needspace{4\baselineskip}
\item \textbf{A --- the local industry outperformed the growth expected from national and industry-mix effects} A positive local competitive effect suggests stronger local performance than national growth and industry mix would predict, but further investigation is needed to explain why.
\needspace{4\baselineskip}
\item \textbf{B --- Engage early and respectfully, understand rights and protocols, involve the proper authorities, and avoid treating consultation as a promotional exercise.} Economic developers should support early, respectful, rights-aware relationships and ensure that the proper legal and governmental authorities lead duties that fall outside the EDO's mandate.
\needspace{4\baselineskip}
\item \textbf{F --- False} Keeping outputs and outcomes separate prevents activity counts from being mistaken for evidence of impact.
\needspace{4\baselineskip}
\item \textbf{F --- False} Public decisions require a broader business case than a single financial metric, even when the calculation is sound.
\needspace{4\baselineskip}
\item \textbf{F --- False} The labour force includes employed people and people classified as unemployed; it excludes people not participating in the labour force.
\needspace{4\baselineskip}
\item \textbf{T --- True} A mature BR\&E program separates confidential case management from responsibly aggregated evidence that can guide community-wide action.
\needspace{4\baselineskip}
\item \textbf{F --- False} The functions reinforce one another but use different objectives, activities, measures, and stages of the prospect journey.
\needspace{4\baselineskip}
\item \textbf{T --- True} Implementation detail connects strategic intent to delivery, accountability, learning, and course correction.
\end{enumerate}
% practice-questions:end


\needspace{5\baselineskip}
\subsection[Short-answer practice sets]{Short-answer practice sets}
For each set:


\begin{itemize}
\item choose one of the three questions;
\item set a 20-minute timer;
\item spend about 2 minutes planning, 15 minutes writing, and 3 minutes checking; and
\item score the response using the rubric below.
\end{itemize}


\needspace{5\baselineskip}
\subsubsection[Mock set A: Professional judgment and evidence]{Mock set A: Professional judgment and evidence}
\begin{enumerate}
\item A fictional regional development office is evaluating a supplier-development pilot that has reached only half its participation target. Senior leaders want to renew it based on several positive testimonials. Explain how the officer should assess the evidence, distinguish outputs from outcomes, identify data limitations, and recommend whether to continue, modify, or end the program.
\item A public dashboard shows that employment in a target sector grew sharply. Statistics Canada later revises the underlying series and the apparent increase becomes much smaller. Prepare a response that addresses correction, communication with decision-makers, documentation, and future quality control.
\item A newly appointed economic-development board asks staff to promise a project approval that legally belongs to another public authority. Explain how an economic developer should clarify roles, preserve momentum, and offer useful support without exceeding the organization's mandate.
\end{enumerate}


Review the response for facts, authority, assumptions, limitations, public interest, evidence quality, documentation, corrective action, and relationship management.


\needspace{5\baselineskip}
\subsubsection[Mock set B: Business intelligence and response]{Mock set B: Business intelligence and response}
\begin{enumerate}
\item A fictional service centre has approximately 140 employers spread across a large rural area and only one staff member available for outreach. Design a twelve-month business-listening program that balances representation, confidentiality, follow-up capacity, and measurable value.
\item Several firms report unrelated symptoms - missed shipments, high turnover, and delayed permits - but staff suspect a shared infrastructure constraint. Describe how you would test the hypothesis, triage firm-specific needs, and convert the findings into a coordinated response.
\item A business survey receives strong participation from large firms but very little response from small, home-based, or newcomer-owned businesses. Explain what can and cannot be concluded, how the evidence gap should be addressed, and how results should be reported.
\end{enumerate}


Review the response for purpose, sampling, accessibility, consent, confidentiality, issue triage, partner referrals, systemic analysis, feedback, implementation, and measures.


\needspace{5\baselineskip}
\subsubsection[Mock set C: Investment readiness and market change]{Mock set C: Investment readiness and market change}
\begin{enumerate}
\item A fictional agricultural region wants to attract food-processing investment, but its cold-storage capacity, wastewater limits, and seasonal workforce are not well documented. Develop a target-market and readiness process before promotional spending begins.
\item An international prospect requests verified information on power capacity, permitting, labour, incentives, and available sites within seventy-two hours. Some information is incomplete and the project's decision process is unclear. Describe how you would qualify the opportunity and organize a credible response.
\item A cluster of small exporters faces simultaneous exchange-rate volatility, changing border requirements, and rising freight costs. Build a scenario-based response that distinguishes immediate firm assistance from longer-term market and supply-chain development.
\end{enumerate}


Review the response for strategic fit, verified data, lead qualification, response ownership, confidentiality, intergovernmental coordination, limitations, timelines, risk, and aftercare.


\needspace{5\baselineskip}
\subsubsection[Mock set D: Governance and measurement]{Mock set D: Governance and measurement}
\begin{enumerate}
\item Three fictional municipalities jointly fund an economic-development organization but disagree about voting rights, service levels, and how regional results should be credited. Propose a governance and accountability framework.
\item A two-year business accelerator reports workshops delivered, attendance, and social-media reach, but the board wants evidence of business value. Develop a logic model and a practical evaluation plan without claiming results the program cannot attribute.
\item A sudden budget reduction requires an economic-development team to pause or redesign several initiatives. Develop a transparent prioritization process that considers mandate, evidence, equity, contractual obligations, opportunity cost, and community impact.
\end{enumerate}


Review the response for mandate, representation, decision rights, service standards, funding, conflicts, data ownership, outputs, outcomes, attribution, review cycles, and public reporting.


\needspace{5\baselineskip}
\subsubsection[Mock set E: Partnerships and program delivery]{Mock set E: Partnerships and program delivery}
\begin{enumerate}
\item A municipality, an Indigenous development corporation, a college, and a utility are exploring a shared employment and training hub. Describe how the parties should establish relationships, define authority, protect information, allocate risk, and decide whether to proceed.
\item A proposed mentoring initiative would match volunteer executives with early-stage businesses. Design the operating model, including participant selection, boundaries, conflicts, accessibility, safeguarding, feedback, and exit provisions.
\item A funding partner asks staff to announce an ambitious job target before the delivery partners have validated it. Explain how you would respond, revise the target-setting process, and preserve the partnership.
\end{enumerate}


Review the response for strategic fit, rights and jurisdiction, governance, written agreements, roles, resources, risk allocation, privacy, conflicts, dispute resolution, exit, and shared outcomes.


\needspace{5\baselineskip}
\subsubsection[Mock set F: Communication and engagement]{Mock set F: Communication and engagement}
\begin{enumerate}
\item A multi-year street reconstruction project will disrupt access to a fictional commercial district during its busiest season. Develop a communication and business-support plan before, during, and after construction.
\item A proposed high-energy-use facility may create investment and tax benefits but raises questions about power, water, land, and community priorities. Design an engagement process that clearly separates information, consultation, technical review, and formal decision authority.
\item A reporter asks about an employer restructuring before the firm has made a public announcement. Prepare the organization's response and a subsequent media plan for the period after verified information becomes public.
\end{enumerate}


Review the response for purpose, audiences, accessible channels, verified messages, spokesperson roles, privacy, stakeholder and rights-holder participation, influence on decisions, feedback, timing, and evaluation.


\needspace{5\baselineskip}
\subsubsection[Mock set G: Planning and location readiness]{Mock set G: Planning and location readiness}
\begin{enumerate}
\item A fictional community markets an employment site as development-ready, but ecological work, transportation access, and utility capacity remain uncertain. Explain how the claim should be corrected and how a defensible readiness program should proceed.
\item An economic strategy was written before major changes in population, commuting, technology, and climate risk. Design an efficient strategy refresh that uses new evidence without discarding useful prior work.
\item A site-selection request compares three communities. Your community performs well on several factors but has a material weakness in one operating requirement. Explain how you would present the evidence, test possible mitigation, and advise local decision-makers.
\end{enumerate}


Review the response for evidence, planning authority, rights and constraints, verified readiness, alternatives, transparent marketing, interdepartmental roles, implementation, risk, and monitoring.


\needspace{5\baselineskip}
\subsubsection[Mock set H: Finance and impact]{Mock set H: Finance and impact}
\begin{enumerate}
\item A growing fictional manufacturer reports higher sales and positive net income but repeatedly uses short-term borrowing to meet payroll. Explain which financial information and ratios you would review, the limits of the EDO's role, and when a qualified adviser is needed.
\item A proposed industrial-park expansion requires municipal servicing, land acquisition, and a long construction period. Develop an impact-analysis framework that separates financial feasibility, fiscal effects, economic effects, environmental and social considerations, distribution, uncertainty, and lifecycle risk.
\item A publicly funded entrepreneurship program requests a substantial renewal after its pilot year. Design the budget review and recommendation process, including relevant costs, performance evidence, alternatives, risk, approval authority, and reporting.
\end{enumerate}


Review the response for financial statements, cash and liquidity, relevant ratios, assumptions, qualified referrals, incremental effects, scenarios, distribution, lifecycle costs, uncertainty, decision criteria, and monitoring.


\needspace{5\baselineskip}
\subsection[Contemporary essay questions]{Contemporary essay questions}
Set 70 minutes and observe the current public word limit. These purpose-built questions combine themes and analytical approaches from the contemporary-issues chapter with constructed local circumstances.


\begin{enumerate}
\item A prolonged disruption affects a major Canadian transportation corridor while firms are also changing suppliers. Analyze how the shock could reach a community familiar to you and recommend immediate, medium-term, and resilience actions.
\item Electricity demand from industrial expansion, housing growth, and digital infrastructure is increasing faster than local capacity. Analyze the economic-development trade-offs and propose a fair, evidence-based response.
\item A regional post-secondary institution substantially reduces enrolment and programming. Analyze effects on labour supply, housing, local businesses, innovation, and community services, then recommend a coordinated adjustment strategy.
\item Small and medium-sized firms are adopting artificial intelligence unevenly. Analyze potential productivity gains, workforce transitions, infrastructure needs, governance risks, and an appropriate local economic-development role.
\item Repeated extreme-weather disruptions are changing insurance availability, infrastructure costs, tourism patterns, and business continuity. Develop an economic resilience strategy for a community familiar to you.
\item Population growth slows while the working-age population ages and employers continue to report skill shortages. Analyze the apparent contradiction and recommend actions that distinguish attraction, retention, participation, productivity, and service capacity.
\item An Indigenous-led economic organization proposes an equity partnership in a major regional infrastructure project. Analyze how relationships, rights, governance, finance, procurement, capacity, benefits, risk, and long-term accountability should be approached.
\item Commercial vacancies rise even though regional employment remains stable. Analyze possible causes, identify the evidence needed, and recommend a place-based response that avoids treating symptoms as causes.
\end{enumerate}


To create additional essay practice, combine one external change, two local transmission channels, one stakeholder tension, one resource constraint, and three time horizons. Do not reuse a memorized thesis.


\needspace{5\baselineskip}
\subsection[Optional calculation practice]{Optional calculation practice}
All figures and situations below were prepared for calculation practice.


\begin{enumerate}
\item A local industry has 720 of 18,000 jobs. The same industry has 24,000 of 800,000 jobs in the reference area. Calculate and interpret the LQ.
\item A local industry grows from 1,000 to 1,120 jobs. Total reference employment grows 3\%, and reference industry employment grows 8\%. Calculate the three shift-share effects.
\item A firm has current assets of 540, inventory of 180, and current liabilities of 270. Calculate current and quick ratios.
\item A firm has cost of goods sold of 960, beginning inventory of 180, and ending inventory of 300. Calculate average inventory and turnover.
\item Average accounts receivable are 180 and annual net credit sales are 1,460. Calculate DSO using 365 days.
\item Interest-bearing debt is 750, assets are 1,250, and equity is 500. Calculate debt-to-assets and interest-bearing debt-to-equity.
\item Sales are 2,000, cost of goods sold is 1,200, and net income is 120. Calculate gross and net margins.
\item A tactic costs 36,000 and produces 18 qualified leads, 6 proposals, and 2 completed investments. Calculate cost per qualified lead and stage conversions.
\end{enumerate}


Answers:


\begin{enumerate}
\item About \texttt{1.33}; the local employment share is about one-third higher than the reference share.
\item Local growth is \texttt{12\%}; reference effect \texttt{30}, industry mix \texttt{50}, local shift \texttt{40}, and total change \texttt{120}.
\item Current ratio \texttt{2.00}; quick ratio about \texttt{1.33}.
\item Average inventory \texttt{240}; turnover \texttt{4.0}.
\item \texttt{45 days}.
\item Debt-to-assets \texttt{60\%}; interest-bearing debt-to-equity \texttt{1.50}.
\item Gross margin \texttt{40\%}; net margin \texttt{6\%}.
\item \texttt{2,000} per qualified lead; lead-to-proposal about \texttt{33.3\%}; proposal-to-completion about \texttt{33.3\%}; lead-to-completion about \texttt{11.1\%}.
\end{enumerate}


\needspace{5\baselineskip}
\subsection[Self-scoring rubric for an 8-mark practice answer]{Self-scoring rubric for an 8-mark practice answer}
Score the practice exam with this independently created rubric:


\begin{center}
\small
\setlength{\LTpre}{0.4em}
\setlength{\LTpost}{0.8em}
\begin{longtable}{>{\RaggedRight\arraybackslash}p{0.290\linewidth}>{\RaggedRight\arraybackslash}p{0.610\linewidth}}
\rowcolor{DeepNavy}
\color{white}\bfseries Criterion & \color{white}\bfseries Marks \\
\endfirsthead
\rowcolor{DeepNavy}
\color{white}\bfseries Criterion & \color{white}\bfseries Marks \\
\endhead
Directly answers every part and establishes context & 1 \\
Uses a correct, logical process & 2 \\
Identifies stakeholders, roles, and authority & 1 \\
Uses evidence and recognizes assumptions & 1 \\
Provides implementation, resources, and timing & 1 \\
Addresses ethics, risk, inclusion, and contingency & 1 \\
Defines outputs, outcomes, and follow-up & 1 \\
\textbf{Total} & \textbf{8} \\
\end{longtable}
\end{center}
After scoring, rewrite only the weakest section so the review stays focused.


\needspace{5\baselineskip}
\marginsection{external-link}{Learn more}
\begin{itemize}
\item \href{https://exam.edac.ca/exam-prep/}{EDAC --- Exam Preparation}: confirm the current public exam format before setting timers or scoring practice.
\item \href{https://edac.ca/about/code-of-ethics/}{EDAC --- Code of Ethics}: use the current code when reviewing ethics and professional-judgment answers.
\item \hyperref[chap:glossary]{Glossary with definition sources}: check terminology and follow the public references after each practice session.
\item \hyperref[chap:calculations]{Calculations and formulas}: review formula conventions before scoring quantitative drills.
\item \hyperref[chap:sources]{Public sources and disclaimer}: verify current exam and contemporary-issue sources.
\end{itemize}


Use the mock exam to practise process, judgment, evidence, and communication. Change the scenarios between attempts and avoid memorizing model wording.


\chapter[Contemporary Issues]{\chaptericon{globe-2}Contemporary Issues}
\label{chap:contemporary}
The contemporary essay should show how a major event could affect the local economy and how an EDO could respond. The analytical questions apply the public issues and study frameworks covered in this guide.



\needspace{5\baselineskip}
\section[Analysis framework]{Analysis framework}
Use LOCAL IMPACT:


\begin{itemize}
\item \textbf{L - Local baseline:} sectors, firms, workers, markets, and existing conditions.
\item \textbf{O - Originating event:} what changed, when, and through which policy or market mechanism.
\item \textbf{C - Channels:} demand, costs, trade, finance, labour, infrastructure, confidence, or regulation.
\item \textbf{A - Affected groups:} firms, workers, residents, governments, rights-holders, and vulnerable groups.
\item \textbf{L - Local evidence:} current quantitative and qualitative evidence.
\item \textbf{I - Immediate response:} urgent information, stabilization, referrals, and coordination.
\item \textbf{M - Medium-term action:} productivity, diversification, workforce, finance, infrastructure, or market development.
\item \textbf{P - Partners and powers:} who has authority, resources, knowledge, and trust.
\item \textbf{A - Accountability:} objectives, measures, reporting, ethics, and equity.
\item \textbf{C - Contingency:} scenarios, triggers, and adjustment.
\item \textbf{T - Timeline:} immediate, 90-day, one-year, and longer-term actions.
\end{itemize}


\needspace{5\baselineskip}
\section[Current Canadian themes as of July 25, 2026]{Current Canadian themes as of July 25, 2026}
These public themes are starting points for independent analysis and mock-exam preparation.


\needspace{5\baselineskip}
\subsection[Trade restructuring and tariffs]{Trade restructuring and tariffs}
Questions to prepare:


\begin{itemize}
\item Which local firms export to or import from the United States?
\item Which inputs and products are tariff-exposed?
\item Can firms diversify markets, suppliers, products, or logistics?
\item What short-term financing, advisory, or worker support exists?
\item How could import replacement create an opportunity without creating an inefficient local strategy?
\end{itemize}


The Government of Canada's \href{https://www.canada.ca/en/department-finance/programs/international-trade-finance-policy/canadas-tariff-responses.html}{tariff-response page} summarizes current sector measures and support. \href{https://ised-isde.canada.ca/site/trade-data-online/en}{Trade Data Online} can identify product and industry trade exposure.


\needspace{5\baselineskip}
\subsection[Growth, inflation, interest rates, and uncertainty]{Growth, inflation, interest rates, and uncertainty}
The \href{https://www.bankofcanada.ca/publications/mpr/mpr-2026-07-15/}{Bank of Canada's July 2026 Monetary Policy Report} described weak recent growth with signs of improvement, elevated near-term inflation, and continued uncertainty related to trade and geopolitical conditions.


Local questions:


\begin{itemize}
\item How do borrowing costs affect construction, business investment, and municipal capital?
\item Which households and firms are sensitive to price changes?
\item Are local conditions stronger or weaker than the national picture?
\item What leading indicators should the EDO monitor?
\end{itemize}


\needspace{5\baselineskip}
\subsection[Housing, workforce, and development feasibility]{Housing, workforce, and development feasibility}
CMHC's \href{https://www.cmhc-schl.gc.ca/media-newsroom/news-releases/2026/economic-uncertainty-continue-weighing-housing-market-cmhc}{July 2026 Housing Market Outlook update} emphasizes economic uncertainty, slower population growth, borrowing costs, and regional differences.


Local questions:


\begin{itemize}
\item Is housing constraining recruitment and retention?
\item Which price points and housing forms are missing?
\item What limits supply: land, servicing, approvals, finance, construction capacity, or demand?
\item How do housing actions connect to transit, infrastructure, services, and employment areas?
\end{itemize}


\needspace{5\baselineskip}
\subsection[Labour, demographics, and skills]{Labour, demographics, and skills}
Use Statistics Canada's \href{https://www.statcan.gc.ca/en/subjects-start/labour\_}{Labour Statistics portal} and local/provincial sources.


Local questions:


\begin{itemize}
\item Which occupations are difficult to recruit or retain?
\item Is the issue skills, wages, housing, transportation, demographics, work conditions, or awareness?
\item What is the commuting shed?
\item Which groups face barriers to participation?
\item What can employers, educators, governments, and community organizations each do?
\end{itemize}


\needspace{5\baselineskip}
\subsection[Artificial intelligence and digital adoption]{Artificial intelligence and digital adoption}
Consider:


\begin{itemize}
\item productivity and new products;
\item task change rather than only job counts;
\item skills and management capacity;
\item data quality, privacy, cybersecurity, bias, and procurement;
\item broadband, computing, and power infrastructure;
\item uneven adoption between large and small firms;
\item local innovation and supplier opportunities.
\end{itemize}


An EDO response could combine awareness, assessment, training, demonstration projects, trusted advisers, financing, and outcome measurement.


\needspace{5\baselineskip}
\subsection[Climate adaptation and economic resilience]{Climate adaptation and economic resilience}
Canada's \href{https://www.canada.ca/en/services/environment/weather/climatechange/climate-plan/national-adaptation-strategy/full-strategy.html}{National Adaptation Strategy} links climate readiness to communities, workers, infrastructure, supply chains, and competitiveness.


Local questions:


\begin{itemize}
\item Which sectors and sites face flood, wildfire, heat, storm, drought, coastal, or permafrost risk?
\item What infrastructure and supply-chain dependencies are critical?
\item How will risk affect insurance, finance, operating continuity, and investment?
\item What adaptation investments produce economic and public-safety benefits?
\item How will the community protect groups with less capacity to adapt?
\end{itemize}


\needspace{5\baselineskip}
\subsection[Indigenous economic development and reconciliation]{Indigenous economic development and reconciliation}
The Government of Canada's \href{https://www.canada.ca/en/services/indigenous-peoples/business-and-economic-development-indigenous-peoples.html}{Business and economic development for Indigenous Peoples} page provides a current program entry point.


Prepare to discuss:


\begin{itemize}
\item rights and jurisdiction;
\item relationship and protocol;
\item Indigenous-led priorities;
\item governance and decision-making;
\item land, infrastructure, skills, procurement, equity participation, and finance;
\item benefit sharing and long-term capacity;
\item data ownership, confidentiality, and cultural knowledge;
\item measurement defined with the Indigenous partner.
\end{itemize}


Avoid treating consultation as a substitute for relationship or consent where the legal context requires more.


\needspace{5\baselineskip}
\section[Seven-day refresh]{Seven-day refresh}
One week before the exam:


\begin{enumerate}
\item review EDAC news and policy updates;
\item read the latest Bank of Canada Monetary Policy Report summary;
\item review the latest Labour Force Survey and local labour data;
\item check Trade Data Online for exposed sectors;
\item read current federal and provincial/territorial budgets or economic updates;
\item review housing, construction, and population indicators;
\item scan current events affecting the candidate's own community;
\item select two events and prepare a one-page LOCAL IMPACT outline for each.
\end{enumerate}


\needspace{5\baselineskip}
\section[Source discipline]{Source discipline}
For every contemporary claim:


\begin{itemize}
\item use a current, authoritative source;
\item note publication and reference dates;
\item distinguish Canada-wide and local conditions;
\item distinguish observed results, forecasts, and scenarios;
\item explain uncertainty;
\item avoid political advocacy unrelated to the professional analysis;
\item show who has authority to act.
\end{itemize}


\needspace{5\baselineskip}
\section[Essay outline template]{Essay outline template}
\begin{Verbatim}[fontsize=\small,breaklines=true,breakanywhere=true,frame=single,framesep=2mm]
Title:
Event and date:
Community/region:
Thesis:

1. Baseline
2. Impact channels
3. Sector and population effects
4. Evidence and uncertainty
5. Immediate actions
6. Medium- and long-term strategy
7. Partners, authority, and resources
8. Communication, ethics, and equity
9. Measures, triggers, and contingency
10. Conclusion
\end{Verbatim}

\needspace{5\baselineskip}
\marginsection{external-link}{Learn more}
\begin{itemize}
\item \href{https://www.bankofcanada.ca/publications/mpr/}{Bank of Canada --- Monetary Policy Report}: growth, inflation, interest rates, trade, productivity, and risk.
\item \href{https://www.statcan.gc.ca/en/subjects-start/labour\_}{Statistics Canada --- Labour statistics}: employment, wages, vacancies, hours, and productivity.
\item \href{https://ised-isde.canada.ca/site/trade-data-online/en}{ISED --- Trade Data Online}: imports and exports by industry, product, partner, province, and territory.
\item \href{https://www.cmhc-schl.gc.ca/professionals/housing-markets-data-and-research}{Canada Mortgage and Housing Corporation --- Housing market information}: housing supply, construction, ownership, rental, and market evidence.
\item \href{https://www.canada.ca/en/services/environment/weather/climatechange/climate-plan/national-adaptation-strategy/full-strategy.html}{Government of Canada --- National Adaptation Strategy}: climate resilience and adaptation.
\item \href{https://www.canada.ca/en/services/indigenous-peoples/business-and-economic-development-indigenous-peoples.html}{Government of Canada --- Business and economic development for Indigenous Peoples}: federal program and information entry points.
\end{itemize}


Contemporary evidence ages quickly. Recheck publication dates, definitions, revisions, and geographic relevance close to the examination rather than relying on this chapter's dated snapshot.


\chapter[Glossary]{\chaptericon{book-a}Glossary}
\label{chap:glossary}
Definitions are plain-language syntheses written for exam review. They are not quotations. Confirm jurisdiction-specific legal terms and follow the linked primary sources for technical detail. Each row pairs the study definition with one or more public sources for deeper learning and comparison.



\needspace{5\baselineskip}
\section[A]{A}
\begin{center}
\small
\setlength{\LTpre}{0.4em}
\setlength{\LTpost}{0.8em}
\begin{longtable}{>{\RaggedRight\arraybackslash}p{0.180\linewidth}>{\RaggedRight\arraybackslash}p{0.400\linewidth}>{\RaggedRight\arraybackslash}p{0.300\linewidth}}
\rowcolor{DeepNavy}
\color{white}\bfseries Term & \color{white}\bfseries Definition & \color{white}\bfseries Sources \\
\endfirsthead
\rowcolor{DeepNavy}
\color{white}\bfseries Term & \color{white}\bfseries Definition & \color{white}\bfseries Sources \\
\endhead
\textbf{Activity} & Work performed by a program, such as business visits or workshops. & \href{https://www.oecd.org/en/publications/glossary-of-key-terms-in-evaluation-and-results-based-management-for-sustainable-development-second-edition\_632da462-en-fr-es.html}{OECD --- Glossary of Key Terms in Evaluation and Results-Based Management} \\
\textbf{Additionality} & The part of an outcome that likely would not have occurred, or would have occurred later or at smaller scale, without the intervention. & \href{https://www.gov.uk/government/publications/the-green-book-appraisal-and-evaluation-in-central-government/the-green-book-2026}{UK Government --- The Green Book: appraisal and evaluation in central government} \\
\textbf{Aftercare} & Structured relationship and support for firms already established in a community, often to encourage retention and reinvestment. & \href{https://www.worldbank.org/en/topic/investment-climate/brief/investment-policy-and-promotion}{World Bank --- Investment Policy and Promotion} \\
\textbf{Arm's-length organization} & A separate corporation or agency that delivers a public mandate with defined independence and accountability. & \href{https://www.canada.ca/en/treasury-board-secretariat/services/guidance-crown-corporations/guiding-principles-management-crown-corporations.html}{Treasury Board of Canada Secretariat --- Guiding Principles for the Management of Crown Corporations} \\
\textbf{Assets} & Present economic resources controlled by an entity as a result of past events. Under IFRS, an economic resource is a right with the potential to produce economic benefits. & \href{https://www.ifrs.org/issued-standards/list-of-standards/conceptual-framework/?headerSearchOption=phrase}{IFRS Foundation --- Conceptual Framework for Financial Reporting} \\
\end{longtable}
\end{center}
\needspace{5\baselineskip}
\section[B]{B}
\begin{center}
\small
\setlength{\LTpre}{0.4em}
\setlength{\LTpost}{0.8em}
\begin{longtable}{>{\RaggedRight\arraybackslash}p{0.180\linewidth}>{\RaggedRight\arraybackslash}p{0.400\linewidth}>{\RaggedRight\arraybackslash}p{0.300\linewidth}}
\rowcolor{DeepNavy}
\color{white}\bfseries Term & \color{white}\bfseries Definition & \color{white}\bfseries Sources \\
\endfirsthead
\rowcolor{DeepNavy}
\color{white}\bfseries Term & \color{white}\bfseries Definition & \color{white}\bfseries Sources \\
\endhead
\textbf{Balance sheet} & A statement of financial position presenting assets, liabilities, and equity at a specified date. & \href{https://www.ifrs.org/issued-standards/list-of-standards/ias-1-presentation-of-financial-statements.html/}{IFRS Foundation --- IAS 1 Presentation of Financial Statements} \\
\textbf{Benchmark} & A baseline, target, historical result, or comparable entity used for interpretation. & \href{https://www.bdc.ca/en/articles-tools/entrepreneur-toolkit/templates-business-guides/glossary/ratio-analysis}{Business Development Bank of Canada --- Ratio analysis} \\
\textbf{BR\&E} & Business retention and expansion; a relationship, intelligence, and action process focused on existing firms. & \href{https://www.ontario.ca/page/business-retention-and-expansion-program}{Ontario Ministry of Agriculture, Food and Agribusiness --- Business Retention and Expansion Program} \\
\textbf{Business attraction} & Effort to identify and secure compatible new investment or business activity. & \href{https://www.worldbank.org/en/topic/investment-climate/brief/investment-policy-and-promotion}{World Bank --- Investment Policy and Promotion} \\
\textbf{Business climate} & The regulatory, cost, infrastructure, service, relationship, and market environment experienced by firms. & \href{https://www.worldbank.org/en/businessready/about-us/faq}{World Bank --- Business Ready frequently asked questions} \\
\end{longtable}
\end{center}
\needspace{5\baselineskip}
\section[C]{C}
\begin{center}
\small
\setlength{\LTpre}{0.4em}
\setlength{\LTpost}{0.8em}
\begin{longtable}{>{\RaggedRight\arraybackslash}p{0.180\linewidth}>{\RaggedRight\arraybackslash}p{0.400\linewidth}>{\RaggedRight\arraybackslash}p{0.300\linewidth}}
\rowcolor{DeepNavy}
\color{white}\bfseries Term & \color{white}\bfseries Definition & \color{white}\bfseries Sources \\
\endfirsthead
\rowcolor{DeepNavy}
\color{white}\bfseries Term & \color{white}\bfseries Definition & \color{white}\bfseries Sources \\
\endhead
\textbf{Capital budget} & Plan for acquisition and financing of long-lived assets. & \href{https://www.tbs-sct.canada.ca/pol/doc-eng.aspx?id=32574}{Treasury Board of Canada Secretariat --- Directive on Corporate Plans, Expenditures and Budgets} \\
\textbf{Cash flow} & Inflows and outflows of cash and cash equivalents during a period, commonly classified as operating, investing, or financing. & \href{https://www.ifrs.org/issued-standards/list-of-standards/ias-7-statement-of-cash-flows.html/}{IFRS Foundation --- IAS 7 Statement of Cash Flows} \\
\textbf{Cluster} & Interconnected firms, suppliers, institutions, workers, and supporting organizations in a related field and geography. & \href{https://www.eda.gov/about/economic-development-glossary}{U.S. Economic Development Administration --- Economic Development Glossary} \\
\textbf{Community profile} & Current, sourced information used by firms and investors to assess a location. & \href{https://www12.statcan.gc.ca/census-recensement/2021/dp-pd/prof/index.cfm?Lang=E}{Statistics Canada --- 2021 Census Profile} \\
\textbf{Competitive analysis} & Quantitative and qualitative assessment of the community's economic position, assets, performance, and opportunities. & \href{https://www.trade.gov/selectusa-reports-and-publications}{SelectUSA --- Reports and Publications} \\
\textbf{Conflict of interest} & A situation in which a private interest, relationship, or competing duty impairs, may impair, or could reasonably appear to impair objective professional judgment. & \href{https://edac.ca/about/code-of-ethics/}{Economic Developers Association of Canada --- Code of Ethics} \\
\textbf{Conversion rate} & Proportion of cases that move from one defined stage to the next. & \href{https://www.ontario.ca/document/performance-measurement-agriculture-agri-food-and-economic-development-organizations/appendix-5-examples-identifying-performance-measures-and-data}{Ontario --- Examples of identifying performance measures and data} \\
\textbf{Cost-benefit analysis} & Comparison of social benefits and costs of alternatives, generally across a defined time horizon. & \href{https://www.tbs-sct.canada.ca/rtrap-parfa/analys/analys-eng.pdf}{Treasury Board of Canada Secretariat --- Canadian Cost-Benefit Analysis Guide} \\
\textbf{Current ratio} & Current assets divided by current liabilities. & \href{https://www.bdc.ca/en/articles-tools/money-finance/manage-finances/financial-ratios-4-ways-assess-business}{Business Development Bank of Canada --- Four ways to assess your business using financial ratios} \\
\end{longtable}
\end{center}
\needspace{5\baselineskip}
\section[D]{D}
\begin{center}
\small
\setlength{\LTpre}{0.4em}
\setlength{\LTpost}{0.8em}
\begin{longtable}{>{\RaggedRight\arraybackslash}p{0.180\linewidth}>{\RaggedRight\arraybackslash}p{0.400\linewidth}>{\RaggedRight\arraybackslash}p{0.300\linewidth}}
\rowcolor{DeepNavy}
\color{white}\bfseries Term & \color{white}\bfseries Definition & \color{white}\bfseries Sources \\
\endfirsthead
\rowcolor{DeepNavy}
\color{white}\bfseries Term & \color{white}\bfseries Definition & \color{white}\bfseries Sources \\
\endhead
\textbf{Days sales outstanding} & Estimated average collection period, normally calculated using average accounts receivable divided by average daily net credit sales. & \href{https://www.bdc.ca/en/articles-tools/money-finance/manage-finances/financial-ratios-what-are-how-use}{Business Development Bank of Canada --- Financial ratios: what they are and how to use them} \\
\textbf{Deadweight} & Benefits that likely would have occurred without the intervention. & \href{https://www.gov.uk/government/publications/the-green-book-appraisal-and-evaluation-in-central-government/the-green-book-2026}{UK Government --- The Green Book: appraisal and evaluation in central government} \\
\textbf{Debt-to-equity} & A leverage ratio comparing a defined debt measure with a defined equity measure. Sources variously use interest-bearing debt or total liabilities, so state the convention. & \href{https://www.bdc.ca/en/articles-tools/money-finance/manage-finances/financial-ratios-what-are-how-use}{Business Development Bank of Canada --- Financial ratios: what they are and how to use them} \\
\textbf{Development impact analysis} & Assessment of fiscal, environmental, socio-economic, transportation, and other consequences of development. & \href{https://www.canada.ca/en/impact-assessment-agency/services/policy-guidance/practitioners-guide-impact-assessment-act/analyzing-health-social-economic-effects-impact-assessment-act.html}{Impact Assessment Agency of Canada --- Analyzing health, social and economic effects} \\
\textbf{Differential shift} & The shift-share component associated with local industry performance relative to the same industry in the reference area; also called the local, regional, or competitive effect. & \href{https://www.dot.mn.gov/ofrw/freight/PDF/d1plan/wp2.pdf}{Minnesota Department of Transportation --- Regional Economic Analysis and Shift-Share Analysis} \\
\textbf{Direct effect} & Initial activity generated by the project or organization being analyzed. & \href{https://www.statcan.gc.ca/en/statistical-programs/document/5115\_D6\_T9\_V1}{Statistics Canada --- Guide to Using the Input-Output Simulation Model} \\
\textbf{Displacement} & Activity that replaces or moves existing activity rather than creating net new benefit. & \href{https://www.gov.uk/government/publications/the-green-book-appraisal-and-evaluation-in-central-government/the-green-book-2026}{UK Government --- The Green Book: appraisal and evaluation in central government} \\
\textbf{Diversification} & Reduced dependence on a narrow set of sectors, firms, customers, markets, or resources. & \href{https://www.eda.gov/resources/comprehensive-economic-development-strategy/content/economic-resilience}{U.S. Economic Development Administration --- Economic Resilience} \\
\end{longtable}
\end{center}
\needspace{5\baselineskip}
\section[E]{E}
\begin{center}
\small
\setlength{\LTpre}{0.4em}
\setlength{\LTpost}{0.8em}
\begin{longtable}{>{\RaggedRight\arraybackslash}p{0.180\linewidth}>{\RaggedRight\arraybackslash}p{0.400\linewidth}>{\RaggedRight\arraybackslash}p{0.300\linewidth}}
\rowcolor{DeepNavy}
\color{white}\bfseries Term & \color{white}\bfseries Definition & \color{white}\bfseries Sources \\
\endfirsthead
\rowcolor{DeepNavy}
\color{white}\bfseries Term & \color{white}\bfseries Definition & \color{white}\bfseries Sources \\
\endhead
\textbf{Economic base} & Industries and activities that bring income into a region by selling to markets outside it or attracting external spending. Economic-base analysis contrasts these with activity serving primarily local demand. & \href{https://www.ers.usda.gov/data-products/county-typology-codes/documentation}{U.S. Department of Agriculture Economic Research Service --- County Typology Codes documentation} \\
\textbf{Economic impact} & A change in measures such as employment, labour income, output, or value added associated with an activity or intervention. Scope, geography, time, counterfactual, and direct, indirect, and induced effects should be stated. & \href{https://www.statcan.gc.ca/en/subjects-start/economic\_accounts/faq}{Statistics Canada --- Economic accounts frequently asked questions} \\
\textbf{Economic resilience} & Capacity to anticipate, withstand, adapt to, and recover from shocks. & \href{https://www.eda.gov/resources/comprehensive-economic-development-strategy/content/economic-resilience}{U.S. Economic Development Administration --- Economic Resilience} \\
\textbf{Equity} & In accounting, the residual interest in assets after deducting liabilities. In community analysis, fairness that recognizes different circumstances and seeks substantively fair access, process, and distribution of effects. State the intended meaning. & \href{https://www.ifrs.org/issued-standards/list-of-standards/conceptual-framework/?headerSearchOption=phrase}{IFRS Foundation --- Conceptual Framework for Financial Reporting}; \href{https://www.canada.ca/en/employment-social-development/corporate/portfolio/labour/programs/employment-equity/reports/act-review-overview-backgrounder-policy-issues.html}{Government of Canada --- Employment equity policy overview} \\
\textbf{Evaluation} & Systematic assessment of program design, delivery, results, and value. & \href{https://www.tbs-sct.canada.ca/pol/doc-eng.aspx?id=31300\&section=glossary}{Treasury Board of Canada Secretariat --- Policy on Results glossary} \\
\end{longtable}
\end{center}
\needspace{5\baselineskip}
\section[F]{F}
\begin{center}
\small
\setlength{\LTpre}{0.4em}
\setlength{\LTpost}{0.8em}
\begin{longtable}{>{\RaggedRight\arraybackslash}p{0.180\linewidth}>{\RaggedRight\arraybackslash}p{0.400\linewidth}>{\RaggedRight\arraybackslash}p{0.300\linewidth}}
\rowcolor{DeepNavy}
\color{white}\bfseries Term & \color{white}\bfseries Definition & \color{white}\bfseries Sources \\
\endfirsthead
\rowcolor{DeepNavy}
\color{white}\bfseries Term & \color{white}\bfseries Definition & \color{white}\bfseries Sources \\
\endhead
\textbf{FDI} & Foreign direct investment; cross-border investment associated with lasting interest and influence in an enterprise. & \href{https://www.oecd.org/en/publications/oecd-benchmark-definition-of-foreign-direct-investment-fifth-edition\_7f05c0a3-en/full-report/main-concepts-and-definitions-of-foreign-direct-investment\_c3240e7d.html}{OECD --- Benchmark Definition of Foreign Direct Investment, fifth edition} \\
\textbf{Fiscal impact} & Effect on public revenue, cost, capital, and financial position. & \href{https://www.calgary.ca/content/dam/www/cfod/finance/documents/corporate-economics/other-reports/municipal-fiscal-impact-model.pdf}{City of Calgary --- A Municipal Fiscal Impact Model} \\
\end{longtable}
\end{center}
\needspace{5\baselineskip}
\section[G]{G}
\begin{center}
\small
\setlength{\LTpre}{0.4em}
\setlength{\LTpost}{0.8em}
\begin{longtable}{>{\RaggedRight\arraybackslash}p{0.180\linewidth}>{\RaggedRight\arraybackslash}p{0.400\linewidth}>{\RaggedRight\arraybackslash}p{0.300\linewidth}}
\rowcolor{DeepNavy}
\color{white}\bfseries Term & \color{white}\bfseries Definition & \color{white}\bfseries Sources \\
\endfirsthead
\rowcolor{DeepNavy}
\color{white}\bfseries Term & \color{white}\bfseries Definition & \color{white}\bfseries Sources \\
\endhead
\textbf{Governance} & Structures and processes for direction, authority, accountability, risk, and oversight. & \href{https://www.tbs-sct.canada.ca/pol/doc-eng.aspx?id=25761}{Treasury Board of Canada Secretariat --- Guideline on Service Agreements} \\
\textbf{Gross margin} & Sales less cost of goods sold, divided by sales. & \href{https://www.bdc.ca/en/articles-tools/money-finance/manage-finances/financial-ratios-what-are-how-use}{Business Development Bank of Canada --- Financial ratios: what they are and how to use them} \\
\end{longtable}
\end{center}
\needspace{5\baselineskip}
\section[I]{I}
\begin{center}
\small
\setlength{\LTpre}{0.4em}
\setlength{\LTpost}{0.8em}
\begin{longtable}{>{\RaggedRight\arraybackslash}p{0.180\linewidth}>{\RaggedRight\arraybackslash}p{0.400\linewidth}>{\RaggedRight\arraybackslash}p{0.300\linewidth}}
\rowcolor{DeepNavy}
\color{white}\bfseries Term & \color{white}\bfseries Definition & \color{white}\bfseries Sources \\
\endfirsthead
\rowcolor{DeepNavy}
\color{white}\bfseries Term & \color{white}\bfseries Definition & \color{white}\bfseries Sources \\
\endhead
\textbf{Impact} & A significant positive or negative, intended or unintended, higher-level effect to which an intervention contributes. Evaluation frameworks differ on whether the term is reserved for longer-term change. & \href{https://www.oecd.org/en/publications/glossary-of-key-terms-in-evaluation-and-results-based-management-for-sustainable-development-second-edition\_632da462-en-fr-es.html}{OECD --- Glossary of Key Terms in Evaluation and Results-Based Management} \\
\textbf{Income statement} & Statement of revenue, expenses, and profit over a period. & \href{https://www.ifrs.org/issued-standards/list-of-standards/ias-1-presentation-of-financial-statements.html/}{IFRS Foundation --- IAS 1 Presentation of Financial Statements} \\
\textbf{Indirect effect} & Activity generated in suppliers by the direct activity. & \href{https://www.statcan.gc.ca/en/statistical-programs/document/5115\_D6\_T9\_V1}{Statistics Canada --- Guide to Using the Input-Output Simulation Model} \\
\textbf{Induced effect} & Activity associated with spending of labour income from direct and indirect activity. & \href{https://www.statcan.gc.ca/en/statistical-programs/document/5115\_D6\_T9\_V1}{Statistics Canada --- Guide to Using the Input-Output Simulation Model} \\
\textbf{Industrial mix effect} & Shift-share component associated with the reference industry's growth relative to overall reference growth. & \href{https://www.dot.mn.gov/ofrw/freight/PDF/d1plan/wp2.pdf}{Minnesota Department of Transportation --- Regional Economic Analysis and Shift-Share Analysis} \\
\textbf{Input} & Resource used by a program, such as staff, money, data, or facilities. & \href{https://www.canada.ca/en/treasury-board-secretariat/services/audit-evaluation/evaluation-government-canada/evaluation-101-backgrounder.html}{Treasury Board of Canada Secretariat --- Evaluation 101} \\
\textbf{Interest coverage} & An earnings measure divided by interest expense; EBIT divided by interest expense is common, but definitions vary. & \href{https://www.bdc.ca/en/articles-tools/money-finance/manage-finances/financial-ratios-what-are-how-use}{Business Development Bank of Canada --- Financial ratios: what they are and how to use them} \\
\textbf{Inventory turnover} & Cost of goods sold divided by average inventory. & \href{https://www.bdc.ca/en/articles-tools/money-finance/manage-finances/financial-ratios-what-are-how-use}{Business Development Bank of Canada --- Financial ratios: what they are and how to use them} \\
\end{longtable}
\end{center}
\needspace{5\baselineskip}
\section[L]{L}
\begin{center}
\small
\setlength{\LTpre}{0.4em}
\setlength{\LTpost}{0.8em}
\begin{longtable}{>{\RaggedRight\arraybackslash}p{0.180\linewidth}>{\RaggedRight\arraybackslash}p{0.400\linewidth}>{\RaggedRight\arraybackslash}p{0.300\linewidth}}
\rowcolor{DeepNavy}
\color{white}\bfseries Term & \color{white}\bfseries Definition & \color{white}\bfseries Sources \\
\endfirsthead
\rowcolor{DeepNavy}
\color{white}\bfseries Term & \color{white}\bfseries Definition & \color{white}\bfseries Sources \\
\endhead
\textbf{Land-use bylaw or zoning bylaw} & Regulation that establishes land uses and development standards; terminology varies by jurisdiction. & \href{https://www.ontario.ca/document/citizens-guide-land-use-planning/zoning-bylaws}{Ontario --- Zoning bylaws} \\
\textbf{Lead} & A possible opportunity requiring validation. & \href{https://www.salesforce.com/sales/prospecting/}{Salesforce --- What is sales prospecting?} \\
\textbf{Leakage} & Benefits that flow outside the intended area or population. & \href{https://www.gov.uk/government/publications/the-green-book-appraisal-and-evaluation-in-central-government/the-green-book-2026}{UK Government --- The Green Book: appraisal and evaluation in central government} \\
\textbf{Liabilities} & Present obligations to transfer an economic resource as a result of past events. & \href{https://www.ifrs.org/issued-standards/list-of-standards/conceptual-framework/?headerSearchOption=phrase}{IFRS Foundation --- Conceptual Framework for Financial Reporting} \\
\textbf{Liquidity} & Ability to meet near-term obligations. & \href{https://www.bdc.ca/en/articles-tools/money-finance/manage-finances/financial-ratios-4-ways-assess-business}{Business Development Bank of Canada --- Four ways to assess your business using financial ratios} \\
\textbf{Location quotient} & Ratio comparing an industry's share of a regional measure with the same industry's share of that measure in a reference economy. Employment is common, but earnings, output, or GDP may also be used. & \href{https://www150.statcan.gc.ca/n1/pub/21-006-x/2008007/def-eng.htm}{Statistics Canada --- Glossary of terms: location quotient} \\
\textbf{Logic model} & Connection among inputs, activities, outputs, and outcomes. & \href{https://www.canada.ca/en/treasury-board-secretariat/services/audit-evaluation/evaluation-government-canada/evaluation-101-backgrounder.html}{Treasury Board of Canada Secretariat --- Evaluation 101} \\
\end{longtable}
\end{center}
\needspace{5\baselineskip}
\section[M]{M}
\begin{center}
\small
\setlength{\LTpre}{0.4em}
\setlength{\LTpost}{0.8em}
\begin{longtable}{>{\RaggedRight\arraybackslash}p{0.180\linewidth}>{\RaggedRight\arraybackslash}p{0.400\linewidth}>{\RaggedRight\arraybackslash}p{0.300\linewidth}}
\rowcolor{DeepNavy}
\color{white}\bfseries Term & \color{white}\bfseries Definition & \color{white}\bfseries Sources \\
\endfirsthead
\rowcolor{DeepNavy}
\color{white}\bfseries Term & \color{white}\bfseries Definition & \color{white}\bfseries Sources \\
\endhead
\textbf{Market-to-book} & Market price per share divided by book value per share. & \href{https://www.cfainstitute.org/insights/professional-learning/refresher-readings/2026/market-based-valuation-price-enterprise-value-multiples}{CFA Institute --- Market-Based Valuation: Price and Enterprise Value Multiples} \\
\textbf{Marketing} & Targeting, positioning, communication, and lead-generation activity. & \href{https://www.bdc.ca/en/articles-tools/entrepreneur-toolkit/templates-business-guides/glossary/marketing-plan}{Business Development Bank of Canada --- Marketing plan} \\
\textbf{Mentoring} & Developmental relationship in which experience supports another person's capability and judgment. & \href{https://justice.canada.ca/eng/rp-pr/other-autre/rr12\_1/p1.html}{Justice Canada --- Three Years On: Mentoring at the Department of Justice and the Public Prosecution Service of Canada} \\
\end{longtable}
\end{center}
\needspace{5\baselineskip}
\section[N]{N}
\begin{center}
\small
\setlength{\LTpre}{0.4em}
\setlength{\LTpost}{0.8em}
\begin{longtable}{>{\RaggedRight\arraybackslash}p{0.180\linewidth}>{\RaggedRight\arraybackslash}p{0.400\linewidth}>{\RaggedRight\arraybackslash}p{0.300\linewidth}}
\rowcolor{DeepNavy}
\color{white}\bfseries Term & \color{white}\bfseries Definition & \color{white}\bfseries Sources \\
\endfirsthead
\rowcolor{DeepNavy}
\color{white}\bfseries Term & \color{white}\bfseries Definition & \color{white}\bfseries Sources \\
\endhead
\textbf{NAICS} & North American Industry Classification System. & \href{https://www.statcan.gc.ca/en/subjects/standard/naics/2022/v1/introduction}{Statistics Canada --- North American Industry Classification System introduction} \\
\textbf{Net margin} & Net income divided by sales. & \href{https://www.bdc.ca/en/articles-tools/money-finance/manage-finances/financial-ratios-what-are-how-use}{Business Development Bank of Canada --- Financial ratios: what they are and how to use them} \\
\end{longtable}
\end{center}
\needspace{5\baselineskip}
\section[O]{O}
\begin{center}
\small
\setlength{\LTpre}{0.4em}
\setlength{\LTpost}{0.8em}
\begin{longtable}{>{\RaggedRight\arraybackslash}p{0.180\linewidth}>{\RaggedRight\arraybackslash}p{0.400\linewidth}>{\RaggedRight\arraybackslash}p{0.300\linewidth}}
\rowcolor{DeepNavy}
\color{white}\bfseries Term & \color{white}\bfseries Definition & \color{white}\bfseries Sources \\
\endfirsthead
\rowcolor{DeepNavy}
\color{white}\bfseries Term & \color{white}\bfseries Definition & \color{white}\bfseries Sources \\
\endhead
\textbf{Official plan or municipal plan} & A statutory or policy plan setting long-term land-use and growth direction. Its name, required content, approval process, and legal effect vary by jurisdiction. & \href{https://www.ontario.ca/document/citizens-guide-land-use-planning/official-plans}{Ontario --- Official plans} \\
\textbf{Operating budget} & Plan for recurring revenues and expenditures. & \href{https://www.tbs-sct.canada.ca/pol/doc-eng.aspx?id=32574}{Treasury Board of Canada Secretariat --- Directive on Corporate Plans, Expenditures and Budgets} \\
\textbf{Outcome} & A short- or medium-term change in condition, behaviour, capacity, or performance that an intervention is expected to influence through its activities and outputs. & \href{https://www.oecd.org/en/publications/glossary-of-key-terms-in-evaluation-and-results-based-management-for-sustainable-development-second-edition\_632da462-en-fr-es.html}{OECD --- Glossary of Key Terms in Evaluation and Results-Based Management} \\
\textbf{Output} & A product, capital good, or service delivered by an intervention, such as firms served, workshops delivered, or reports completed. & \href{https://www.oecd.org/en/publications/glossary-of-key-terms-in-evaluation-and-results-based-management-for-sustainable-development-second-edition\_632da462-en-fr-es.html}{OECD --- Glossary of Key Terms in Evaluation and Results-Based Management} \\
\end{longtable}
\end{center}
\needspace{5\baselineskip}
\section[P]{P}
\begin{center}
\small
\setlength{\LTpre}{0.4em}
\setlength{\LTpost}{0.8em}
\begin{longtable}{>{\RaggedRight\arraybackslash}p{0.180\linewidth}>{\RaggedRight\arraybackslash}p{0.400\linewidth}>{\RaggedRight\arraybackslash}p{0.300\linewidth}}
\rowcolor{DeepNavy}
\color{white}\bfseries Term & \color{white}\bfseries Definition & \color{white}\bfseries Sources \\
\endfirsthead
\rowcolor{DeepNavy}
\color{white}\bfseries Term & \color{white}\bfseries Definition & \color{white}\bfseries Sources \\
\endhead
\textbf{P3} & Public-private partnership; an arrangement allocating project functions and risks between public and private parties. & \href{https://www.infrastructureontario.ca/en/what-we-do/capital-delivery/faqs---public-private-partnerships-p3s/}{Infrastructure Ontario --- Public-private partnerships frequently asked questions} \\
\textbf{Performance measure} & Indicator used to monitor implementation, efficiency, output, outcome, or impact. & \href{https://www.tbs-sct.canada.ca/pol/doc-eng.aspx?id=31300\&section=glossary}{Treasury Board of Canada Secretariat --- Policy on Results glossary} \\
\textbf{Price-earnings ratio} & Market price per share divided by earnings per share. & \href{https://www.investor.gov/introduction-investing/investing-basics/glossary/price-earnings-pe-ratio}{Investor.gov --- Price-earnings ratio} \\
\textbf{Project} & In this guide's investment funnel, a qualified and actively managed investment opportunity. Pipeline terminology varies among organizations and customer-relationship systems. & \href{https://www.trade.gov/selectusa-investor}{SelectUSA --- Investor services} \\
\textbf{Prospect} & A lead judged sufficiently suitable for active follow-up. Sales and investment-pipeline terminology varies by organization. & \href{https://www.salesforce.com/sales/prospecting/}{Salesforce --- What is sales prospecting?} \\
\end{longtable}
\end{center}
\needspace{5\baselineskip}
\section[Q]{Q}
\begin{center}
\small
\setlength{\LTpre}{0.4em}
\setlength{\LTpost}{0.8em}
\begin{longtable}{>{\RaggedRight\arraybackslash}p{0.180\linewidth}>{\RaggedRight\arraybackslash}p{0.400\linewidth}>{\RaggedRight\arraybackslash}p{0.300\linewidth}}
\rowcolor{DeepNavy}
\color{white}\bfseries Term & \color{white}\bfseries Definition & \color{white}\bfseries Sources \\
\endfirsthead
\rowcolor{DeepNavy}
\color{white}\bfseries Term & \color{white}\bfseries Definition & \color{white}\bfseries Sources \\
\endhead
\textbf{Qualified lead/project} & A verified opportunity meeting the organization's documented criteria, normally including fit, need, decision authority, resources, and timing. & \href{https://www.salesforce.com/blog/sales/sales-qualified-lead/}{Salesforce --- What is a sales-qualified lead?} \\
\textbf{Quick ratio} & Quick assets divided by current liabilities. Quick assets commonly include cash, short-term investments, and receivables; simplified formulas may subtract inventory and less-liquid prepaids from current assets. & \href{https://www.bdc.ca/en/articles-tools/money-finance/manage-finances/financial-ratios-4-ways-assess-business}{Business Development Bank of Canada --- Four ways to assess your business using financial ratios} \\
\end{longtable}
\end{center}
\needspace{5\baselineskip}
\section[R]{R}
\begin{center}
\small
\setlength{\LTpre}{0.4em}
\setlength{\LTpost}{0.8em}
\begin{longtable}{>{\RaggedRight\arraybackslash}p{0.180\linewidth}>{\RaggedRight\arraybackslash}p{0.400\linewidth}>{\RaggedRight\arraybackslash}p{0.300\linewidth}}
\rowcolor{DeepNavy}
\color{white}\bfseries Term & \color{white}\bfseries Definition & \color{white}\bfseries Sources \\
\endfirsthead
\rowcolor{DeepNavy}
\color{white}\bfseries Term & \color{white}\bfseries Definition & \color{white}\bfseries Sources \\
\endhead
\textbf{Reference-area growth effect} & Shift-share component associated with overall employment growth in the reference economy. & \href{https://www.dot.mn.gov/ofrw/freight/PDF/d1plan/wp2.pdf}{Minnesota Department of Transportation --- Regional Economic Analysis and Shift-Share Analysis} \\
\textbf{Return on assets} & Net income divided by average total assets. & \href{https://www.bdc.ca/en/articles-tools/money-finance/manage-finances/financial-ratios-what-are-how-use}{Business Development Bank of Canada --- Financial ratios: what they are and how to use them} \\
\textbf{Return on equity} & Net income available to common shareholders divided by average common equity. & \href{https://www.bdc.ca/en/articles-tools/money-finance/manage-finances/financial-ratios-what-are-how-use}{Business Development Bank of Canada --- Financial ratios: what they are and how to use them} \\
\textbf{Rights-holder} & A person or collective holding legal, constitutional, treaty, inherent, or other recognized rights relevant to a decision. In the Crown's duty-to-consult context, the term refers to Indigenous rights-holding collectives, not merely interested parties. & \href{https://justice.canada.ca/eng/cons/def.html}{Justice Canada --- Consultation terminology}; \href{https://our-languages.canada.ca/en/blogue-blog/rethinking-the-word-stakeholder-eng}{Language Portal of Canada --- Rethinking the word ``stakeholder''} \\
\end{longtable}
\end{center}
\needspace{5\baselineskip}
\section[S]{S}
\begin{center}
\small
\setlength{\LTpre}{0.4em}
\setlength{\LTpost}{0.8em}
\begin{longtable}{>{\RaggedRight\arraybackslash}p{0.180\linewidth}>{\RaggedRight\arraybackslash}p{0.400\linewidth}>{\RaggedRight\arraybackslash}p{0.300\linewidth}}
\rowcolor{DeepNavy}
\color{white}\bfseries Term & \color{white}\bfseries Definition & \color{white}\bfseries Sources \\
\endfirsthead
\rowcolor{DeepNavy}
\color{white}\bfseries Term & \color{white}\bfseries Definition & \color{white}\bfseries Sources \\
\endhead
\textbf{Sales} & Relationship and decision work that moves a qualified project toward commitment. & \href{https://www.bdc.ca/en/articles-tools/entrepreneur-toolkit/templates-business-guides/sales-and-marketing-toolkit/sales-and-marketing-toolkit-download}{Business Development Bank of Canada --- Sales and marketing toolkit} \\
\textbf{Sector strategy} & Coordinated plan to strengthen a defined industry or economic field. & \href{https://www.oecd.org/en/topics/sub-issues/strategic-planning-and-regional-development.html}{OECD --- Strategic planning and regional development} \\
\textbf{Shift-share} & Method decomposing local employment change into reference growth, industry mix, and local shift. & \href{https://www.dot.mn.gov/ofrw/freight/PDF/d1plan/wp2.pdf}{Minnesota Department of Transportation --- Regional Economic Analysis and Shift-Share Analysis} \\
\textbf{Site selection} & Business decision process used to compare, eliminate, and select locations. & \href{https://www.trade.gov/selectusa-investor-guide}{SelectUSA --- Investor Guide} \\
\textbf{SMART} & Specific, measurable, achievable, relevant, and time-bound. Variants use words such as realistic, results-focused, or audience-specific; define the version being used. & \href{https://www.canada.ca/en/environment-climate-change/services/environmental-funding/tools-for-applying/writing-smart-objectives.html}{Environment and Climate Change Canada --- Writing SMART objectives} \\
\textbf{Stakeholder} & A person, group, or organization affected by, interested in, or able to influence a decision. Do not use the term as a substitute for Indigenous rights-holder. & \href{https://www.canada.ca/en/health-canada/services/publications/health-system-services/health-canada-public-health-agency-canada-guidelines-public-engagement/glossary-terms.html}{Health Canada and Public Health Agency of Canada --- Public engagement glossary}; \href{https://our-languages.canada.ca/en/blogue-blog/rethinking-the-word-stakeholder-eng}{Language Portal of Canada --- Rethinking the word ``stakeholder''} \\
\textbf{Strategic plan} & Evidence-based framework that links mission, priorities, actions, resources, and measurement. & \href{https://www.oecd.org/en/topics/sub-issues/strategic-planning-and-regional-development.html}{OECD --- Strategic planning and regional development} \\
\textbf{SWOT} & Strengths, weaknesses, opportunities, and threats; usually strengths and weaknesses are internal, while opportunities and threats arise externally. & \href{https://www.canada.ca/en/revenue-agency/corporate/about-canada-revenue-agency-cra/sustainable-development/sustainable-development-strategies/sustainable-development-strategy-2007-2010-10.html}{Canada Revenue Agency --- Environmental Scan and SWOT Analysis} \\
\end{longtable}
\end{center}
\needspace{5\baselineskip}
\section[T]{T}
\begin{center}
\small
\setlength{\LTpre}{0.4em}
\setlength{\LTpost}{0.8em}
\begin{longtable}{>{\RaggedRight\arraybackslash}p{0.180\linewidth}>{\RaggedRight\arraybackslash}p{0.400\linewidth}>{\RaggedRight\arraybackslash}p{0.300\linewidth}}
\rowcolor{DeepNavy}
\color{white}\bfseries Term & \color{white}\bfseries Definition & \color{white}\bfseries Sources \\
\endfirsthead
\rowcolor{DeepNavy}
\color{white}\bfseries Term & \color{white}\bfseries Definition & \color{white}\bfseries Sources \\
\endhead
\textbf{Target sector} & Industry or business function selected for focused development based on evidence and fit. & \href{https://www.trade.gov/selectusa-target-industries-dashboard}{SelectUSA --- Target Industries Dashboard} \\
\textbf{Theory of change} & Explanation of how and why activities are expected to create outcomes. & \href{https://www.canada.ca/en/treasury-board-secretariat/services/audit-evaluation/evaluation-government-canada/evaluation-101-backgrounder.html}{Treasury Board of Canada Secretariat --- Evaluation 101} \\
\textbf{Triple bottom line} & A framework that considers economic, social, and environmental performance or value rather than financial performance alone. & \href{https://ised-isde.canada.ca/site/corporate-social-responsibility/en/sme-sustainability-roadmap}{Innovation, Science and Economic Development Canada --- SME Sustainability Roadmap} \\
\end{longtable}
\end{center}
\needspace{5\baselineskip}
\section[V]{V}
\begin{center}
\small
\setlength{\LTpre}{0.4em}
\setlength{\LTpost}{0.8em}
\begin{longtable}{>{\RaggedRight\arraybackslash}p{0.180\linewidth}>{\RaggedRight\arraybackslash}p{0.400\linewidth}>{\RaggedRight\arraybackslash}p{0.300\linewidth}}
\rowcolor{DeepNavy}
\color{white}\bfseries Term & \color{white}\bfseries Definition & \color{white}\bfseries Sources \\
\endfirsthead
\rowcolor{DeepNavy}
\color{white}\bfseries Term & \color{white}\bfseries Definition & \color{white}\bfseries Sources \\
\endhead
\textbf{Value proposition} & Evidence-based explanation of why a target should choose or engage with a community. & \href{https://www.trade.gov/selectusa-reports-and-publications}{SelectUSA --- Reports and Publications} \\
\end{longtable}
\end{center}
\needspace{5\baselineskip}
\section[W]{W}
\begin{center}
\small
\setlength{\LTpre}{0.4em}
\setlength{\LTpost}{0.8em}
\begin{longtable}{>{\RaggedRight\arraybackslash}p{0.180\linewidth}>{\RaggedRight\arraybackslash}p{0.400\linewidth}>{\RaggedRight\arraybackslash}p{0.300\linewidth}}
\rowcolor{DeepNavy}
\color{white}\bfseries Term & \color{white}\bfseries Definition & \color{white}\bfseries Sources \\
\endfirsthead
\rowcolor{DeepNavy}
\color{white}\bfseries Term & \color{white}\bfseries Definition & \color{white}\bfseries Sources \\
\endhead
\textbf{Working capital} & Current assets minus current liabilities. & \href{https://www.bdc.ca/en/articles-tools/money-finance/manage-finances/using-working-capital-ratio}{Business Development Bank of Canada --- Using the working capital ratio} \\
\end{longtable}
\end{center}

\part{Sources and Use}
\chapter[Public Sources and Disclaimer]{\chaptericon{library}Public Sources and Disclaimer}
\label{chap:sources}

\needspace{5\baselineskip}
\marginsection{calendar-check}{Verification date}
The public requirements and links in this guide were checked on August 12, 2026. Certification rules can change, so candidates should verify them again before registering or writing.


\needspace{5\baselineskip}
\section[Official EDAC sources]{Official EDAC sources}
\begin{center}
\small
\setlength{\LTpre}{0.4em}
\setlength{\LTpost}{0.8em}
\begin{longtable}{>{\RaggedRight\arraybackslash}p{0.290\linewidth}>{\RaggedRight\arraybackslash}p{0.610\linewidth}}
\rowcolor{DeepNavy}
\color{white}\bfseries Public source & \color{white}\bfseries How it is used in this guide \\
\endfirsthead
\rowcolor{DeepNavy}
\color{white}\bfseries Public source & \color{white}\bfseries How it is used in this guide \\
\endhead
\href{https://edac.ca/professional-development/ecd-designation/}{Ec.D. Designation} & Public eligibility routes, professional-development expectations, and the designation's three current pillars \\
\href{https://exam.edac.ca/exam-prep/}{Ec.D. Exam Preparation} & Membership and eligibility summary, preparation recommendations, technical requirements, exam timing, section structure, and marks \\
\href{https://exam.edac.ca/registration-process/}{Exam Registration Process} & Public registration and scheduling sequence \\
\href{https://exam.edac.ca/policies/}{Exam Policies} & Candidate agreement, result handling, Board approval, re-write policy, fees, and membership requirements \\
\href{https://edac.ca/about/code-of-ethics/}{EDAC Code of Ethics} & Professional duties concerning competence, integrity, confidentiality, conflicts, public interest, and professional conduct \\
\href{https://edac.ca/professional-development/ec-d-recertification/}{Ec.D. Re-certification} & Continuing membership, professional activity, and re-certification expectations \\
\end{longtable}
\end{center}
EDAC is the authority for eligibility and examination rules. Its public designation page describes a baseline route and alternatives, while the separate exam portal summarizes the three-year and 45-point route. Administrative details also differ across pages. Candidates should therefore ask EDAC to confirm their individual eligibility, fee, schedule, permitted materials, accommodation process, technical setup, and current contact information.


\needspace{5\baselineskip}
\section[Canadian economic and community data]{Canadian economic and community data}
\begin{center}
\small
\setlength{\LTpre}{0.4em}
\setlength{\LTpost}{0.8em}
\begin{longtable}{>{\RaggedRight\arraybackslash}p{0.290\linewidth}>{\RaggedRight\arraybackslash}p{0.610\linewidth}}
\rowcolor{DeepNavy}
\color{white}\bfseries Public source & \color{white}\bfseries Useful applications \\
\endfirsthead
\rowcolor{DeepNavy}
\color{white}\bfseries Public source & \color{white}\bfseries Useful applications \\
\endhead
\href{https://www12.statcan.gc.ca/census-recensement/2021/dp-pd/prof/index.cfm?Lang=E}{Statistics Canada Census Profile} & Population, age, income, education, housing, commuting, language, immigration, and labour-force profiles \\
\href{https://www12.statcan.gc.ca/census-recensement/2021/dp-pd/dv-vd/cpdv-vdpr/cpdv-vdpr-eng.cfm}{Statistics Canada Census Program Data Viewer} & Geographic comparison and visualization of census indicators \\
\href{https://www.statcan.gc.ca/en/subjects-start/labour\_}{Statistics Canada Labour Statistics} & Employment, unemployment, wages, vacancies, hours, productivity, and labour-market trends \\
\href{https://www23.statcan.gc.ca/imdb-bmdi/pub/1105-eng.htm}{Statistics Canada Business Register and Canadian Business Counts} & Business-location and establishment counts by industry, geography, and employment size \\
\href{https://ised-isde.canada.ca/site/trade-data-online/en}{ISED Trade Data Online} & Canadian imports and exports by product, industry, country, province, and territory \\
\href{https://innovation.ised-isde.canada.ca/s/?language=en\_CA}{Canada Business Benefits Finder} & Public programs and services potentially relevant to firms and entrepreneurs \\
\end{longtable}
\end{center}
National sources should be supplemented with current provincial, territorial, Indigenous, municipal, regional, utility, real-estate, workforce, and local business information.


\needspace{5\baselineskip}
\section[Contemporary-issue sources]{Contemporary-issue sources}
\begin{center}
\small
\setlength{\LTpre}{0.4em}
\setlength{\LTpost}{0.8em}
\begin{longtable}{>{\RaggedRight\arraybackslash}p{0.290\linewidth}>{\RaggedRight\arraybackslash}p{0.610\linewidth}}
\rowcolor{DeepNavy}
\color{white}\bfseries Public source & \color{white}\bfseries Topics supported \\
\endfirsthead
\rowcolor{DeepNavy}
\color{white}\bfseries Public source & \color{white}\bfseries Topics supported \\
\endhead
\href{https://www.bankofcanada.ca/publications/mpr/}{Bank of Canada Monetary Policy Report} & Inflation, interest rates, growth, investment, household demand, trade, productivity, and economic risk \\
\href{https://www.bankofcanada.ca/publications/mpr/mpr-2026-07-15/}{July 2026 Monetary Policy Report} & A dated economic snapshot for contemporary-issue practice \\
\href{https://www.cmhc-schl.gc.ca/media-newsroom/news-releases/2026/economic-uncertainty-continue-weighing-housing-market-cmhc}{CMHC 2026 Housing Market Outlook update} & Housing supply, sales, prices, rental conditions, uncertainty, and regional differences \\
\href{https://www.canada.ca/en/department-finance/programs/international-trade-finance-policy/canadas-tariff-responses.html}{Government of Canada tariff responses} & Current trade measures and Canadian policy responses \\
\href{https://www.canada.ca/en/services/indigenous-peoples/business-and-economic-development-indigenous-peoples.html}{Business and economic development for Indigenous Peoples} & Federal programs and entry points related to Indigenous entrepreneurship, business development, and economic participation \\
\href{https://www.canada.ca/en/services/environment/weather/climatechange/climate-plan/national-adaptation-strategy/full-strategy.html}{Canada's National Adaptation Strategy} & Climate resilience, infrastructure, disaster risk, health, nature, and economic adaptation \\
\end{longtable}
\end{center}
The contemporary-issues chapter uses these sources to identify themes, not to predict the examination question. Refresh the evidence close to the exam date.


\needspace{5\baselineskip}
\section[Interpretation notes]{Interpretation notes}
\begin{itemize}
\item Every substantive chapter contains a topic-specific \textbf{Learn more} section.
\item The \hyperref[chap:glossary]{glossary} pairs every definition with one or more authoritative public references.
\item The currently published essay limit is a maximum of 3,000 words.
\item The public exam page labels eight short-answer sections A through H but does not publicly state the subject assigned to each section.
\item EDAC's current designation page identifies business retention and expansion, investment attraction, and strategic planning as three pillars. The exam page does not publicly map the eight short-answer sections to named topics, so this guide studies a broader professional knowledge base.
\item Planning systems and legal terminology differ across provinces and territories. Verify the relevant jurisdiction.
\item Market-to-book is market price per share divided by book value per share.
\item Days sales outstanding may use 360 or 365 days. State the convention and use it consistently with the benchmark.
\item Economic-development outcomes should be supported by current evidence rather than treated as universal percentages or guarantees.
\end{itemize}


\needspace{5\baselineskip}
\section[Mock-exam design]{Mock-exam design}
The multiple-choice, true/false, written, and calculation practice questions draw on the concepts, frameworks, terminology, and public references used throughout this study guide. Constructed scenarios and figures provide practice organized around the publicly described examination structure.


\needspace{5\baselineskip}
\marginsection{sparkles}{AI assistance disclosure}
Xiao prepared the underlying study notes and made the final decisions about the guide's scope, content, and wording. OpenAI's Codex, using the GPT-5.6 model, helped synthesize the notes, research additional public sources, edit and format the Markdown and LaTeX files, and produce the PDF edition.


AI-assisted content can contain errors or omissions. Xiao reviewed the guide and takes responsibility for its published content. Readers should verify examination requirements and substantive information with EDAC and the cited public sources.


\needspace{5\baselineskip}
\marginsection{user-round}{Feedback and contact}
This guide was created by \href{https://www.linkedin.com/in/xxiao27/}{Xiao}, who publishes projects on GitHub as \href{https://github.com/B1ackCat7}{B1ackCat7}. Constructive corrections, comments, and professional connections are welcome through \href{https://www.linkedin.com/in/xxiao27/}{LinkedIn}.


\needspace{5\baselineskip}
\marginsection{badge-check}{Copyright and use}
This guide contains original educational explanations and purpose-built mock questions. It links to authoritative public sources, while EDAC retains rights in its publications, examination material, name, marks, and artwork.


Copyright \textcopyright{} 2026 \href{https://www.linkedin.com/in/xxiao27/}{Xiao}. The original content in this guide is licensed under the \href{https://creativecommons.org/licenses/by/4.0/}{Creative Commons Attribution 4.0 International License} (CC BY 4.0). Readers may copy, redistribute, remix, transform, and build upon it for any purpose, including commercially, provided they give appropriate credit, link to the licence, and indicate whether changes were made. See \hyperref[chap:license]{LICENSE.md}.


The licence does not grant rights in EDAC's name, marks, artwork, publications, examination material, or any third-party content linked or separately identified in this guide.


\needspace{5\baselineskip}
\marginsection{shield-alert}{Disclaimer}
This is an independently created educational resource and is not endorsed by EDAC. It is not legal, financial, accounting, planning, investment, or examination advice. EDAC is the authority for certification rules. Applicable governments and rights-holders are the authorities for law, jurisdiction, and approval requirements.


\needspace{5\baselineskip}
\marginsection{external-link}{Learn more}
\begin{itemize}
\item Start with \href{https://edac.ca/professional-development/ecd-designation/}{EDAC's Ec.D. Designation page} and \href{https://exam.edac.ca/exam-prep/}{Exam Preparation page} for current official examination information.
\item Use the \hyperref[chap:about]{chapter-level reference sections} to research a topic in context.
\item Use the \hyperref[chap:glossary]{term-by-term glossary} when checking a definition or competing convention.
\end{itemize}


\chapter*{\chaptericon{badge-check}Creative Commons Attribution 4.0 International}
\addcontentsline{toc}{chapter}{Creative Commons Attribution 4.0 International}
\label{chap:license}

Copyright \textcopyright{} 2026 \href{https://www.linkedin.com/in/xxiao27/}{Xiao}


Except where otherwise noted, the original content in this repository is licensed under the Creative Commons Attribution 4.0 International Public License (CC BY 4.0).


You are free to share and adapt the licensed material for any purpose, including commercially, subject to the licence terms. Those terms include giving appropriate credit, providing a link to the licence, and indicating whether changes were made.


\begin{itemize}
\item Licence summary: \href{https://creativecommons.org/licenses/by/4.0/}{https://creativecommons.org/licenses/by/4.0/}
\item Full legal code: \href{https://creativecommons.org/licenses/by/4.0/legalcode.en}{https://creativecommons.org/licenses/by/4.0/legalcode.en}
\end{itemize}


Suggested attribution:


\begin{guidequote}
\emph{EDAC Ec.D. Exam Study Guide} by \href{https://www.linkedin.com/in/xxiao27/}{Xiao}, published on GitHub by \href{https://github.com/B1ackCat7}{B1ackCat7} and licensed under \href{https://creativecommons.org/licenses/by/4.0/}{CC BY 4.0}. Changes were made.

\end{guidequote}


Remove ``Changes were made'' when redistributing an unmodified copy.


\needspace{5\baselineskip}
\section[Scope]{Scope}
This licence applies only to original content in this repository. It does not grant rights in third-party names, trademarks, artwork, publications, examination material, linked resources, or content that is quoted, incorporated, or otherwise separately identified. Those materials remain subject to their respective rights and terms.


The EDAC name, Ec.D. designation, associated marks, and EDAC materials belong to their respective rights-holders. Their appearance is for identification, commentary, and educational reference and does not imply sponsorship or endorsement.


This repository and its contents are provided as-is, without warranties of any kind. The Creative Commons licence's disclaimer of warranties and limitation of liability apply.


\needspace{5\baselineskip}
\section[Lucide icon assets]{Lucide icon assets}
The icons in \texttt{assets/icons/} and their rendered counterparts in \texttt{print-edition/icons/} are from \href{https://lucide.dev/}{Lucide}. They are not covered by this repository's CC BY 4.0 licence. Lucide is released under the ISC License, and some Lucide icons are derived from Feather Icons under the MIT License.


\needspace{5\baselineskip}
\subsection[ISC License]{ISC License}
Copyright \textcopyright{} 2026 Lucide Icons and Contributors


Permission to use, copy, modify, and/or distribute this software for any purpose with or without fee is hereby granted, provided that the above copyright notice and this permission notice appear in all copies.


THE SOFTWARE IS PROVIDED "AS IS" AND THE AUTHOR DISCLAIMS ALL WARRANTIES WITH REGARD TO THIS SOFTWARE INCLUDING ALL IMPLIED WARRANTIES OF MERCHANTABILITY AND FITNESS. IN NO EVENT SHALL THE AUTHOR BE LIABLE FOR ANY SPECIAL, DIRECT, INDIRECT, OR CONSEQUENTIAL DAMAGES OR ANY DAMAGES WHATSOEVER RESULTING FROM LOSS OF USE, DATA OR PROFITS, WHETHER IN AN ACTION OF CONTRACT, NEGLIGENCE OR OTHER TORTIOUS ACTION, ARISING OUT OF OR IN CONNECTION WITH THE USE OR PERFORMANCE OF THIS SOFTWARE.


\needspace{5\baselineskip}
\subsection[MIT License for Feather-derived icons]{MIT License for Feather-derived icons}
Copyright \textcopyright{} 2013-present Cole Bemis


Permission is hereby granted, free of charge, to any person obtaining a copy of this software and associated documentation files (the "Software"), to deal in the Software without restriction, including without limitation the rights to use, copy, modify, merge, publish, distribute, sublicense, and/or sell copies of the Software, and to permit persons to whom the Software is furnished to do so, subject to the following conditions:


The above copyright notice and this permission notice shall be included in all copies or substantial portions of the Software.


THE SOFTWARE IS PROVIDED "AS IS", WITHOUT WARRANTY OF ANY KIND, EXPRESS OR IMPLIED, INCLUDING BUT NOT LIMITED TO THE WARRANTIES OF MERCHANTABILITY, FITNESS FOR A PARTICULAR PURPOSE AND NONINFRINGEMENT. IN NO EVENT SHALL THE AUTHORS OR COPYRIGHT HOLDERS BE LIABLE FOR ANY CLAIM, DAMAGES OR OTHER LIABILITY, WHETHER IN AN ACTION OF CONTRACT, TORT OR OTHERWISE, ARISING FROM, OUT OF OR IN CONNECTION WITH THE SOFTWARE OR THE USE OR OTHER DEALINGS IN THE SOFTWARE.



\backmatter
\chapter*{Closing Note}
\addcontentsline{toc}{chapter}{Closing Note}
This guide is designed to support structured review, applied thinking, and professional discussion. Use it as a starting point, verify current requirements with EDAC, and continue learning from practitioners, communities, and public sources.

\vspace{1em}
Constructive feedback, comments, and professional connections are welcome through \href{https://www.linkedin.com/in/xxiao27/}{Xiao's LinkedIn profile}.

\vfill
\begin{center}
\sffamily\color{Moss}
\href{https://www.linkedin.com/in/xxiao27/}{linkedin.com/in/xxiao27}\\[0.5em]
\href{https://github.com/B1ackCat7/ecd-exam-study-guide}{github.com/B1ackCat7/ecd-exam-study-guide}
\end{center}
\end{document}
